\documentclass[a4paper,12pt]{report}
\usepackage[utf8]{inputenc}
\usepackage[T1]{fontenc}
\usepackage[margin=2.5cm]{geometry}
\usepackage{comment}
\usepackage{graphicx}
\usepackage{setspace}
\usepackage{array}
\usepackage{tabularx} 
\usepackage{longtable}
\usepackage{booktabs}
\usepackage{array}
\usepackage{ragged2e}
\usepackage{xcolor}
\usepackage{colortbl}
\usepackage{float}
\usepackage{mdframed} 
%
\definecolor{primaryblue}{RGB}{30,80,140}
\definecolor{accentblue}{RGB}{70,130,180}
\definecolor{lightgray}{RGB}{245,247,250}
%
\newcolumntype{C}[1]{>{\Centering\arraybackslash}p{#1}}
%
\setlength{\LTpre}{0.3cm}
\setlength{\LTpost}{0.3cm}
\renewcommand{\arraystretch}{1.25}
\setlength{\tabcolsep}{2pt}

\usepackage{longtable}
\usepackage{float}
\usepackage{xcolor}
\usepackage[inline]{enumitem}
\usepackage{multirow}
\usepackage{ragged2e}
\usepackage{amssymb}
\usepackage{xparse}
\usepackage[plainpages=false,pdfpagelabels,bookmarksnumbered,hidelinks]{hyperref}
\usepackage{booktabs}          % Professional table rules (\toprule, \midrule, \bottomrule)
\usepackage{colortbl}          % Colored table rows/cells
\usepackage{caption}           % Custom caption styling
\usepackage{subcaption}        % Subfigures
\usepackage{tikz}              % Diagrams and decorations
\usepackage{mdframed}          % Framed/shaded environments
\usepackage{fontawesome5}      % Icons (optional, for visual polish)
\usepackage{tcolorbox}         % Colored boxes for callouts
\usepackage{multirow}          % Merged rows in tables
\usepackage{makecell}          % Line breaks inside cells
\usepackage{listings}          % Code listings
\usepackage{parskip}           % Paragraph spacing
\usepackage{microtype}         % Microtypography improvements
\emergencystretch=2em
\hbadness=10000
\hfuzz=50pt
\vfuzz=50pt
\newcolumntype{L}[1]{>{\RaggedRight\arraybackslash}p{#1}}
\newcolumntype{F}{>{\bfseries\RaggedRight\arraybackslash}p{0.26\linewidth}}
\newcolumntype{V}{>{\RaggedRight\arraybackslash}p{0.68\linewidth}}

\newcommand{\stacklogo}[1]{%
  \raisebox{-.35\height}{\includegraphics[width=2.35cm,height=0.78cm,keepaspectratio]{#1}}%
}
\setlength{\LTpre}{0.45cm}
\setlength{\LTpost}{0.45cm}
\let\PrismIncludeGraphics\includegraphics
\RenewDocumentCommand{\includegraphics}{O{} m}{%
  \IfFileExists{#2}{%
    \PrismIncludeGraphics[#1]{#2}%
  }{%
    \fbox{\parbox[c][0.18\textheight][c]{0.8\linewidth}{\centering
      \textcolor{red}{\textbf{MISSING IMAGE --- ADD BEFORE DEFENSE}}\\[0.4em]
      \footnotesize\texttt{\detokenize{#2}}}}%
  }%
}
\definecolor{ucheader}{RGB}{30,80,140}
\definecolor{uctitle}{RGB}{220,235,255}
\definecolor{ucalt}{RGB}{245,248,255}
\definecolor{ucwarn}{RGB}{255,245,220}
\definecolor{ucgreen}{RGB}{235,248,240}

\definecolor{primaryblue}{HTML}{1A5276}
\definecolor{accentblue}{HTML}{2E86C1}
\definecolor{lightgray}{HTML}{F2F3F4}
\definecolor{medgray}{HTML}{D5D8DC}
\definecolor{successgreen}{HTML}{1E8449}
\definecolor{warnorange}{HTML}{D35400}
\definecolor{sprintBlue}{HTML}{1A5276}

\captionsetup{font=small, labelfont={bf,color=primaryblue}, margin=0.5cm}
\captionsetup[table]{position=bottom,skip=6pt}
\setlength{\abovecaptionskip}{6pt}
\setlength{\belowcaptionskip}{4pt}

\usepackage{titlesec}
\titleformat{\section}{\Large\bfseries\color{primaryblue}}{{\thesection}}{1em}{}[\color{accentblue}\titlerule]
\titleformat{\subsection}{\large\bfseries\color{primaryblue}}{\thesubsection}{1em}{}
\titlespacing*{\section}{0pt}{18pt}{8pt}
\titlespacing*{\subsection}{0pt}{12pt}{6pt}
\newcommand{\comparisonlist}[1]{%
  \begin{minipage}[t]{\linewidth}%
  \begin{itemize}[leftmargin=*, nosep, topsep=0pt]
  #1
  \end{itemize}%
  \end{minipage}%
}
\begin{document}
\hypersetup{pageanchor=false}
\raggedbottom
\sloppy

% ================= COVER PAGE =================
\thispagestyle{empty}

\begin{tabular}{ >{\centering\arraybackslash}m{0.3\textwidth} >{\centering\arraybackslash}m{0.3\textwidth} >{\centering\arraybackslash}m{0.3\textwidth} }
\centering
\small
Tunisian Republic\\
Ministry of Higher Education\\
and\\
Scientific Research
&
\centering
\includegraphics[width=3cm]{logoIsima.png}
&
\centering
\small
University of Monastir\\
Higher Institute of Computer\\
Science of Mahdia\\
THESIS\\
Order Number: \rule{2cm}{0.2pt}
\end{tabular}

\vspace{0.5cm}
\hrule
\vspace{1.5cm}

\begin{center}

{\LARGE \textbf{DISSERTATION OF GRADUATION}}

\vspace{1cm}

Presented at\\
{\large \textbf{Higher Institute of Computer Science of Mahdia}}

\vspace{1cm}

In view of obtaining the diploma of\\
Bachelor in Computer Engineering: Embedded Systems and Internet of Things

\vspace{1.5cm}

Prepared by:\\
{\large \textbf{Maiza Yazid, Rais Mohamed Anas}}

\vspace{0.3cm}

\includegraphics[width=6cm]{psychplatfromlogo.jpeg}

\vspace{0.2cm}
\hrule
\vspace{0.5cm}

{\large \textbf{PsychPlatform}}

\vspace{0.5cm}
\hrule

\vspace{0.8cm}

Defended on: 12 June 2026

\vspace{0.3cm}

\begin{tabular}{ll}
Academic Supervisor & : Ben Aichaa Takwa \\
\end{tabular}

\vfill

Academic Year: 2025-2026

\end{center}

\clearpage

% ================= FRONT MATTER =================
\pagenumbering{arabic}

% ================= ACKNOWLEDGEMENTS =================

\chapter*{Acknowledgements}
\phantomsection
\addcontentsline{toc}{chapter}{Acknowledgements}

Before presenting our work, we would like to express our deep gratitude to all those who contributed, directly or indirectly, to the realization of this project. Their support, encouragement, and guidance were essential throughout this journey.

We would like to extend our sincere thanks to \textbf{Dr. Takwa Ben Aïcha Gader}, our academic supervisor. Thanks to her valuable feedback, rigorous guidance, and insightful advice, we were able to progress in a structured and constructive manner. Her availability and dedication were truly invaluable.

We also wish to thank the members of the jury for the time they dedicate to evaluating our work. Their critical perspective and feedback will be of great benefit for the continuation of our journey.

Finally, we would like to thank our families, friends, and classmates who supported us, directly or indirectly, throughout this rewarding academic experience.

% ================= DEDICATION =================

\chapter*{Dedication}
\phantomsection
\addcontentsline{toc}{chapter}{Dedication}

\thispagestyle{empty}

\begin{center}
    {\itshape
        \vspace*{0.5cm}%
        To my parents, the pillars of my life.\\
        No words can truly capture the depth of my gratitude for your \\
        unconditional support and the silent sacrifices you made to pave my way.\\
        Thank you for believing in me even when I doubted myself, \\
        and for teaching me that perseverance is the quiet engine of success. \\[0.8em]
        To my sisters, my first friends and constant companions.\\
        Thank you for your boundless kindness, your laughter that lightened \\
        the heaviest days, and your unwavering encouragement that kept me moving forward. \\[0.8em]
        To my teachers and mentors, who opened the doors of knowledge.\\
        I am deeply grateful for your passion, your dedication to the craft of teaching, \\
        and for showing me that learning is not just a destination, but a lifelong journey. \\[1.5em]

        \textbf{Yazid Maiza}
    }

    \vspace{0.5cm} % Visual separator between the two dedications
    \rule{0.4\textwidth}{0.4pt} % Subtle horizontal line
    \vspace{0.5cm}

    {\itshape
        To my family, the bedrock of my existence.\\
        You have been my pillar of strength through every challenge and \\
        an endless source of inspiration. Thank you for the countless \\
        ways you have supported my journey and for providing the foundation \\
        upon which I build my future. \\[0.8em]
        
        To my loyal friends, the brothers I chose for myself.\\
        I am deeply grateful for your unwavering moral support and your \\
        steadfast presence during the most difficult times. Thank you for the \\
        shared laughter, the long nights of work, and for always \\
        reminding me of what is possible. \\[0.8em]
        
        To the visionaries and advocates of mental well-being.\\
        To all those who strive daily to democratize access to care and \\
        break the barriers of stigma. This work is a tribute to your \\
        commitment to making psychological support a universal right. \\[1.5em]
    }
        \textbf{Mohamed Anas Rais}
    
\end{center}

% ================= ABSTRACT =================

\chapter*{Abstract}
\phantomsection
\addcontentsline{toc}{chapter}{Abstract}

Mental health care faces a structural problem that technology has not yet solved:
patients struggle to articulate their concerns at the start of a session, and
psychologists arrive without the context they need to make that session
immediately productive. This project addresses that problem directly.

\emph{PsychPlatform} is a full-stack AI-assisted psychological consultation
platform developed as a final-year project for the Bachelor's degree in Embedded
Systems and IoT. The platform connects patients with verified psychologists
through a structured digital workflow that spans the entire consultation
lifecycle, from first contact to post-session follow-up.

At its core is a conversational intake assistant powered by two LLaMA~3 models
served through the Groq inference API and enriched by a dual
Retrieval-Augmented Generation pipeline. One retriever draws from a curated
index of Darija psychological expressions, enabling the chatbot to provide initial
support for North African Arabic dialects while acknowledging that this dialectal
coverage is still evolving and not yet exhaustive. The other queries a clinical
knowledge base of twenty condition-specific psychoeducation documents. When the
patient finishes their pre-session conversation, the system automatically
produces a structured clinical summary containing the dominant emotion, urgency
score, sentiment trend, key themes, and tailored follow-up questions --- giving
the psychologist a clinically grounded briefing before the session begins.

Beyond the intake pipeline, the platform implements a complete consultation
infrastructure: AI-assisted psychologist credential verification using OCR and
LLM-based document analysis; geospatial practitioner discovery; an interactive
calendar and booking lifecycle; a one-time session activation code delivered by
email; real-time Socket.IO messaging with voice message support; PDF session
report generation; longitudinal emotional indicator tracking across sessions;
permission-scoped document question-answering via RAG; a ten-question
post-session rating system; and a multilingual interface supporting English,
French, and Arabic with full right-to-left layout rendering.

The platform is built on the MERN stack and developed across five SCRUM sprints,
with role-based access control enforced throughout for patients, psychologists,
and administrators.

\textbf{Keywords:} Artificial Intelligence, Retrieval-Augmented Generation,
Large Language Models, Natural Language Processing, Mental Health, MERN Stack,
Socket.IO, Groq API, Emotional Analysis, Darija, Multilingual NLP,
Psychologist Verification.


% ================= TABLE OF CONTENTS =================
\tableofcontents
\clearpage

% ================= MAIN CONTENT =================
\hypersetup{pageanchor=true}
\pagenumbering{arabic}

% ================= GENERAL INTRODUCTION =================

\chapter*{General Introduction}
\phantomsection
\addcontentsline{toc}{chapter}{General Introduction}

\vspace{0.5cm}

The digital transformation of the healthcare sector has led to the emergence
of numerous \textit{e-health} solutions aimed at improving accessibility,
efficiency, and quality of care. In the field of mental health, these
technologies provide innovative ways to facilitate communication between
patients and psychologists, while also offering intelligent assistance tools
capable of delivering personalized support and guidance.

In this context, our project, \textbf{PsychPlatform}, proposes an
AI-assisted psychological consultation platform designed to enhance the
interaction between patients and mental health professionals. The platform
integrates modern technologies such as intelligent chatbots, semantic
retrieval systems, and personalized recommendation mechanisms in order to
deliver a more adaptive and human-centered experience.

\vspace{0.3cm}

Although this project is submitted in fulfilment of the Bachelor's degree in
Computer Engineering with a specialisation in \emph{Embedded Systems and
Internet of Things}, its scope deliberately extends toward AI-integrated
full-stack web engineering. This extension is motivated by the growing
convergence between IoT architectures and AI-driven software systems in
digital health, where event-driven communication, real-time data streaming,
and resource-aware API design form a shared engineering foundation.
The non-blocking, event-driven runtime of Node.js, the bidirectional
WebSocket messaging layer (Socket.IO), the asynchronous AI inference pipeline,
and the background indexing workers implemented in PsychPlatform reflect
architectural principles that are equally central to IoT data pipelines,
edge computing systems, and connected health device ecosystems.
The project is therefore positioned as an applied engineering contribution
that leverages the core competencies of the degree programme — concurrent
systems design, real-time communication, and resource-efficient software
architecture — within the emerging domain of AI-assisted digital health.

\vspace{0.4cm}

This report presents the different phases of analysis, design, and
implementation carried out throughout the development process. It is organized
as follows:

\vspace{0.3cm}

\begin{center}
\renewcommand{\arraystretch}{1.4}
\begin{tabular}{|p{3cm}|p{10cm}|}
\hline
\textbf{Chapter} & \textbf{Content} \\ \hline
Chapter 1 & General framework of the project: context, objectives, requirements specification, SCRUM methodology, product backlog, and global architecture. \\ \hline
Chapter 2 & State of the art: scientific foundations, key concepts, and comparative analysis of existing digital mental health platforms. \\ \hline
Chapter 3 & Sprint 1 --- Authentication, RBAC, and initial setup: JWT-based authentication, role-based access control, and base MERN infrastructure. \\ \hline
Chapter 4 & Sprint 2 --- Psychologist onboarding system: professional profile creation, credential upload, AI-assisted verification pipeline, and administrative approval workflow. \\ \hline
Chapter 5 & Sprint 3 --- Search, scheduling, and booking lifecycle: psychologist discovery, geospatial search, calendar management, and appointment booking. \\ \hline
Chapter 6 & Sprint 4 --- Session management, messaging, and clinical summarisation: session activation, real-time Socket.IO communication, voice messaging, PDF report generation, AI chatbot integration with dual RAG, and clinical summary generation. \\ \hline
Chapter 7 & Sprint 5 --- Platform enrichment, governance, and continuous engagement: psychologist dashboard, document AI querying, ratings system, extended admin panel, internationalisation with RTL, floating assistant, and notifications subsystem. \\ \hline
\end{tabular}
\end{center}

\vspace{0.5cm}

The objective of this report is to demonstrate the methodological approach,
technical architecture, and innovative solutions implemented within
PsychPlatform, while highlighting the contribution of artificial intelligence
to the modernization of psychological consultation services.


% ================= CHAPTER 1 =================

\chapter{General Framework of the Project}

% ============================================================
\section{Introduction}
% ============================================================

The digital transformation of healthcare services has opened new possibilities for delivering psychological support remotely, yet the gap between patient need and structured professional care remains significant. Existing platforms either offer generic conversational agents disconnected from clinical follow-up, or teleconsultation tools that rely on static intake forms and provide the psychologist with no structured pre-session intelligence. Neither model makes full use of the capabilities that modern artificial intelligence, real-time communication, and multilingual interfaces now make available.

This project addresses that gap by designing and implementing \emph{PsychPlatform}, an intelligent full-stack psychological consultation platform built on the MERN stack. The platform integrates a Retrieval-Augmented Generation chatbot that supports three written languages and includes preliminary Darija handling, an automated credential verification pipeline, a complete appointment and session lifecycle with secure one-time access codes, real-time Socket.IO messaging, LLM-generated clinical summaries with longitudinal emotional tracking, and a permission-scoped document AI querying system. Its purpose is not to replace the psychologist but to make every minute of the consultation more clinically productive by ensuring that both patient and practitioner arrive at the session better prepared.

This chapter introduces the general framework of the project. It presents the context and problem statement, the platform objectives, the target user categories, the functional and non-functional requirements, the SCRUM-based development methodology and its justification, the product backlog, the sprint planning, the general planning artifacts, the technical environment, and the global application architecture.

% ============================================================
\section{Project Presentation}
% ============================================================

\subsection{Context and Problem Statement}

Mental health support has become an increasingly urgent concern in contemporary healthcare systems. Psychological distress affects individuals across age groups and social contexts, yet access to professional support remains constrained by a persistent set of barriers: social stigma, geographical distance, financial cost, long waiting times, and the difficulty of taking the first step toward seeking help. These barriers create a structural gap between the patient's need for timely psychological assistance and the actual availability of organised professional support.

\begin{figure}[H]
  \centering
  \includegraphics[width=0.85\textwidth]{assets/s1/1.jpeg}
  \caption{Transition from fragmented psychological care access to the integrated PsychPlatform model. The left side illustrates the barriers patients face --- stigma, distance, delays, and lack of structured preparation --- while the right side shows how the platform addresses each barrier through AI-assisted intake, verified practitioner discovery, multilingual support, and a structured consultation workflow.}
  \label{fig:mental-health-context}
\end{figure}

By providing online communication, appointment scheduling, and distant information access, digital health platforms help to mitigate these challenges. However, as the comparative analysis in Chapter 2 shows, the majority of current solutions are still restricted to one of two insufficient models: teleconsultation platforms that rely on static intake forms and deprive the psychologist of organized pre-session context, or generic conversational bots that are not connected to professional follow-up. In both cases, the patient experience remains fragmented, and the clinical value of the session depends entirely on what the patient can articulate spontaneously in the first few minutes.

PsychPlatform responds to this gap by integrating the complete consultation workflow within a single coherent platform. A conversational AI chatbot guided by retrieval-augmented generation conducts a structured pre-session intake in the patient's preferred language, with initial support for Tunisian Darija that remains improvable and not yet fully mature. The resulting conversation is automatically transformed into a structured clinical summary --- including dominant emotion, urgency score, sentiment trend, and key themes --- that the psychologist receives before the session begins. Credential verification, appointment management, real-time messaging, and post-session reporting are all handled within the same platform, eliminating the fragmentation and data isolation that characterise the current generation of digital mental health tools.
\clearpage
\subsection{Main Objectives}

The main objective of the project is to develop an intelligent, secure, and user-centred psychological consultation platform that connects patients, psychologists, and administrators through a coherent digital workflow. Beyond simple online interaction, the platform aims to structure the entire consultation process, from the initial expression of psychological concerns to the delivery of a secure and well-organised professional follow-up.

\begin{figure}[H]
  \centering
  \includegraphics[width=0.85\textwidth]{assets/s1/2.jpeg}
  \caption{The six strategic objectives of PsychPlatform, each addressing a specific limitation of existing solutions: AI-powered conversational intake, automated clinical summarisation, credential verification, full session lifecycle management, multilingual and RTL accessibility, and role-based access control.}
  \label{fig:main-objectives-platform}
\end{figure}

The specific objectives of the project are summarised as follows:

\begin{itemize}[leftmargin=1.5cm, itemsep=6pt]
    \item[--] To design and implement an AI-powered conversational intake module that enables patients to express their psychological concerns through a natural and guided dialogue, enriched by Retrieval-Augmented Generation over a preliminary Darija psychological context index and a clinical PDF knowledge base, rather than through rigid static questionnaires.

    \item[--] To automatically generate structured clinical summaries from chatbot conversations, including emotional indicators, urgency score, sentiment trend, key themes, and actionable recommendations, providing the psychologist with AI-generated pre-session intelligence.

    \item[--] To implement a psychologist verification workflow in which uploaded CV and diploma documents are processed through OCR and LLM-based analysis before being reviewed and validated by the administrator, with the AI serving as a structured briefing tool rather than an autonomous decision-maker.

    \item[--] To support the complete consultation lifecycle, including appointment request, availability management, payment confirmation, one-time session code verification, real-time Socket.IO messaging, voice messaging, session completion, PDF report generation, and post-session rating.

    \item[--] To ensure multilingual accessibility through English, French, and Arabic interfaces with dynamic right-to-left layout adaptation for Arabic, and to include partial Tunisian Darija support in the conversational AI pipeline through a dedicated semantic retrieval index that can be expanded over time.

    \item[--] To enforce role-based access control across patients, psychologists, and administrators, protecting sensitive clinical data and guaranteeing a coherent separation of privileges throughout the platform.
\end{itemize}

\subsection{Target Audience}

The proposed platform is designed to serve three principal categories of users: patients, psychologists, and administrators. Each category interacts with the system through a dedicated set of functionalities while remaining connected to the same secure digital workflow. This organisation ensures that every actor has access only to the services and information that correspond to their responsibilities within the consultation process.

\begin{figure}[H]
  \centering
  \includegraphics[width=0.85\textwidth]{assets/s1/3.jpeg}
  \caption{The three platform actors and their interaction flows. The patient initiates the process through psychologist discovery, AI chatbot preparation, booking, payment, and online consultation. The psychologist manages availability, conducts sessions, and consults AI-generated summaries. The administrator supervises credential verification, user management, and platform governance.}
  \label{fig:target-audience-workflow}
\end{figure}

The platform addresses three principal user categories. Table~\ref{tab:personas} summarises their responsibilities and their expected interactions with the system.


\begin{longtable}{@{} L{3.2cm} L{11.2cm} @{}}
\toprule
\bfseries Persona &
\bfseries Description \\
\midrule
\endfirsthead
\toprule
\bfseries Persona &
\bfseries Description \\
\midrule
\endhead

Patient &
Individuals seeking psychological support through a secure and accessible digital environment. Patients can browse verified psychologist profiles, interact with the AI chatbot before a consultation, submit booking requests, confirm payments through a secure session code, participate in real-time consultations, access their session history, and submit post-session ratings. \\[4pt]
\midrule

Psychologist &
Licensed mental health professionals who register on the platform and submit professional credentials for administrator verification. Once approved, psychologists can manage their availability, accept or reject booking requests, conduct online sessions, consult AI-generated patient summaries, maintain private clinical notes, exchange real-time messages with patients, and generate structured session reports in PDF format. \\[4pt]
\midrule

Administrator &
Governance and supervision actors responsible for maintaining platform trust, security, and operational consistency. Administrators review psychologist verification requests, examine AI-assisted document summaries, approve or reject professional accounts, manage users and roles, monitor aggregated platform statistics, and ensure that each actor operates within the appropriate access privileges. \\

\bottomrule
\caption{Platform user personas and responsibilities}
\label{tab:personas}\\
\end{longtable}

% ============================================================
\section{Requirements Specification}
% ============================================================

The requirements specification translates the project objectives into verifiable system obligations. It distinguishes between functional requirements, which define what the platform must provide, and non-functional requirements, which define the quality attributes expected from the implemented solution.

\subsection{Functional Requirements}

The functional requirements are organised by module in order to maintain traceability between user needs, product backlog items, and implementation tasks.

\vspace{1cm}

\begin{longtable}{@{} L{0.24\textwidth} L{0.14\textwidth} L{0.54\textwidth} @{}}
\toprule
\textbf{Module} &
\textbf{ID} &
\textbf{Requirement} \\
\midrule
\endfirsthead

\toprule
\textbf{Module} &
\textbf{ID} &
\textbf{Requirement} \\
\midrule
\endhead

\emph{Authentication and Account Management} & FR-01 & Users shall be able to register an account with a specified role, either \textit{Patient} or \textit{Psychologist}. \\
\midrule
\emph{Authentication and Account Management} & FR-02 & Users shall be able to authenticate securely using email and password, with JWT-based session management. \\
\midrule
\emph{Authentication and Account Management} & FR-03 & Users shall be able to view and update their personal profile information. \\
\midrule
\emph{Psychologist Profile and Credential Verification} & FR-04 & Psychologists shall be able to create and update their professional profile, including specialisations, languages, city, biography, and session pricing. \\
\midrule
\emph{Psychologist Profile and Credential Verification} & FR-05 & Psychologists shall be able to upload their CV and diploma documents through a guided multi-step setup wizard. \\
\midrule
\emph{Psychologist Profile and Credential Verification} & FR-06 & The system shall extract text from uploaded documents using OCR with Tesseract.js or PDF parsing when applicable, then generate a structured AI summary for administrator review. \\
\midrule
\emph{Psychologist Profile and Credential Verification} & FR-07 & The administrator shall be able to review AI-generated summaries and approve or reject psychologist verification requests. \\
\midrule
\emph{Search and Discovery} & FR-08 & Patients shall be able to browse the list of approved psychologists using multi-criteria filtering. \\
\midrule
\emph{Search and Discovery} & FR-09 & Patients shall be able to search for nearby psychologists using geospatial queries based on their current location. \\
\midrule
\emph{Search and Discovery} & FR-10 & Patients shall be able to consult a psychologist's public profile, including specialisations, ratings, and reviews. \\
\midrule
\emph{Calendar and Booking} & FR-11 & Psychologists shall be able to add, view, and delete availability time slots through an interactive calendar interface. \\
\midrule
\emph{Calendar and Booking} & FR-12 & Patients shall be able to view available slots for a psychologist and submit a booking request. \\
\midrule
\emph{Calendar and Booking} & FR-13 & Psychologists shall be able to confirm or reject incoming booking requests. \\
\midrule
\emph{Session Management} & FR-14 & Upon booking confirmation, patients shall select a session type among \textit{Preparation}, \textit{Follow-up}, or \textit{Free Expression}. \\
\midrule
\emph{Session Management} & FR-15 & The system shall generate a unique six-digit access code after payment confirmation and deliver it to the patient by email. \\
\midrule
\emph{Session Management} & FR-16 & The patient shall enter the session code to activate the shared messaging interface and advance the session to \texttt{active} status. \\
\midrule
\emph{Session Management} & FR-17 & The psychologist shall be able to terminate an active session and generate a downloadable PDF session report. \\
\midrule
\emph{Session Management} & FR-18 & Both parties shall be able to access the history of their past sessions. \\
\midrule
\emph{Session Management} & FR-19 & Patients shall be able to submit a ten-question post-session satisfaction rating. \\
\midrule
\emph{AI Chatbot} & FR-20 & Patients shall have access to an AI-powered chatbot using LLaMA~3.3 through the Groq Cloud API, enriched by RAG over a preliminary Darija psychological context index and a clinical PDF knowledge base. \\
\midrule
\emph{AI Chatbot} & FR-21 & The chatbot shall maintain a persistent, user-scoped conversation history stored in the database. \\
\midrule
\emph{AI Chatbot} & FR-22 & Upon session end, the system shall generate and persist a \texttt{ChatbotSummary} document containing emotional indicators, urgency score, sentiment trend, key themes, and clinical recommendations. \\
\midrule
\emph{Messaging} & FR-23 & Session participants shall be able to exchange real-time text messages through a WebSocket-based interface using Socket.IO. \\
\midrule
\emph{Messaging} & FR-24 & Users shall be able to record and send voice messages through the integrated microphone interface. \\
\midrule
\emph{Messaging} & FR-25 & The system shall track unread message counts and allow messages to be marked as read. \\
\midrule
\emph{Psychologist Dashboard, Admin Panel, and Platform Utilities} & FR-26 to FR-29 & Psychologists shall access a patient dashboard showing session histories, AI summaries, private clinical notes, and uploadable patient documents queryable via permission-scoped RAG. \\
\midrule
\emph{Psychologist Dashboard, Admin Panel, and Platform Utilities} & FR-30 to FR-31 & Administrators shall manage all user accounts and access aggregated platform usage statistics. \\
\midrule
\emph{Psychologist Dashboard, Admin Panel, and Platform Utilities} & FR-32 & All users shall have access to a floating navigation assistant powered by an LLM. \\
\midrule
\emph{Psychologist Dashboard, Admin Panel, and Platform Utilities} & FR-33 to FR-34 & Users shall receive and manage in-app notifications for relevant platform events. \\
\bottomrule
\caption{Functional requirements}
\label{tab:functional-requirements}\\
\end{longtable}

\subsection{Non-Functional Requirements}

Non-functional requirements define the expected quality level of the platform. They are particularly important in this project because psychological data, user authentication, and AI-generated summaries must be handled with caution.

\vspace{1cm}

\begin{longtable}{@{} L{1.6cm} L{2.8cm} L{10.2cm} @{}}
\toprule
\bfseries ID &
\bfseries Category &
\bfseries Requirement \\
\midrule
\endfirsthead
\toprule
\bfseries ID &
\bfseries Category &
\bfseries Requirement \\
\midrule
\endhead

NFR-01 & Security & Passwords shall be hashed using \texttt{bcryptjs} with a work factor of 10 before storage. \\
\midrule
NFR-02 & Security & All protected API routes shall use JWT middleware, and role-based access control shall be enforced at route level. \\
\midrule
NFR-03 & Security & Rate limiting shall be applied globally, including 10 requests per 15 minutes on authentication endpoints, 100 requests per hour on API routes, and 50 messages per hour on the chatbot. \\
\midrule
NFR-04 & Performance & The backend shall rely on Node.js non-blocking I/O to support concurrent session and messaging requests. \\
\midrule
NFR-05 & Scalability & MongoDB Atlas shall be used as the database engine to support cloud-hosted storage and future scaling needs. \\
\midrule
NFR-06 & Availability & The platform shall expose a health-check endpoint and support graceful shutdown in case of database connection failure. \\
\midrule
NFR-07 & Usability & The user interface shall be responsive, built with React~19 and TailwindCSS, and compatible with modern browsers. \\
\midrule
NFR-08 & Localisation & The platform shall support dynamic language switching between English, French, and Arabic without page reload, including RTL rendering for Arabic. \\
\midrule
NFR-09 & Data Integrity & Mongoose schemas shall enforce type validation, uniqueness constraints, and ObjectId references where required. \\
\midrule
NFR-10 & Compliance & The platform shall avoid storing raw credentials or sensitive identifiers in plaintext, and uploaded files shall be stored with unique server-side identifiers. \\

\bottomrule
\caption{Non-functional requirements}
\label{tab:nfr}\\
\end{longtable}

% ============================================================
\section{SCRUM Development Methodology}
% ============================================================

\subsection{Methodological Choice and Justification}

The selection of a development methodology is not a formality. It is an architectural decision that shapes how complexity is managed, how risk is surfaced, and how the delivered system relates to the original requirements. For a project of the nature and scope of PsychPlatform, this choice deserves explicit justification --- including a consideration of the alternative that was not selected.

The principal alternative to SCRUM for an AI-intensive project is CRISP-ML (Cross-Industry Standard Process for Machine Learning), a methodology derived from the CRISP-DM process and specifically adapted for machine learning system development. CRISP-ML structures a project around six sequential phases: business understanding, data understanding, data preparation, model engineering, model evaluation, and model deployment and monitoring. It is the appropriate choice when the primary deliverable is a trained machine learning model, when the project's central uncertainty concerns data quality and model performance, and when the development cycle is dominated by iterative experimentation on datasets and evaluation metrics.

PsychPlatform is not primarily a machine learning project. It does not train models from data. Its AI components --- the LLaMA~3 chatbot, the clinical summarisation engine, the credential verification pipeline, and the RAG retrieval system --- use pre-trained foundation models through inference APIs. The platform's primary complexity is not model engineering but \emph{system integration}: connecting authentication, onboarding, scheduling, real-time communication, AI inference, document processing, multilingual rendering, and administrative governance into a coherent, secure, and user-facing product. CRISP-ML provides no guidance for this class of problem. Its sequential phases would impose a linear structure on a project whose components are interdependent and whose requirements --- particularly around AI-assisted workflows and clinical UX --- can only be validated through working increments.

SCRUM was therefore selected as the methodologically appropriate choice for the following concrete reasons, each tied to a specific challenge of this project:

\begin{itemize}[leftmargin=1.5cm, itemsep=6pt]
  \item[--] \textbf{Interdependent components with evolving integration requirements.} PsychPlatform integrates eight distinct technical subsystems: authentication and RBAC, psychologist onboarding and AI verification, geospatial search, calendar and booking, session management with one-time codes, real-time messaging, the RAG chatbot pipeline, and platform governance features. The interfaces between these subsystems cannot be fully specified before implementation begins, because each sprint reveals integration constraints that are not apparent from requirements documents alone. SCRUM's sprint-by-sprint delivery allows these interfaces to be negotiated and tested incrementally rather than committed to in a single upfront design phase.

  \item[--] \textbf{High-risk AI components requiring isolated validation.} Three platform components involve AI inference under uncertain production conditions: the LLM chatbot, the clinical summarisation engine, and the credential verification pipeline. Each of these components can fail in ways that are difficult to anticipate --- malformed JSON output, provider timeouts, hallucinated content, or prompt injection. Addressing these risks in dedicated sprints (Sprint~2 for verification, Sprint~4 for the chatbot and summarisation) allows the team to design fallback strategies, implement the \texttt{parseJsonFromModel} defensive parser, and validate acceptance criteria in isolation before these components are integrated with patient-facing interfaces.

  \item[--] \textbf{Sensitivity of clinical workflows to usability feedback.} The acceptability of a mental health platform cannot be determined from a requirements document. Features such as the chatbot tone, the clinical summary layout, the booking confirmation flow, and the session code entry interface must be validated by testing working implementations. SCRUM's definition of done requires that every sprint deliverable renders correctly and satisfies its acceptance criteria, ensuring that usability issues are caught at sprint level rather than at final integration.

  \item[--] \textbf{Transparent academic progress tracking.} In the academic context of a supervised final-year project, SCRUM ceremonies provide structured checkpoints for progress reporting, scope management, and retrospective learning. The sprint-by-sprint organisation of this report directly reflects the sprint delivery structure, allowing the academic supervisor and jury to evaluate each functional increment independently.
\end{itemize}

\subsection{SCRUM Roles}

Table~\ref{tab:scrum-roles} presents the SCRUM roles adopted in the project. In this academic context, the responsibilities are adapted to the supervision and development conditions of a final-year engineering project.

\vspace{1cm}

\begin{longtable}{@{} L{3.5cm} L{11cm} @{}}
\toprule
\bfseries Role &
\bfseries Responsibilities \\
\midrule
\endfirsthead
\toprule
\bfseries Role &
\bfseries Responsibilities \\
\midrule
\endhead

Product Owner &
Defines and prioritises the Product Backlog, clarifies acceptance criteria, and validates delivered features against the project vision. In the academic context, this role is supported by the project supervisor and the project initiator. \\[4pt]
\midrule

SCRUM Master &
Facilitates sprint ceremonies, helps remove obstacles, supports process discipline, and contributes to maintaining a realistic development scope during each sprint. \\[4pt]
\midrule

Development Team &
Responsible for design, implementation, AI integration, testing, documentation, and technical validation. The team handles both frontend and backend responsibilities across the different sprints. \\

\bottomrule
\caption{SCRUM role assignments}
\label{tab:scrum-roles}\\
\end{longtable}

% ============================================================
\section{Product Backlog}
% ============================================================

The Product Backlog structures the work to be delivered throughout the project. It is organised by functional epics and ordered according to business value, technical dependency, and implementation complexity. The effort estimation follows a Fibonacci-inspired scale from 1 to 8 story points.

\footnotesize
\begin{longtable}{@{} L{1.1cm} L{2.7cm} L{7.2cm} L{1.8cm} L{1.6cm} @{}}
\toprule
\bfseries ID &
\bfseries Epic &
\bfseries User Story &
\bfseries Priority &
\bfseries Effort \\
\midrule
\endfirsthead
\toprule
\bfseries ID &
\bfseries Epic &
\bfseries User Story &
\bfseries Priority &
\bfseries Effort \\
\midrule
\endhead

PB-01 & Authentication & As a user, I want to register and log in securely so that my account is protected. & High & 3 SP \\
\midrule
PB-02 & Authentication & As a user, I want JWT-based session management so that I remain authenticated across navigations. & High & 2 SP \\
\midrule
PB-03 & Psychologist & As a psychologist, I want to create a detailed professional profile so that patients can evaluate my expertise. & High & 5 SP \\
\midrule
PB-04 & Verification & As a psychologist, I want to upload my CV and diploma for AI-assisted verification and administrator approval. & High & 8 SP \\
\midrule
PB-05 & Verification & As an administrator, I want to review AI-generated summaries of psychologist documents to approve or reject applications. & High & 3 SP \\
\midrule
PB-06 & Search & As a patient, I want to browse and filter approved psychologists to find one matching my needs. & High & 5 SP \\
\midrule
PB-07 & Search & As a patient, I want to search for psychologists near my location to prioritise geographical convenience. & Medium & 5 SP \\
\midrule
PB-08 & Calendar & As a psychologist, I want to manage my availability calendar so that patients can book suitable slots. & High & 8 SP \\
\midrule
PB-09 & Calendar & As a patient, I want to request a booking for an available time slot with a psychologist. & High & 5 SP \\
\midrule
PB-10 & Session & As a patient, I want to confirm payment and receive an access code by email to activate my session. & High & 8 SP \\
\midrule
PB-11 & Session & As a psychologist, I want to end a session and generate a PDF report for consultation records. & Medium & 5 SP \\
\midrule
PB-12 & Messaging & As a session participant, I want to exchange real-time text messages within the shared interface. & High & 8 SP \\
\midrule
PB-13 & Messaging & As a session participant, I want to send voice messages through a microphone interface. & Medium & 5 SP \\
\midrule
PB-14 & AI Chatbot & As a patient, I want to interact with an AI chatbot to express my mental state before a session. & High & 8 SP \\
\midrule
PB-15 & AI Chatbot & As the system, I want to generate a structured clinical summary from the patient's chatbot history. & High & 8 SP \\
\midrule
PB-16 & Dashboard & As a psychologist, I want a patient dashboard showing session histories, AI summaries, and clinical notes. & High & 8 SP \\
\midrule
PB-17 & Documents & As a psychologist, I want to upload and query patient documents using AI to enrich clinical records. & Medium & 5 SP \\
\midrule
PB-18 & Rating & As a patient, I want to submit a ten-question post-session rating to help maintain service quality. & Medium & 3 SP \\
\midrule
PB-19 & Admin Panel & As an administrator, I want to manage user accounts, roles, and platform statistics. & Medium & 5 SP \\
\midrule
PB-20 & i18n & As a user, I want English, French, and Arabic support with RTL layout so that I can use the platform in my language. & Medium & 8 SP \\
\midrule
PB-21 & Assistant & As a user, I want access to a floating AI-powered platform navigation assistant. & Low & 3 SP \\
\midrule
PB-22 & Notifications & As a user, I want to receive and manage in-app notifications for important platform events. & Low & 3 SP \\

\bottomrule
\caption{Product Backlog}
\label{tab:backlog}\\
\end{longtable}
\normalsize

% ============================================================
\section{User Story Acceptance Criteria}
% ============================================================

Acceptance criteria clarify how selected user stories are validated. The following criteria focus on the most technically sensitive features: credential verification, AI chatbot summarisation, and session code verification.

\noindent
\begin{mdframed}[backgroundcolor=lightgray, linecolor=primaryblue, linewidth=1.2pt,
                 innertopmargin=8pt, innerbottommargin=8pt, innerleftmargin=10pt]
PB-04 --- Psychologist Credential Upload and AI Verification
\end{mdframed}

\begin{longtable}{@{} L{2cm} L{12.3cm} @{}}
\toprule
\rowcolor{accentblue!20}
\bfseries Criterion & \bfseries Condition \\
\midrule
AC-1 & The upload form shall accept PDF files for both CV and diploma fields. \\
\midrule
AC-2 & Upon submission, Multer shall save files and Tesseract.js or \texttt{pdf-parse} shall extract text automatically according to file type. \\
\midrule
AC-3 & The Groq LLM shall generate a structured summary of the extracted text within an acceptable response time under normal conditions. \\
\midrule
AC-4 & The administrator panel shall display the AI summary alongside uploaded document links. \\
\midrule
AC-5 & Pending psychologists shall be blocked from patient-facing professional features until administrator approval. \\
\bottomrule
\end{longtable}
\clearpage
\noindent
\begin{mdframed}[backgroundcolor=lightgray, linecolor=primaryblue, linewidth=1.2pt,
                 innertopmargin=8pt, innerbottommargin=8pt, innerleftmargin=10pt]
PB-14 / PB-15 --- AI Chatbot Integration and Clinical Summary Generation
\end{mdframed}

\begin{longtable}{@{} L{2cm} L{12.3cm} @{}}
\toprule
\rowcolor{accentblue!20}
\bfseries Criterion & \bfseries Condition \\
\midrule
AC-1 & The chatbot shall maintain a full user-scoped conversation history persisted in MongoDB. \\
\midrule
AC-2 & System prompts shall enforce an empathetic tone and explicitly prevent the AI from producing clinical diagnoses. \\
\midrule
AC-3 & Upon session-end trigger, the system shall generate a \texttt{ChatbotSummary} containing \texttt{dominantEmotion}, \texttt{urgencyScore}, \texttt{sentimentTrend}, \texttt{keyThemes[]}, \texttt{rawSummary}, and \texttt{recommendations[]}. \\
\midrule
AC-4 & The generated summary shall be accessible to the assigned psychologist through the patient dashboard. \\
\midrule
AC-5 & If the LLM returns malformed JSON, a fallback parsing strategy shall prevent system failure and preserve the valid parts of the response. \\
\bottomrule
\end{longtable}

\noindent
\begin{mdframed}[backgroundcolor=lightgray, linecolor=primaryblue, linewidth=1.2pt,
                 innertopmargin=8pt, innerbottommargin=8pt, innerleftmargin=10pt]
PB-10 --- Session Payment and Code Verification
\end{mdframed}

\begin{longtable}{@{} L{2cm} L{12.3cm} @{}}
\toprule
\rowcolor{accentblue!20}
\bfseries Criterion & \bfseries Condition \\
\midrule
AC-1 & A unique six-digit access code shall be generated and stored after payment confirmation. \\
\midrule
AC-2 & The code shall be delivered to the patient's registered email address using Nodemailer. \\
\midrule
AC-3 & The patient shall enter the code on the \texttt{VerifyCode} page to advance the session status. \\
\midrule
AC-4 & The code shall expire after a configurable duration and become invalid thereafter. \\
\midrule
AC-5 & An incorrect code shall display a clear error message without causing unintended session state modification. \\
\bottomrule
\end{longtable}
\clearpage
% ============================================================
\section{Sprint Planning}
% ============================================================

The implementation was organised into five sprints. Each sprint groups coherent backlog items and aims to reduce technical risk progressively, beginning with core infrastructure and ending with enrichment features. The sprint ordering reflects deliberate risk management: high-risk AI components are introduced in Sprint~2 and Sprint~4, after the authentication boundary and data model are stable, so that fallback strategies can be designed with full knowledge of the surrounding system.

\begin{longtable}{@{} L{1.7cm} L{2.0cm} L{7.0cm} L{3.2cm} @{}}
\toprule
\bfseries Sprint &
\bfseries Weeks &
\bfseries Goal &
\bfseries Backlog Items \\
\midrule
\endfirsthead
\toprule
\bfseries Sprint &
\bfseries Weeks &
\bfseries Goal &
\bfseries Backlog Items \\
\midrule
\endhead

Sprint 1 & 3 weeks & Core infrastructure, authentication, role system, and base MERN architecture. & PB-01, PB-02 \\
\midrule
Sprint 2 & 3 weeks & Psychologist onboarding, professional profile creation, credential upload, AI-assisted verification, and administrator approval flow. & PB-03, PB-04, PB-05 \\
\midrule
Sprint 3 & 3 weeks & Patient engagement, search and discovery, geospatial search, calendar management, and booking lifecycle. & PB-06 to PB-09 \\
\midrule
Sprint 4 & 3 weeks & Core consultation engine, payment confirmation, session code system, Socket.IO messaging, voice messages, session lifecycle, PDF reports, and AI chatbot integration with RAG. & PB-10 to PB-15 \\
\midrule
Sprint 5 & 3 weeks & Platform enrichment, psychologist dashboard, permission-scoped document AI querying, ratings, admin panel, i18n with RTL, floating assistant, and notifications. & PB-16 to PB-22 \\

\bottomrule
\caption{Sprint planning overview}
\label{tab:sprints}\\
\end{longtable}

\noindent
\begin{mdframed}[backgroundcolor=accentblue!10, linecolor=accentblue, linewidth=1pt,
                 innertopmargin=8pt, innerbottommargin=8pt, innerleftmargin=12pt]
A backlog item is considered \textit{Done} when all acceptance criteria are verified, the feature is integrated without regression, relevant API endpoints return expected HTTP status codes under test, and the user interface renders correctly on desktop and mobile viewports.
\end{mdframed}
\clearpage
% ============================================================
\section{General Planning}
% ============================================================

\subsection{Class Diagram}

The class diagram in Figure~\ref{fig:class-diagram} represents the core data model entities and their associations within the MongoDB database layer. The \texttt{User} entity acts as the central identity record, extended by role-specific profile documents for patients and psychologists. Sessions, messages, chatbot histories, clinical summaries, emotional indicators, verification records, notifications, and ratings are linked through ObjectId references, allowing the platform to accumulate a rich longitudinal dataset per patient while preserving clear ownership boundaries between roles.

\begin{figure}[H]
  \centering
  \includegraphics[width=0.98\textwidth]{assets/chapter1/global class diagram.png}
  \caption{Platform class diagram showing the core data model entities and their ObjectId-based associations. The \texttt{User} entity is the central identity record, with role-specific extensions for patients and psychologists and domain entities for sessions, messaging, AI summaries, emotional indicators, and verification data.}
  \label{fig:class-diagram}
\end{figure}

\subsection{Use Case Diagrams by Actor}
The global use case diagram provides a synthetic view of the interactions between the system actors and the main functional domains. The primary actors are the Patient, the Psychologist, and the Administrator, while the System/AI actor represents automated processing: chatbot responses, RAG retrieval, clinical summary generation, document extraction, credential analysis, and notification delivery.

The actor-to-domain participation matrix in Table~\ref{tab:usecase-matrix} complements the diagram by clarifying the responsibility of each actor across the platform's functional domains.



\begin{longtable}{@{} L{5.0cm} C{2.0cm} C{2.8cm} C{1.9cm} C{2.5cm} @{}}
\toprule
\bfseries Functional Domain &
\bfseries Patient &
\bfseries Psych. &
\bfseries Admin &
\bfseries System \\
\midrule
\endfirsthead
\toprule
\bfseries Functional Domain &
\bfseries Patient &
\bfseries Psych. &
\bfseries Admin &
\bfseries System \\
\midrule
\endhead

Authentication and Account & \checkmark & \checkmark & \checkmark & \\
\midrule
Psychologist Profile and Verification & & \checkmark & \checkmark & \checkmark \\
\midrule
Browse and Search & \checkmark & & & \\
\midrule
Calendar and Booking & \checkmark & \checkmark & & \\
\midrule
Session Management & \checkmark & \checkmark & \checkmark & \\
\midrule
Messaging & \checkmark & \checkmark & & \\
\midrule
AI Chatbot and RAG & \checkmark & & & \checkmark \\
\midrule
Ratings and Reviews & \checkmark & & & \\
\midrule
Psychologist Dashboard & & \checkmark & & \\
\midrule
Document AI Querying & & \checkmark & & \checkmark \\
\midrule
Admin Panel & & & \checkmark & \\
\midrule
Platform Assistant & \checkmark & \checkmark & & \checkmark \\
\midrule
Notifications & \checkmark & \checkmark & & \\

\bottomrule
\caption{Actor--domain participation matrix}
\label{tab:usecase-matrix}\\
\end{longtable}

\vspace{1em}
% --- Patient ---
\begin{figure}[H]
  \centering
  \includegraphics[width=0.98\textwidth]{assets/chapter1/patient GUC.png}
  \caption{Use case diagram for the \textbf{Patient} actor. This diagram covers the seven functional domains accessible to patients: authentication and account management, browsing and searching for psychologists, calendar and appointment booking, session participation, secure messaging, AI chatbot interaction via RAG, and rating and reviewing psychologists.}
  \label{fig:use-case-patient}
\end{figure}

% --- Psychologist ---
\begin{figure}[H]
  \centering
  \includegraphics[width=0.98\textwidth]{assets/chapter1/psy GUC.png}
  \caption{Use case diagram for the \textbf{Psychologist} actor. This diagram covers the seven functional domains accessible to psychologists: authentication and account management, profile creation and credential submission, calendar and availability management, session oversight, secure messaging with patients, AI-assisted document querying, and the personal analytics dashboard.}
  \label{fig:use-case-psychologist}
\end{figure}

% --- Administrator ---
\begin{figure}[H]
  \centering
  \includegraphics[width=0.88\textwidth]{assets/chapter1/admin GUC.png}
  \caption{Use case diagram for the \textbf{Administrator} actor. This diagram covers the four functional domains under administrator responsibility: authentication and account management, psychologist credential verification and profile moderation, session oversight and dispute resolution, and platform governance through the admin panel.}
  \label{fig:use-case-admin}
\end{figure}

% --- System / AI ---
\begin{figure}[H]
  \centering
  \includegraphics[width=0.98\textwidth]{assets/chapter1/AI GUC.png}
  \caption{Use case diagram for the \textbf{System\,/\,AI} actor. This diagram covers the five automated functional domains handled without direct user action: credential analysis and verification support, RAG-based chatbot response generation, clinical summary generation, AI-powered document extraction and querying, and automated notification delivery.}
  \label{fig:use-case-system}
\end{figure}
\clearpage

\subsection{Gantt Chart}

The project timeline spans fifteen weeks of active development across five SCRUM sprints of three weeks each, as illustrated in Figure~\ref{fig:gantt}. Sprint~4 carries the highest technical load, integrating the session activation engine, real-time messaging, voice messages, PDF generation, the RAG chatbot pipeline, and the clinical summarisation service within a single sprint cycle.

\begin{figure}[H]
  \centering
    \includegraphics[width=\textwidth]{assets/chapter1/psychplatform_gantt.png}%
  \caption{Project Gantt chart showing the fifteen-week sprint timeline across five SCRUM sprints. Sprint~4 carries the highest integration load, covering the session engine, real-time communication, PDF generation, and the full AI chatbot and clinical summarisation pipeline.}
  \label{fig:gantt}
\end{figure}


% ============================================================
\section{Development Tools and Environment}
% ============================================================

\subsection{Hardware Environment}

The hardware environment illustrated in Figure~\ref{fig:hardware} corresponds to a standard development setup for a MERN-based application integrating OCR, vector embeddings, and external AI inference services.

\begin{figure}[H]
  \centering
  \includegraphics[width=0.9\textwidth]{images/laptop-comparison.png}
  \caption{Hardware environment setup.}
  \label{fig:hardware}
\end{figure}

\subsection{Technical Stack}

The technical stack combines a MERN foundation with additional services for real-time messaging, OCR, vector embeddings, AI inference, and multilingual interfaces. Tables~\ref{tab:backend-stack}, \ref{tab:frontend-stack}, and \ref{tab:devtools} summarise the main tools used.

\vspace{0.3cm}
\noindent\emph{Backend Technologies}

\begin{table}[H]
\centering

\begingroup
\renewcommand{\arraystretch}{1.8}
\setlength{\extrarowheight}{4pt}
\begin{tabularx}{\textwidth}{@{} C{3cm} L{11.5cm} @{}}
\toprule
\bfseries Technology &
\bfseries Role \\
\midrule
\stacklogo{assets/logos/nodejs.png} & Node.js~\cite{NodeOfficial}: JavaScript runtime for the server-side application. \\
\midrule
\stacklogo{assets/logos/express.png} & Express.js~\cite{ExpressOfficial}: Web framework used for RESTful API development and middleware orchestration. \\
\midrule
\stacklogo{assets/logos/mongodb.png} & MongoDB Atlas~\cite{MongoDBAtlasOfficial}: NoSQL document database used for persistent data storage and vector search indexing. \\
\midrule
\stacklogo{assets/logos/mongoose.png} & Mongoose~\cite{MongooseOfficial}: Object Data Modeling layer supporting schema validation and query construction. \\
\midrule
\stacklogo{assets/logos/socketio.png} & Socket.IO~\cite{SocketIOOfficial}: WebSocket library used for bidirectional real-time messaging and live notification delivery. \\
\bottomrule
\end{tabularx}
\caption{Backend technology stack}
\label{tab:backend-stack}
\endgroup
\end{table}

\vspace{0.3cm}
\noindent\emph{Frontend Technologies}

\begin{table}[H]
\centering

\begingroup
\renewcommand{\arraystretch}{1.8}
\setlength{\extrarowheight}{4pt}
\begin{tabularx}{\textwidth}{@{} C{3cm} L{11.5cm} @{}}
\toprule
\bfseries Technology &
\bfseries Role \\
\midrule
\stacklogo{assets/logos/react.png} & React.js~\cite{ReactOfficial}: Component-based framework used for the single-page application. \\
\midrule
\stacklogo{assets/logos/axios.png} & Axios~\cite{AxiosOfficial}: HTTP client used for REST API communication. \\
\midrule
\stacklogo{assets/logos/socketio.png} & Socket.IO Client~\cite{SocketIOOfficial}: WebSocket client used for real-time messaging. \\
\midrule
\stacklogo{assets/logos/tailwind.png} & Tailwind CSS~\cite{TailwindOfficial}: Utility-first CSS framework used for responsive styling and RTL layout adaptation. \\
\midrule
\stacklogo{assets/logos/i18next.png} & i18next~\cite{I18nextOfficial}: Internationalisation framework supporting English, French, and Arabic with RTL rendering. \\
\bottomrule
\end{tabularx}
\caption{Frontend technology stack}
\label{tab:frontend-stack}
\endgroup
\end{table}

\vspace{0.3cm}
\noindent\emph{Development and DevOps Tools}

\begin{table}[H]
\centering

\begingroup
\renewcommand{\arraystretch}{1.8}
\setlength{\extrarowheight}{4pt}
\begin{tabularx}{\textwidth}{@{} C{3cm} L{11.5cm} @{}}
\toprule
\bfseries Tool &
\bfseries Purpose \\
\midrule
\stacklogo{assets/logos/vscode.png} & Visual Studio Code~\cite{VSCodeOfficial}: Primary integrated development environment for full-stack implementation. \\
\midrule
\stacklogo{assets/logos/gitgithub.png} & Git/GitHub~\cite{GitOfficial,GitHubOfficial}: Version control and remote repository management. \\
\midrule
\stacklogo{assets/logos/postman.png} & Postman~\cite{PostmanOfficial}: REST API testing and endpoint validation. \\
\midrule
\stacklogo{assets/logos/compass.png} & MongoDB Compass~\cite{MongoDBCompassOfficial}: Graphical interface for database inspection and query building. \\
\midrule
\stacklogo{assets/logos/plantuml.png} & PlantUML~\cite{PlantUMLOfficial}: UML diagram generation, including use case, class, and sequence diagrams. \\
\bottomrule
\end{tabularx}
\caption{Development and DevOps tooling}
\label{tab:devtools}
\endgroup
\end{table}

% ============================================================
\section{General Application Architecture}
% ============================================================

\subsection{Architectural Choice}

PsychPlatform adopts the MERN stack implemented as a layered RESTful API server complemented by a WebSocket layer for real-time communication. The backend follows a \emph{modular monolith} organisation: each domain --- authentication, chatbot, sessions, verification, messaging, notifications, and dashboard management --- is encapsulated in dedicated route, controller, and model modules, but all modules share the same deployment unit and process space.

This choice was made deliberately in preference to a microservice architecture. A microservice decomposition would have introduced distributed systems complexity --- inter-service networking, independent deployment pipelines, distributed transaction management --- that would have been disproportionate to the team size, the academic timeline, and the integration-heavy nature of the platform. The modular monolith achieves a clear separation of concerns within a single deployable unit, which is more maintainable under the constraints of a two-person team and more straightforward to validate against the acceptance criteria defined in the sprint backlogs.

The platform's AI and external service dependencies --- Groq Cloud API for LLM inference, Google Generative AI for vector embeddings, Tesseract.js for OCR, and Nodemailer for email delivery --- are deliberately kept at the service layer rather than embedded in route controllers. This isolation means that a provider outage or API key rotation affects only the service module, not the surrounding business logic, and fallback strategies can be added without modifying controller code.

\subsection{Architecture Diagram}

The global architecture diagram in Figure~\ref{fig:architecture} illustrates the three-tier organisation of the platform: the React client layer, the Express.js API and Socket.IO server, and the data and external services layer.

\begin{figure}[H]
  \centering
  \fbox{\begin{minipage}{0.92\textwidth}
    \centering
    \includegraphics[width=0.92\textwidth]{images/global archi.png}
  \end{minipage}}
  \caption{Global application architecture of PsychPlatform. The React client communicates with the Express.js API server over HTTPS and connects to Socket.IO for real-time events. The server coordinates business logic, authentication middleware, and calls to external services: Groq Cloud API for LLM inference, Google Generative AI for vector embeddings, MongoDB Atlas for persistent storage and vector search, Tesseract.js for OCR, and Nodemailer for email delivery.}
  \label{fig:architecture}
\end{figure}

\subsection{Key Data Flow Narratives}

The main data flows clarify how the system behaves during the most important user scenarios and verify the consistency between frontend actions, backend services, database persistence, and external AI processing.

\vspace{0.3cm}
\noindent\emph{Patient Onboarding via AI Chatbot}

The patient accesses the chatbot interface and sends a message through the React client. The frontend calls \texttt{POST /api/chatbot/message}. The backend stores the message in the \texttt{ChatbotMessage} collection, then executes two parallel RAG retrieval operations: \texttt{RetrievePsychologicalContext} queries the preliminary Darija psychological context index, and \texttt{RetrieveKnowledgeChunks} queries the clinical PDF knowledge base. Both retrieval results are merged with the conversation history and assembled into a structured prompt before being forwarded to the Groq Cloud API. The generated response is stored and returned to the interface, rendered using \texttt{react-markdown}. At the session-end trigger, the backend submits the full conversation transcript to LLaMA~3.1 8B for clinical summarisation, and stores the resulting \texttt{ChatbotSummary} for psychologist access.

\vspace{0.3cm}
\noindent\emph{Session Lifecycle}

The patient selects an available slot and submits a booking request through \texttt{POST /api/sessions/request}. The psychologist confirms through \texttt{PUT /api/sessions/:id/confirm}. After payment confirmation, the backend generates a cryptographically random six-digit access code, stores it with an expiry timestamp in the \texttt{SessionCode} collection, and delivers it to the patient by email using Nodemailer. The patient enters the code on the \texttt{VerifyCode} page; on successful validation, the session transitions to \texttt{active} status and Socket.IO creates a scoped communication room identified by the session identifier. At session end, the psychologist triggers PDF generation via PDFKit and the session is marked as completed.

\vspace{0.3cm}
\noindent\emph{Psychologist Credential Verification}

The psychologist uploads CV and diploma documents through the guided setup wizard. Multer stores the files on the server. The verification controller extracts document text using \texttt{pdf-parse} for machine-readable PDFs or Tesseract.js for scanned documents. The extracted text is submitted to LLaMA~3.1 8B with a structured verification prompt requesting qualifications assessment, experience estimation, specialisation identification, diploma legitimacy evaluation, and a three-option recommendation: Approve, Review, or Reject. The \texttt{aiVerificationSummary} is stored in the psychologist profile and surfaced in the administrator's review queue alongside the original document files. The administrator makes the final decision, which updates the psychologist's approval status and triggers a notification.

% ============================================================
\section*{Conclusion}
\phantomsection
\addcontentsline{toc}{section}{Conclusion}
% ============================================================

This chapter established the general framework of PsychPlatform by presenting its problem context, platform objectives, user categories, functional and non-functional requirements, development methodology, product backlog, sprint organisation, technical environment, and global architecture. The analysis demonstrates that the platform addresses a genuine and documented gap in the digital mental health landscape: the absence of a unified system that connects AI-assisted patient preparation, verified practitioner discovery, structured clinical intelligence, and secure session management within a single coherent workflow. The SCRUM methodology was selected over process-model alternatives such as CRISP-ML on the basis of a specific architectural argument: PsychPlatform's primary complexity is system integration rather than model training, and SCRUM's incremental delivery model is the appropriate instrument for managing the interdependencies, AI integration risks, and clinical usability constraints that characterise this class of platform. The implementation and validation of each sprint increment are presented in the chapters that follow.


% ================= CHAPTER 2 =================

\chapter{State of the Art}
\label{chap:soa}

\section*{Introduction}
\phantomsection
\addcontentsline{toc}{section}{Introduction}

This chapter presents the scientific and technological foundations of \emph{PsychPlatform}, an AI-assisted psychological consultation platform designed to connect patients with verified psychologists through a secure and structured digital workflow. It reviews the key domains supporting the project, including web engineering, artificial intelligence, natural language processing, real-time communication, geospatial search, retrieval-augmented generation, and credential verification, before analysing existing solutions and identifying their main limitations. This review clarifies the positioning of the proposed platform and highlights the opportunities addressed by its integrated approach.

\section{Scientific Domains and Key Concepts}
\label{sec:domains}

The development of PsychPlatform draws upon a broad and interdisciplinary set of scientific fields. This section introduces each major domain and explains the specific concepts within that domain that are directly relevant to the architecture and functionality of the system.

% -----------------------------------------------------------
\subsection{Web Engineering and Full-Stack Development}
\label{subsec:web-eng}
% -----------------------------------------------------------

Web engineering is the discipline concerned with the systematic application of engineering principles to the design, development, testing, and maintenance of web-based software systems \cite{Murugesan2001}. Modern web engineering has evolved significantly from the era of static HTML documents to the present paradigm of dynamic, data-driven, reactive applications served over distributed architectures.

\subsubsection{The Client-Server Architecture}

The client-server model is the foundational architectural paradigm of modern web applications. Under this model, a \emph{client} (typically a browser running a JavaScript application) communicates with a \emph{server} (a process running on remote hardware) over the HTTP or HTTPS protocol. The server is responsible for processing business logic, authenticating requests, and managing persistent data. The client is responsible for rendering the user interface and managing localised application state.

PsychPlatform follows this paradigm strictly, with the client implemented in React and the server implemented in Node.js with Express. This separation is clinically significant: all sensitive operations --- credential verification decisions, session code generation, clinical summary storage, and role-based access enforcement --- execute exclusively on the server side, where they are protected from client-side manipulation. The React client receives only the data it is authorised to display, and every request carries a JWT token that the server verifies before processing. This architecture ensures that a patient can never access a psychologist's private clinical notes, and that an unapproved psychologist cannot appear in search results, regardless of what a client-side actor might attempt.

\begin{figure}[htbp]
  \centering

  \includegraphics[width=0.8\textwidth]{assets/stateofart/Client-ServerArchitectureofPsychPlatform2.png}
  \caption{High-level client-server communication model in PsychPlatform.
           The React single-page application communicates with the Node.js/Express
           backend exclusively over HTTPS, attaching a JWT token on each protected
           request. MongoDB serves as the persistence layer on the server side.}
  \label{fig:client-server}
\end{figure}

\subsubsection{The MERN Stack}

The MERN stack is a collection of four open-source technologies --- MongoDB, Express.js, React, and Node.js --- that together form a comprehensive JavaScript-based ecosystem for building full-stack web applications. Its principal advantage is the uniformity of the programming language (JavaScript and its superset TypeScript) across both the client and server tiers, which reduces context-switching for developers and facilitates the sharing of validation logic, data models, and utility functions.

\begin{itemize}[leftmargin=1.5em, itemsep=4pt]

  \item[--] \textbf{MongoDB} is a document-oriented NoSQL database that stores data as flexible JSON-like documents (BSON). Its schemaless nature allows for rapid, iterative development of data models, which is particularly valuable in early-stage projects where requirements evolve frequently. PsychPlatform stores all persistent entities in MongoDB: users, psychologist profiles, sessions, messages, chatbot conversation histories, clinical summaries, ratings, calendar slots, notification records, and credential verification documents. The flexible document model is particularly well-suited to this platform because entities such as clinical summaries and AI verification results have heterogeneous, nested structures that would require complex joins in a relational database.

  \item[--] \textbf{Express.js} is a minimal and unopinionated web application framework for Node.js. It provides routing, middleware composition, and HTTP utility methods, enabling the development of structured RESTful APIs. PsychPlatform's backend exposes a comprehensive REST API organised into domain-specific route groups: \texttt{/api/auth}, \texttt{/api/sessions}, \texttt{/api/chatbot}, \texttt{/api/psychologists}, \texttt{/api/admin}, \texttt{/api/notifications}, and \texttt{/api/ratings}. Each route group is protected by authentication middleware and further guarded by role-specific authorisation checks, so that a patient endpoint is structurally unreachable by a psychologist token and vice versa.

  \item[--] \textbf{React} is a declarative, component-based JavaScript library developed by Meta for building user interfaces \cite{Fedosejev2015}. React's virtual DOM diffing algorithm allows for efficient, incremental updates to the user interface without full page reloads. PsychPlatform's frontend is a single-page application (SPA) built with React 19, using React Router for client-side navigation. The role-aware routing layer redirects users to their appropriate dashboard upon login --- patient, psychologist, or administrator --- and prevents access to routes that do not correspond to the authenticated role. The component architecture also enables the chatbot interface, the booking calendar, the clinical summary viewer, and the admin panel to be developed and maintained as independent, reusable modules.

  \item[--] \textbf{Node.js} is a JavaScript runtime built on the V8 engine that allows JavaScript to be executed outside of a browser, enabling server-side development. Its event-driven, non-blocking I/O model makes it particularly well-suited to I/O-intensive applications that handle many concurrent connections. This property is directly relevant to PsychPlatform, which must simultaneously manage active WebSocket sessions between patients and psychologists, outgoing email delivery for session codes, asynchronous AI inference calls to the Groq API, and background document indexing jobs --- all without blocking the main event loop.

\end{itemize}

\begin{figure}[H]
  \centering
  \includegraphics[width=0.85\textwidth]{assets/stateofart/ClientServerInteraction2.png}
  \caption{MERN stack layer diagram showing the interaction between the client-side
           (React), server-side (Node.js and Express.js), and the database (MongoDB).
           JavaScript is used consistently across all layers, enabling seamless
           data flow and development.}
  \label{fig:mern-stack}
\end{figure}

\subsubsection{RESTful API Design}

Representational State Transfer (REST) is an architectural style for distributed hypermedia systems \cite{Fielding2000}. A RESTful API adheres to REST constraints, including stateless communication, a uniform interface (CRUD operations mapped to HTTP verbs: \texttt{GET}, \texttt{POST}, \texttt{PUT}, \texttt{DELETE}), and resource-based URIs.

PsychPlatform exposes a RESTful API with endpoints organised by resource domain. Each endpoint is subject to role-based access control enforced through JWT middleware. For instance, \texttt{PUT /api/sessions/:id/confirm} is restricted to psychologists who own the session, \texttt{GET /api/admin/users} is restricted to administrators, and \texttt{POST /api/chatbot/message} is restricted to authenticated patients. This resource-oriented structure makes the API predictable and reduces the risk of privilege escalation, since access rules are enforced at the route level before any controller logic executes.

\subsubsection{Authentication and Authorisation}

Authentication is the process of verifying the identity of a user or system. Authorisation is the process of determining what an authenticated entity is permitted to do. These two mechanisms are distinct but complementary.

JSON Web Tokens (JWT) are an open standard (RFC~7519) for securely transmitting claims between parties as a digitally signed JSON object. JWTs are stateless: the server does not maintain a session table; instead, the token itself carries the user's identity and role, and its integrity is guaranteed by a cryptographic signature. Upon a valid login, PsychPlatform issues a JWT encoding the user's identifier and role (\texttt{patient}, \texttt{psychologist}, or \texttt{admin}). This token is attached to all subsequent API requests in the \texttt{Authorization} header and verified by the server's authentication middleware before any protected endpoint is accessed. The stateless design is particularly appropriate for a platform that integrates WebSocket connections, REST endpoints, and background AI jobs, since session state does not need to be shared across these execution contexts.

Password security is enforced through \texttt{bcrypt} hashing with a work factor of 10. Bcrypt is a computationally expensive adaptive hashing function designed specifically for password storage, incorporating a work factor to resist brute-force attacks \cite{Provos1999}. Plaintext passwords are never stored or logged at any point in the system.

\subsubsection{Real-Time Communication with WebSockets}

The HTTP request-response protocol is fundamentally unidirectional: the client initiates a request, and the server responds. This model is insufficient for real-time applications where the server must push data to connected clients asynchronously.

WebSocket (RFC~6455) establishes a persistent, full-duplex communication channel over a single TCP connection. Socket.IO abstracts over WebSocket with fallbacks to HTTP long-polling when WebSocket is unavailable, providing room-based broadcasting, event-driven APIs, and automatic reconnection logic. PsychPlatform uses Socket.IO for two distinct real-time features: session messaging between patients and psychologists, scoped to a room identified by the session identifier so that messages are never broadcast to unrelated parties; and live notification delivery, which pushes booking confirmations, session status changes, and report availability events to the appropriate user's private room as soon as the triggering event occurs on the server.

\begin{figure}[H]
  \centering
  \includegraphics[width=0.85\textwidth]{assets/stateofart/WebSocket2.png}
  \caption{Comparison of HTTP request-response polling (left) and the persistent
WebSocket/Socket.IO channel (right) used in PsychPlatform for real-time session messaging.
Both the patient and the psychologist client join a Socket.IO room keyed by the session
identifier.}
  \label{fig:webSocket}
\end{figure}

\subsubsection{Internationalisation (i18n)}

Internationalisation refers to the engineering practice of designing a software system to support multiple languages and locales without requiring changes to the core codebase \cite{W3C2005}. The \texttt{react-i18next} library provides a declarative API for declaring translatable strings using lookup keys and for switching languages at runtime.

PsychPlatform supports English, French, and Arabic, including full right-to-left (RTL) layout rendering for Arabic. This is not a cosmetic addition: research consistently shows that patients disclose emotional content more naturally and accurately in their native language \cite{JurafskyMartin2023}. Implementing RTL support required adopting CSS logical properties throughout the frontend (\texttt{margin-inline-start}, \texttt{padding-inline-end}, etc.) so that layout directions flip automatically on locale switch, and mirroring directional icons such as arrows and progress indicators. The AI chatbot is also designed to detect the patient's input language and respond in kind, ensuring that the most clinically sensitive feature of the platform is equally accessible to Arabic, French, and English speakers. User language preferences are persisted in the \texttt{UserPreference} document and applied on every subsequent login.

% -----------------------------------------------------------
\subsection{Artificial Intelligence and Machine Learning}
\label{subsec:ai-ml}
% -----------------------------------------------------------

Artificial Intelligence (AI) encompasses a broad field of computer science concerned with the development of systems that exhibit behaviours associated with human intelligence, such as reasoning, learning, natural language understanding, and perception \cite{Russell2020}. Machine Learning (ML) is a subfield of AI in which systems learn patterns from data rather than being programmed with explicit rules \cite{Mitchell1997}.

Within PsychPlatform, AI is not a peripheral feature added at the end of development --- it is a core architectural component present across the entire consultation lifecycle. The platform integrates AI at five distinct functional layers: the patient-facing conversational agent; the automated clinical summarisation and emotional analysis engine; the credential verification assistant for psychologist onboarding; the document question-answering pipeline for psychologists; and the platform navigation assistant available to all users. This breadth of AI integration distinguishes PsychPlatform from existing solutions, which typically apply AI to at most one of these layers.

\subsubsection{The Role of AI in Mental Health Applications}

The application of AI to mental health support is a rapidly growing area of research and commercial development. Broadly speaking, AI in this domain serves two non-exclusive roles. The first is \emph{screening and triage}: AI systems analyse patient-provided data to detect indicators of mental health conditions and assess severity. The second is \emph{therapeutic augmentation}: AI systems supplement, extend, or prepare the patient for human-led therapeutic sessions by providing a safe, non-judgmental conversational space.

PsychPlatform belongs firmly to the second category. Its chatbot system prompt explicitly prohibits the model from producing clinical diagnoses or medication suggestions, and the acceptance criteria for the chatbot module (AC-2, PB-14) enforce this constraint as a testable condition rather than a best-effort guideline. The platform's purpose is not to replace the psychologist but to make the time the patient spends with the psychologist more productive by ensuring both parties arrive at the session better prepared.

% -----------------------------------------------------------
\subsection{Natural Language Processing}
\label{subsec:nlp}
% -----------------------------------------------------------

Natural Language Processing (NLP) is the branch of AI concerned with enabling computers to understand, interpret, generate, and respond to human language \cite{JurafskyMartin2023}. NLP spans a wide spectrum of tasks, from low-level morphological analysis to high-level semantic understanding such as question answering, sentiment analysis, and discourse modelling.

PsychPlatform applies NLP across several functional layers, each addressing a distinct clinical or operational need.

\begin{itemize}[leftmargin=1.5em, itemsep=6pt]

  \item[--] \textbf{Sentiment Analysis.} The computational identification of subjective information --- opinions, emotions, and attitudes --- from textual data. In PsychPlatform, sentiment analysis is performed by the LLM at the end of each chatbot session. The full conversation transcript is submitted with a structured prompt requesting an assessment of the patient's emotional tone and a directional trend classification: improving, stable, or declining. This trend is surfaced to the psychologist as the \texttt{sentimentTrend} field of the clinical summary, providing a longitudinal signal that a static mood questionnaire could not capture.

  \item[--] \textbf{Emotion Detection.} A more fine-grained form of sentiment analysis that classifies text according to discrete emotional categories. PsychPlatform computes quantitative scores across four emotional dimensions --- anxiety, sadness, anger, and positivity --- for each patient session, storing them in the \texttt{EmotionalIndicator} model on a normalised scale of 0 to 100. These per-session scores accumulate over time to form a longitudinal emotional trajectory that the psychologist can inspect as a time-series chart in the patient dashboard.

  \item[--] \textbf{Abstractive Text Summarisation.} The automatic production of a condensed narrative from a longer text. Unlike extractive summarisation, which selects existing sentences, abstractive summarisation generates novel sentences that preserve the essential meaning. PsychPlatform employs abstractive summarisation to produce the \texttt{rawSummary} field of the clinical summary, transforming a potentially lengthy chatbot conversation into a concise clinical narrative that the psychologist can read in under a minute before the session.

  \item[--] \textbf{Key Theme Extraction.} The identification of the most salient topics discussed in a body of text. The \texttt{keyThemes} array in the clinical summary is generated by prompting the LLM to identify the principal concerns the patient raised during the chatbot session. This allows the psychologist to immediately see whether the patient's conversation centred on relationship difficulties, work stress, sleep problems, or other themes, without reading the full transcript.

  \item[--] \textbf{Multilingual and Dialectal Processing.} PsychPlatform's chatbot is designed to detect the language of the patient's input automatically and respond in kind, supporting English, French, standard Arabic, and partial Tunisian Darija handling within a single conversational pipeline. Darija support is particularly significant because it is a widely spoken but linguistically under-resourced dialect that most commercial NLP systems do not handle adequately. In this version, PsychPlatform addresses this through a dedicated Darija retrieval index, but its coverage remains preliminary and requires further enrichment to reach more reliable dialectal understanding, as described further in Section~\ref{subsec:RAG}.

\end{itemize}

% -----------------------------------------------------------
\subsection{Large Language Models}
\label{subsec:llm}
% -----------------------------------------------------------

Large Language Models (LLMs) are a class of neural network-based AI systems trained on massive corpora of text with the objective of modelling the probability distribution of token sequences. The transformer architecture \cite{Vaswani2017}, with its self-attention mechanism, is the foundational building block of all modern LLMs. The pretraining phase exposes the model to billions of text tokens, enabling it to acquire broad linguistic and world knowledge. Models such as GPT-4 \cite{OpenAI2023}, LLaMA~3 \cite{MetaAI2024}, Claude \cite{Anthropic2024}, and Gemini \cite{GoogleDeepMind2024} have demonstrated remarkable capabilities in instruction following, multi-turn dialogue, code generation, and complex reasoning.

\begin{figure}[H]
  \centering
  \includegraphics[width=0.85\textwidth]{assets/stateofart/LLMtraining2.png}
  \caption{Simplified overview of the transformer-based LLM architecture \cite{Vaswani2017}.
Input tokens are projected into Query, Key, and Value matrices; self-attention scores
determine contextual representations; and a softmax output head produces the next-token
probability distribution. PsychPlatform employs two LLaMA~3 variants served via the
Groq inference API.}
  \label{fig:LLM_Pre-Training}
\end{figure}

PsychPlatform integrates two LLaMA~3 model variants served through the Groq inference API. The choice of two distinct models rather than a single unified model reflects a deliberate engineering trade-off between response quality and latency, matched to the requirements of each task.

\begin{itemize}[leftmargin=1.5em, itemsep=6pt]

  \item[--] \textbf{LLaMA~3.3 70B Versatile} is used for the patient-facing conversational chatbot and for the floating platform navigation assistant. Its 70 billion parameters enable nuanced, contextually appropriate, and empathetic dialogue generation --- qualities that are essential when a patient is expressing distress or vulnerability. The higher latency of this model is acceptable in a conversational context where users expect a brief pause between message and reply. It is also this model that handles multilingual input and assists with Darija interpretation, although Darija reliability depends on the quality and coverage of the supporting retrieval index.

  \item[--] \textbf{LLaMA~3.1 8B Instant} is used for all structured JSON generation tasks: clinical summary production, credential document analysis, and the generation of tailored follow-up questions for psychologists. For these tasks, output format compliance and low latency are more important than conversational depth. The 8B model produces well-structured JSON reliably and responds within the time constraints of a synchronous API request, making it suitable for operations that block the user interface while they execute.

\end{itemize}

The use of LLMs in PsychPlatform represents a significant departure from earlier chatbot paradigms. Rule-based systems respond according to predefined scripts and decision trees --- they are robust but brittle, failing when user input deviates from anticipated patterns. Retrieval-based models select responses from a fixed candidate pool, offering slightly greater flexibility but still constrained by the quality and scope of the response database. LLMs, by contrast, generate responses dynamically, enabling them to handle the full unpredictability of natural human emotional expression \cite{Brown2020}. This generative capability is what makes it possible for the chatbot to respond appropriately when a patient raises an unexpected topic, switches language mid-conversation, or expresses distress in culturally specific terms.
\clearpage
% -----------------------------------------------------------
\subsection{Prompt Engineering}
\label{subsec:prompt}
% -----------------------------------------------------------

Prompt engineering is the discipline of designing, structuring, and iterating upon the textual instructions provided to an LLM in order to elicit specific, reliable, and appropriately constrained outputs \cite{White2023}. Since LLMs are highly sensitive to the phrasing and structure of their input, prompt engineering is a critical engineering skill in any application that relies on LLM inference.

\begin{figure}[H]
  \centering
  \includegraphics[width=0.85\textwidth]{assets/stateofart/PromptEngPipeline2.png}
  \caption{End-to-end prompt engineering pipeline in PsychPlatform. A structured
           system prompt,comprising role assignment, behavioural guardrails,
           output format constraints, and a calibrated temperature value,wraps
           the user message before submission to the Groq inference API. The raw
           model response is then cleaned and parsed by \texttt{parseJsonFromModel}
           to yield a reliable, structured JSON object.}
  \label{fig:PromptEngPipeline}
\end{figure}

PsychPlatform employs a layered prompt engineering strategy that differs by task type. Each LLM call is prefixed with a system prompt that assigns a specific role to the model and establishes its behavioural boundaries. For conversational tasks, the chatbot system prompt assigns the role of compassionate listening assistant, explicitly prohibits clinical diagnosis and medication suggestions, instructs the model to detect and match the patient's language on every turn, and sets a temperature of \texttt{0.7} to introduce natural variation in phrasing and prevent the conversation from feeling scripted. For structured analytical tasks --- clinical summarisation, credential verification, and follow-up question generation --- the system prompt assigns a clinical or verification analyst role, constrains the output to a precisely specified JSON schema, and sets a temperature of \texttt{0.3} to ensure deterministic, format-compliant responses.

Because LLMs may deviate from the requested output format by wrapping JSON in Markdown code fences or inserting explanatory prose, PsychPlatform implements a defensive parsing function (\texttt{parseJsonFromModel}) that strips known formatting artefacts before applying \texttt{JSON.parse}, and falls back to boundary-scanning the text for the outermost brace pair when standard parsing fails. This fallback is essential in a clinical context: a malformed summary that crashes the application is more harmful than a partial summary that degrades gracefully. The acceptance criterion AC-5 of PB-14 formally specifies this requirement: if the LLM returns malformed JSON, the fallback parsing strategy must prevent system failure and preserve any valid fields in the response.

% -----------------------------------------------------------
\subsection{Sentiment Analysis and Emotional State Modelling}
\label{subsec:sentiment}
% -----------------------------------------------------------

Sentiment analysis occupies a central functional role in PsychPlatform through the generation of clinical summaries. When a patient's chatbot session is concluded, the full conversation transcript is submitted to LLaMA~3.1 8B with a structured prompt requesting a multi-dimensional assessment of the patient's emotional state. The output is a JSON document containing the following fields, each of which serves a distinct clinical purpose:

\begin{itemize}[leftmargin=1.5em, itemsep=4pt]
  \item[--] \textbf{Dominant Emotion} (\texttt{dominantEmotion}): A single-word classification of the most prevalent emotional state expressed in the conversation, giving the psychologist an immediate high-level signal before reading further.
  \item[--] \textbf{Urgency Score} (\texttt{urgencyScore}): A numerical score from 1 to 5 representing the clinical urgency of the patient's situation. A score of 4 or 5 is intended to prompt the psychologist to prioritise the upcoming session or consider an earlier contact.
  \item[--] \textbf{Sentiment Trend} (\texttt{sentimentTrend}): A categorical assessment of whether the patient's emotional tone improved, remained stable, or declined over the course of the chatbot conversation --- a longitudinal signal not available from a static intake form.
  \item[--] \textbf{Key Themes} (\texttt{keyThemes}): An array of the main topics raised by the patient, enabling the psychologist to prepare targeted questions before the session begins.
  \item[--] \textbf{Raw Summary} (\texttt{rawSummary}): A free-text narrative summary of the full conversation, written as a clinical briefing rather than a transcript.
\end{itemize}

In addition to per-session summaries, PsychPlatform maintains an \texttt{EmotionalIndicator} model that stores quantitative scores for anxiety, sadness, anger, and positivity on a normalised 0--100 scale for each completed session. These scores accumulate across sessions to form a longitudinal time series, enabling the psychologist to observe whether therapeutic progress is occurring over multiple consultations --- an outcome-oriented capability absent from all surveyed platforms.

\begin{figure}[H]
  \centering
  \includegraphics[width=0.85\textwidth]{assets/stateofart/ClinicalSummary2.png}
  \caption{Left: example structured clinical summary JSON produced by LLaMA~3.1 8B
after a patient's chatbot session, showing the five output fields stored in the database.
Right: illustrative longitudinal \texttt{EmotionalIndicator} chart for a patient across five
consecutive sessions, showing per-session scores for anxiety, sadness, anger, and positivity
on a normalised 0--100 scale.}
  \label{fig:ClinicalSummary}
\end{figure}

% -----------------------------------------------------------
\subsection{Geospatial Computing and Location-Based Services}
\label{subsec:geo}
% -----------------------------------------------------------

Geospatial computing involves the processing of geographic information to solve location-based problems. In the context of web applications, location-based services allow users to discover entities of interest within a defined geographic proximity.

PsychPlatform integrates geospatial search to enable patients to locate psychologists near their physical location. The \texttt{Psychologist} schema stores a GeoJSON \texttt{Point} object containing latitude and longitude coordinates. A \texttt{2dsphere} index is created on this field in MongoDB, which enables the geospatial query operators \texttt{\$near} and \texttt{\$geoWithin} to compute proximity-based search results using spherical Earth geometry rather than planar approximations. When a patient enables location sharing in the browser, their coordinates are sent to the \texttt{GET /api/psychologists/nearby} endpoint, which executes a \texttt{\$near} query and returns a list of approved psychologists sorted in ascending order of distance. This feature is particularly relevant in the Tunisian and North African context, where patients may need to determine whether a psychologist is close enough to visit in person, or whether an online session is their only practical option.

\begin{figure}[H]
  \centering
  \includegraphics[width=0.85\textwidth]{assets/stateofart/2dsphere.png}
  \caption{Geospatial proximity search in PsychPlatform. A patient's GPS coordinates
  are used as the origin of a \texttt{\$near} query against the \texttt{2dsphere}-indexed
  location field of the Psychologist collection. Results are returned in ascending order
  of distance, enabling the patient to discover the nearest available practitioners.}
  \label{fig:2dsphere}
\end{figure}

% -----------------------------------------------------------
\subsection{Document Parsing and AI-Assisted Credential Verification}
\label{subsec:cred}
% -----------------------------------------------------------

The onboarding of psychologists to the platform requires validation of professional credentials to protect patients from unqualified practitioners. PsychPlatform implements an automated preliminary screening workflow that uses AI to assist the administrator in this process.

When a psychologist submits their CV and diploma as PDF documents, the server extracts the textual content using the \texttt{pdf-parse} library, with Tesseract.js as an OCR fallback for scanned documents. The extracted text is then submitted to LLaMA~3.1 8B with a structured verification prompt requesting an assessment of the candidate's qualifications, estimated years of experience, identified specialisations, the apparent legitimacy of the diploma, and an overall recommendation from three options: Approve, Review, or Reject. The resulting \texttt{aiVerificationSummary} is stored in the psychologist's profile and displayed in the administrator's pending verifications panel alongside the original document files. The administrator then makes the final approval or rejection decision, with the AI summary serving as a structured briefing rather than an autonomous verdict.

This design reflects a deliberate human-in-the-loop architecture. The AI dramatically reduces the cognitive effort required for routine review --- a clearly complete and consistent credential set will typically receive an Approve recommendation, allowing the administrator to confirm it in seconds. Cases flagged as Review or Reject receive focused human attention. This approach reflects the broader principle that AI should augment and accelerate human judgment rather than replace it in safety-critical decisions \cite{Dellermann2019}.

\begin{figure}[H]
  \centering
  \includegraphics[width=0.85\textwidth]{assets/stateofart/CredentialVerificationWorkflow2.png}
  \caption{Step-by-step AI-assisted credential verification workflow in PsychPlatform.
PDF documents are parsed server-side using \texttt{pdf-parse} or Tesseract.js, the extracted
text is submitted to LLaMA~3.1 8B for structured analysis, and the resulting
\texttt{aiVerificationSummary} is surfaced to the administrator. The final approval
decision is made by a human, preserving oversight in this safety-critical process
\cite{Dellermann2019}.}
  \label{fig:CredentialVerificationWorkflow}
\end{figure}

% -----------------------------------------------------------
\subsection{Rate Limiting and API Security}
\label{subsec:rate-limiting}
% -----------------------------------------------------------

In distributed web systems, API endpoints are subject to abuse through denial-of-service attacks, credential stuffing, and excessive automated requests. Rate limiting restricts the number of requests a client can make to an endpoint within a defined time window, protecting both the backend infrastructure and external paid services such as the Groq inference API.

PsychPlatform applies three distinct rate limiters, each calibrated to the specific abuse risk of its endpoint category. Authentication endpoints are limited to 10 requests per 15 minutes per IP address, which blocks automated credential-stuffing attacks while allowing legitimate users to retry a forgotten password. General API endpoints are limited to 100 requests per hour, protecting MongoDB and server resources from scraping or denial-of-service attempts. The chatbot endpoint is limited to 50 messages per hour, which prevents both abuse of the Groq API (which charges per token) and adversarial attempts to exhaust the platform's inference budget through automated message floods. These limits are summarised in Table~\ref{tab:rate-limiters}.

\begin{table}[H]
  \centering
  \caption{Rate limiting configuration per endpoint category in PsychPlatform.}
  \label{tab:rate-limiters}
  \renewcommand{\arraystretch}{1.35}
  \begin{tabularx}{0.85\linewidth}{@{} l X r r @{}}
    \toprule
    \rowcolor{medgray}
    \textbf{Endpoint Category} & \textbf{Purpose} & \textbf{Limit} & \textbf{Window} \\
    \midrule
    Authentication             & Prevent brute-force login attacks & 10 req  & 15 min \\
\midrule
    \rowcolor{lightgray}
    General API                & Protect backend resources          & 100 req & 1 hour \\
\midrule
    Chatbot                    & Prevent LLM inference abuse        & 50 msg  & 1 hour \\
    \bottomrule
  \end{tabularx}
\end{table}

%=====================================================================
\subsection{Retrieval-Augmented Generation}
\label{subsec:RAG}
%=====================================================================

Retrieval-Augmented Generation (RAG) is an architectural pattern introduced by Lewis et al. \cite{Lewis2020} that addresses a fundamental limitation of large language models: their knowledge is static, encoded in weights at training time and inaccessible to update without retraining. This creates two problems for clinical applications. First, a model cannot access information postdating its training cutoff. Second, and more critically, it cannot reason about private or domain-specific documents --- such as a patient's uploaded clinical file or a curated corpus of Darija psychological expressions --- that were never part of its training data.

RAG resolves both problems by decoupling knowledge storage from language generation. Rather than relying on parametric memory alone, a RAG system retrieves relevant passages from an external index at inference time and injects them into the model's prompt context. The model then generates a response grounded in the retrieved evidence, which substantially reduces hallucination and makes every claim traceable to a verifiable source.

\subsubsection{The RAG Pipeline}

A RAG system operates through three sequential stages. During \emph{ingestion}, documents are parsed, split into overlapping text chunks of manageable size, and converted into dense vector embeddings using a pre-trained embedding model. These vectors are stored in a vector database alongside their source metadata. During \emph{retrieval}, a user's query is embedded with the same model and compared against all stored vectors using a similarity metric such as cosine distance; the most semantically relevant chunks are returned. During \emph{generation}, the retrieved chunks are assembled into a structured prompt alongside the original query, and the language model produces a response that draws from the provided evidence rather than from its weights alone.

\subsubsection{Two Distinct RAG Applications in PsychPlatform}

PsychPlatform implements RAG in two architecturally distinct functional contexts, serving different user needs and employing different retrieval strategies.

The \textbf{first application} is proactive context enrichment in the intake chatbot. Every time a patient sends a message, the system automatically executes two parallel retrieval operations before generating a reply. The first retriever (\texttt{RetrievePsychologicalContext.js}) performs a semantic similarity search against a curated index of Darija psychological expressions and culturally situated support content, giving the chatbot a first layer of support for North African Arabic dialects and culturally specific emotional registers that a general-purpose LLM would handle poorly. However, this Darija layer is not considered complete: it improves dialect awareness but still requires broader lexical coverage, more real-world examples, and further evaluation before it can be treated as fully reliable. The second retriever (\texttt{RetrieveKnowledgeChunks.js}) queries a clinical PDF knowledge base and returns relevant therapeutic guidance chunks. Both retrieval results are merged with the conversation history and injected into the system prompt before the response is generated by LLaMA~3.3 70B. This use of RAG is \emph{proactive}: retrieval is triggered automatically on every message turn, and the patient is unaware of the underlying mechanism. A background process in \texttt{chatbotController.js} performs embedding-based gap detection and reseeding to maintain the quality and coverage of the knowledge index over time.

The \textbf{second application} is reactive document question-answering, introduced in Sprint~5. When a psychologist or patient explicitly queries an uploaded document, the system embeds the question using the Google Gemini \texttt{text-embedding-004} model, applies a permission-scoped pre-filter to the vector index to ensure that only chunks within the user's authorised document scope are considered, and retrieves the top-$k$ most semantically relevant chunks. The retrieved chunks are assembled into a grounded prompt and submitted to LLaMA~3.1 8B, which generates a cited answer alongside source excerpt references. Every query is persisted as a \path{DocumentQueryLog} entry for auditability. Unlike the chatbot pipeline, this use of RAG is \emph{reactive}: retrieval is triggered by an explicit user question, and the sourcing of the answer is made visible to the user so that every claim can be verified.

Both pipelines share the same embedding infrastructure (Gemini \texttt{text-embedding-004} and MongoDB Atlas Vector Search) and implement a three-tier degradation fallback --- Atlas Vector Search, followed by MongoDB text-index search, followed by regex scanning --- to ensure continued partial functionality under infrastructure failure.

This dual-RAG architecture is one of the most technically distinctive aspects of PsychPlatform. The two pipelines address the same underlying limitation of LLMs --- lack of access to private, domain-specific knowledge --- but apply the solution in opposite interaction modes: one enriches every conversational turn invisibly, the other responds to explicit queries with transparent, auditable sourcing.

\begin{figure}[H]
  \centering
  \includegraphics[width=0.98\textwidth]{assets/stateofart/dual_rag_psychplatform.png}
  \caption{The dual RAG architecture of PsychPlatform. Left: the intake chatbot pipeline
executes two parallel retrievers (preliminary Darija psychological context and clinical PDF chunks)
automatically on every patient message turn before generating a reply with LLaMA~3.3 70B.
Right: the document Q\&A pipeline responds to explicit user queries with permission-scoped
vector retrieval and cited, auditable answers from LLaMA~3.1 8B. Both pipelines share the
Gemini \texttt{text-embedding-004} model and MongoDB Atlas Vector Search infrastructure.}
  \label{fig:DualRAGPipeline}
\end{figure}

%=====================================================================
\subsection{Clinical knowledge base}
\label{subsec:CKB}
%=====================================================================

The clinical PDF knowledge base indexed by \texttt{RetrieveKnowledgeChunks.js} consists of twenty condition-specific psychoeducation documents covering depression, bipolar disorder, post-traumatic stress disorder, obsessive-compulsive disorder, schizophrenia, eating disorders, perinatal depression, psychosis, stress, anxiety, and adolescent mental health. Several documents originate from evidence-based intervention programmes, including the \emph{Problem Management Plus} (PM+) training manual developed by the World Health Organization \cite{WHO_PMP2016} and the Mental Health Foundation UK anxiety report \cite{MHF2023}. This deliberate curation ensures that the chatbot's retrieved context is clinically grounded rather than generic, and that responses to patients expressing condition-specific distress are informed by recognised therapeutic frameworks rather than by the LLM's parametric knowledge alone.


\subsection*{Introduction to the Comparative Review}
\phantomsection
\addcontentsline{toc}{subsection}{Introduction to the Comparative Review}
 
The preceding sections established the scientific and technological foundations upon
which PsychPlatform is built. Before presenting the platform's design and implementation,
it is necessary to examine what already exists in the digital mental health landscape and
to identify, precisely and critically, the limitations that motivated the present work.
This section analyses five representative platforms drawn from two categories: general
teleconsultation and appointment management tools, and AI-driven conversational mental
health assistants. The analysis is structured around the seven criteria that directly
correspond to PsychPlatform's core design objectives, and concludes with a synthesis
that derives the platform's architectural decisions from the gaps observed.
 
% -----------------------------------------------------------
\subsection{Platforms Under Review}
\label{subsec:platforms}
% -----------------------------------------------------------
 
The five platforms selected for review were chosen to represent the most relevant
categories of prior art. They are not exhaustive but are sufficiently representative
to expose the structural limitations of existing approaches.
 
\subsubsection{BetterHelp}
 
BetterHelp~\cite{BetterHelp2023} is one of the largest English-language online therapy
platforms, connecting users with licensed therapists through text, video, and voice
sessions. Users complete an initial questionnaire upon registration, which the platform
uses to match them with an available therapist. The platform's principal strength is
its scale and accessibility: it removes geographical barriers and offers asynchronous
messaging between sessions.
 
However, BetterHelp presents several significant limitations for the present project's
context. Its intake questionnaire is a static form that produces a matching recommendation
but no structured clinical document readable by the therapist. The therapist has no
AI-generated pre-session briefing and must reconstruct the patient's concerns from
scratch at the start of the first session. The platform offers no psychologist credential
verification pipeline; therapists are vetted through a manual process conducted
internally by BetterHelp staff. There is no multilingual support for North African
Arabic dialects, and the platform operates exclusively in English, making it inaccessible
to Tunisian patients who express themselves in Darija or French. Furthermore, the
platform does not offer geospatial discovery, as it is designed for exclusively remote
interaction.
 
\subsubsection{Doctolib}
 
Doctolib~\cite{Doctolib2023} is a European medical appointment booking platform
operating primarily in France, Germany, and Italy, with a growing presence in the
broader francophone world including parts of North Africa. It allows patients to search
for healthcare practitioners, including psychologists, by location, specialty, and
availability, and to book appointments online. Its geospatial search and availability
calendar are well-engineered and widely adopted.
 
Doctolib, however, is a scheduling platform, not a consultation support system. It
provides no AI-assisted intake, no chatbot interaction, no pre-session clinical
briefing, and no session report generation. The practitioner verification process
relies on the French national health system's identifier (\emph{RPPS number}), which
is not applicable in Tunisia or other North African countries where no equivalent
national registry API exists. Doctolib offers no support for Darija or Arabic RTL
interfaces, and has no post-session emotional tracking capability. Its contribution
to the digital mental health workflow ends at the moment the appointment is confirmed.
 
\subsubsection{Woebot}
 
Woebot~\cite{Woebot2023} is an AI-powered mental health chatbot developed from research
at Stanford University and deployed as a standalone mobile application. It engages
users in structured conversations grounded in cognitive-behavioural therapy (CBT) and
dialectical behaviour therapy (DBT) principles. Woebot represents the most technically
relevant prior art for PsychPlatform's conversational intake module.
 
Its principal limitation is that it is a closed loop: it operates independently of
any human psychologist. Woebot does not connect users to verified practitioners, does
not generate clinical summaries for a professional, and does not support the full
consultation lifecycle. Its conversations are structured around CBT exercises rather
than free-form psychological expression, which limits its ability to serve as a
general-purpose pre-session intake tool. Woebot does not support Arabic or Darija,
and it has no credential verification, booking, or session management capabilities.
It addresses only the conversational layer of the mental health workflow, leaving the
professional consultation entirely outside its scope.
 
\subsubsection{Wysa}
 
Wysa~\cite{Wysa2023} is a multilingual AI mental health platform available in English,
Hindi, and several other languages. It offers guided CBT and mindfulness exercises
through a conversational interface, and in some commercial configurations provides
a pathway to connect with human coaches or therapists. Clinical evidence for its
efficacy in reducing mild to moderate anxiety and depression symptoms has been
published in peer-reviewed literature, making it one of the more academically
grounded commercial tools in this space.
 
Wysa's limitations are structurally similar to those of Woebot. It does not generate
structured clinical summaries, does not support a full consultation lifecycle with
a verified psychologist, and does not provide session activation codes, real-time
messaging, or PDF report generation. Its connection to human practitioners is
described as a referral pathway rather than an integrated consultation workflow.
Wysa does not support Darija or Arabic RTL interfaces, and its AI pipeline is
not designed for RAG-based enrichment with domain-specific clinical documents.
 
\subsubsection{Replika}
 
Replika~\cite{Replika2023} is a conversational AI companion application designed for
emotional support and social interaction. It uses large language model inference to
maintain an ongoing personalised conversation with the user over time, building a
longitudinal interaction history.
 
Replika is included in this review not as a direct competitor but as an illustration
of the risks associated with general-purpose LLM companions in a mental health
context. The platform has no clinical guardrails, no psychologist involvement,
no credential verification, and no prohibition on producing content that could
be clinically harmful. Several published reports and community observations have
documented cases in which users in distress received responses that reinforced
rather than redirected their concerns. Replika demonstrates that access to a
capable LLM alone is insufficient for a responsible mental health application:
the system prompt, safety constraints, human-in-the-loop oversight, and clinical
knowledge grounding that PsychPlatform implements are not optional enhancements
but essential design requirements.
 
% -----------------------------------------------------------
\subsection{Structured Comparative Analysis}
\label{subsec:comparison-table}
% -----------------------------------------------------------
 
Table~\ref{tab:platform-comparison} presents a structured comparison of the five surveyed
platforms against PsychPlatform across seven criteria derived directly from the functional
requirements introduced in Chapter~1.
 
The criteria are defined as follows:
 
\begin{itemize}[leftmargin=1.5em, itemsep=4pt]
  \item[--] \textbf{AI-assisted pre-session intake}: Does the platform conduct a
    structured conversational intake that produces a clinical document before the
    session?
  \item[--] \textbf{Darija / dialectal NLP}: Does the platform support Tunisian
    Darija or similar North African Arabic dialects in its conversational pipeline?
  \item[--] \textbf{Psychologist credential verification}: Does the platform
    implement a credential verification workflow that validates practitioner
    qualifications before granting patient-facing access?
  \item[--] \textbf{Clinical summary generation}: Does the platform automatically
    produce a structured clinical briefing document for the psychologist?
  \item[--] \textbf{Full consultation lifecycle}: Does the platform cover the
    complete workflow from booking to session activation, real-time communication,
    report generation, and post-session rating?
  \item[--] \textbf{Multilingual with RTL}: Does the platform support Arabic with
    correct right-to-left layout rendering, alongside English and French?
  \item[--] \textbf{RAG over private documents}: Does the platform implement
    retrieval-augmented generation over private, permission-scoped clinical
    documents?
\end{itemize}
 
\begin{table}[H]
  \centering
  \caption{Comparative analysis of existing platforms and PsychPlatform across
           seven criteria derived from the project's functional requirements.
           \checkmark\ = fully supported, $\sim$ = partially supported,
           \texttimes\ = not supported.}
  \label{tab:platform-comparison}
  \renewcommand{\arraystretch}{1.5}
  \begin{tabularx}{\linewidth}{@{} l *{6}{>{\centering\arraybackslash}X}
                               >{\centering\arraybackslash}X @{}}
    \toprule
    \rowcolor{medgray}
    \textbf{Criterion}
      & \textbf{Better-Help}
      & \textbf{Doctolib}
      & \textbf{Woebot}
      & \textbf{Wysa}
      & \textbf{Replika}
      & \textbf{Psych-Platform} \\
    \midrule
    AI pre-session intake
      & \texttimes & \texttimes & $\sim$ & $\sim$ & \texttimes & \checkmark \\
    \rowcolor{lightgray}
    Darija / dialectal NLP
      & \texttimes & \texttimes & \texttimes & \texttimes & \texttimes & \checkmark \\
    Credential verification
      & $\sim$ & $\sim$ & \texttimes & \texttimes & \texttimes & \checkmark \\
    \rowcolor{lightgray}
    Clinical summary generation
      & \texttimes & \texttimes & \texttimes & \texttimes & \texttimes & \checkmark \\
    Full consultation lifecycle
      & $\sim$ & $\sim$ & \texttimes & \texttimes & \texttimes & \checkmark \\
    \rowcolor{lightgray}
    Multilingual with RTL
      & \texttimes & \texttimes & \texttimes & \texttimes & \texttimes & \checkmark \\
    RAG over private documents
      & \texttimes & \texttimes & \texttimes & \texttimes & \texttimes & \checkmark \\
    \bottomrule
  \end{tabularx}
  \smallskip
 
  \noindent\footnotesize
  Note: BetterHelp and Doctolib receive $\sim$ on credential verification because
  they perform practitioner vetting, but through manual or national-registry processes
  that are not applicable in the Tunisian context and do not produce an AI-assisted
  structured summary for administrative review. Woebot and Wysa receive $\sim$ on
  AI pre-session intake because they offer conversational interaction, but the output
  is not structured as a clinical document and is not transmitted to a psychologist.
\end{table}
 
% -----------------------------------------------------------
\subsection{Synthesis: Deriving PsychPlatform's Design Decisions from Observed Gaps}
\label{subsec:synthesis}
% -----------------------------------------------------------
 
The comparative analysis reveals a consistent pattern across all five surveyed platforms:
each addresses at most one layer of the digital mental health workflow in depth, while
leaving the surrounding infrastructure unimplemented. This fragmentation has a concrete
clinical consequence. A patient using Doctolib to book a session with a psychologist
arrives at the session with no structured preparation; the psychologist has no pre-session
intelligence and must spend the first minutes of a paid consultation reconstructing the
patient's concerns from scratch. A patient using Woebot or Wysa benefits from structured
conversational support but cannot transition that support into a professional consultation
— the two worlds remain disconnected. A patient using BetterHelp accesses a therapist
but in English only, through a static intake questionnaire, with no longitudinal emotional
tracking across sessions.
 
Three specific gaps are particularly consequential for the design of PsychPlatform,
and each maps directly to an architectural decision made in the present project.
 
\textbf{The absence of a structured clinical handoff between patient preparation and
professional consultation.} No surveyed platform transforms patient conversational data
into a structured, psychologist-readable clinical document before the session begins. This
gap motivates PsychPlatform's clinical summarisation pipeline: the five-field
\texttt{ChatbotSummary} document (dominant emotion, urgency score, sentiment trend,
key themes, and raw summary) is designed specifically to fill this handoff gap, giving
the psychologist a grounded briefing that replaces the unproductive first minutes of
a session.
 
\textbf{The complete absence of Darija and North African dialect support.} Every
surveyed platform is designed for linguistically homogeneous markets. None supports
Tunisian Darija, a spoken dialect of approximately twelve million speakers that is
linguistically distinct from Modern Standard Arabic and is not adequately handled by
general-purpose multilingual NLP systems. This gap motivates the dedicated Darija
retrieval index implemented in PsychPlatform's first RAG pipeline
(\texttt{RetrievePsychologicalContext.js}), which helps the chatbot better interpret
dialectal psychological expressions that would otherwise be misunderstood or ignored by a
general-purpose model, while still requiring further enrichment to reach robust coverage.
 
\textbf{The absence of an AI-assisted, locally applicable credential verification
workflow.} Both BetterHelp and Doctolib perform psychologist vetting, but through
processes that are either internal and opaque (BetterHelp) or dependent on national
health system identifiers that do not exist in Tunisia (Doctolib's RPPS system). Neither
produces an AI-generated structured briefing for the reviewing administrator. This gap
motivates PsychPlatform's OCR-and-LLM verification pipeline, which is designed to
function without dependency on external registries by extracting and analysing
qualification evidence directly from uploaded documents, while preserving human
oversight through a human-in-the-loop administrative approval step~\cite{Dellermann2019}.
 
In summary, the state of the art establishes that the individual technologies required
to build PsychPlatform are well-founded and already deployed at scale in adjacent
applications. What does not exist is a platform that integrates them into a unified,
clinically coherent, and linguistically inclusive workflow. PsychPlatform is designed
to close that gap.


\chapter{Sprint 1: Authentication, RBAC, and Initial Setup}
\label{chap:sprint1-auth}

\section*{Introduction}
\phantomsection
\addcontentsline{toc}{section}{Introduction}

Sprint~1 establishes the technical and security foundation of \emph{PsychPlatform}. It focuses on project initialization, database connectivity, user identity management, authentication, role-based access control, protected routes, and the first interface shells. Since the platform is intended to handle sensitive user information, this sprint deliberately limits its scope to the creation of a stable trust boundary before introducing psychologist onboarding, booking, messaging, or AI-assisted services in later iterations.

\section{Sprint Objectives and Scope}
\label{sec:sprint1-objectives}

This section defines the purpose and boundaries of Sprint~1. The sprint is not intended to deliver advanced business functionality; rather, it prepares the secure entry point required by all future modules of the platform.

\begin{figure}[H]
  \centering
  \includegraphics[width=0.84\textwidth]{assets/s1/objective1.jpg}
  \caption{Sprint~1 objectives and scope overview.}
  \label{fig:sprint1-objectives-scope}
\end{figure}

\subsection{Primary Objectives}

The primary objectives of Sprint~1 are centred on initialization, authentication, and authorization. They ensure that the platform can safely identify users and restrict their access according to their role.

\begin{itemize}[leftmargin=*]
  \item[--] Initialize the frontend and backend workspaces and establish the base project structure.
  \item[--] Configure MongoDB connectivity and define the initial \texttt{User} schema through Mongoose.
  \item[--] Implement registration, login, email verification, and authenticated profile retrieval using JWT-based authentication.
  \item[--] Enforce role-based access control for \emph{patient}, \emph{psychologist}, and \emph{admin} roles.
  \item[--] Create foundational interfaces for login, registration, email verification, password recovery, and role-aware dashboard entry pages.
\end{itemize}

\subsection{Out-of-Scope Elements}

To preserve the coherence of the sprint, several features are intentionally excluded from this iteration. Psychologist onboarding verification, search and booking, live sessions, messaging, notifications, document processing, and AI-assisted functionalities are postponed to later sprints. This limitation is methodologically justified, since these modules depend on a reliable authentication and authorization layer.

\subsection{Success Criteria}

The sprint is considered successful if the authentication workflow is functional, secure, and reusable by subsequent modules. The expected validation criteria are listed below.

\begin{itemize}[leftmargin=*]
  \item[--] A new user can register with a valid role and securely log in.
  \item[--] Passwords are never stored in plaintext and are verified through \texttt{bcrypt} comparison.
  \item[--] Protected routes reject unauthenticated requests and unauthorized role transitions.
  \item[--] The frontend preserves session state and redirects users to the appropriate dashboard shell according to role.
\end{itemize}

\section{Sprint Planning}
\label{sec:sprint1-planning}

Sprint planning clarifies the implementation priorities and anticipates the risks associated with the first technical iteration. In this sprint, the main challenge is to define a clean foundation that remains extensible for the following modules.

\subsection{Sprint Context}

Sprint~1 is the first implementation sprint of the project. Its contribution is mainly infrastructural: it prepares the project structure, authentication services, database connection, and access-control mechanisms. Although this sprint provides limited visible business value, it has a decisive impact on the maintainability and security of the whole platform.

\subsection{Planned Deliverables}

The planned deliverables focus on the minimum components required to authenticate users and protect platform access.

\begin{itemize}[leftmargin=*]
  \item[--] Base MERN repository structure with separated client and server layers.
  \item[--] MongoDB connection configuration and initial environment setup.
  \item[--] \texttt{User} model with role field and password hashing hook.
  \item[--] Authentication API endpoints for registration, login, email verification, and authenticated profile retrieval.
  \item[--] Frontend authentication context, protected routing, and initial login, registration, verification, and recovery page shells.
\end{itemize}

\subsection{Dependencies and Risks}

Although Sprint~1 is functionally limited, it involves important architectural risks. Weak session handling, incomplete backend route protection, inconsistent role naming between the client and the server, or premature coupling between authentication and future business modules could affect later development. To mitigate these risks, the implementation is based on simple contracts, centralized middleware, explicit role checks, and a clear separation between authentication logic and domain-specific services.

\section{Sprint Backlog}
\label{sec:sprint1-backlog}

The sprint backlog presented in Table~\ref{tab:sprint1-backlog} summarizes the development items selected for Sprint~1. Each item contributes to the construction of the authentication and platform initialization layer.

\begin{longtable}{@{} L{0.10\textwidth} L{0.18\textwidth} L{0.42\textwidth} C{0.10\textwidth} C{0.11\textwidth} @{}}
\caption{Sprint~1 backlog focused on authentication and platform initialization.}
\label{tab:sprint1-backlog}\\
\toprule
\textbf{ID} &
\textbf{Title} &
\textbf{Description} &
\textbf{Priority} &
\textbf{Status} \\
\midrule
\endfirsthead

\toprule
\textbf{ID} &
\textbf{Title} &
\textbf{Description} &
\textbf{Priority} &
\textbf{Status} \\
\midrule
\endhead

PB-01.1 & Project Initialization & Set up the frontend and backend workspaces, base dependencies, and development scripts. & High & Done \\
\midrule
PB-01.2 & Database Connection & Configure MongoDB connectivity and reusable server initialization. & High & Done \\
\midrule
PB-01.3 & User Model & Define the user schema with identity fields, role management, and password hashing hooks. & High & Done \\
\midrule
PB-01.4 & Auth Endpoints & Implement register, login, verify-email, and authenticated profile retrieval endpoints. & High & Done \\
\midrule
PB-01.5 & Frontend Auth State & Persist JWT session state, restore profile state, and attach tokens to protected requests. & High & Done \\
\midrule
PB-01.6 & Initial UI Shells & Create login, registration, email verification, recovery, and role-aware dashboard shell pages. & Medium & Done \\
\midrule
PB-01.7 & Route Protection & Restrict frontend and backend access according to authentication state, verification state, and role. & High & Done \\

\bottomrule
\end{longtable}

\section{Requirement Analysis}
\label{sec:sprint1-analysis}

The requirement analysis of Sprint~1 focuses on access management. It identifies the expectations of the main stakeholders and translates them into functional, non-functional, and business rules related to identity and authorization.

\subsection{Stakeholder Perspective}

From a stakeholder perspective, Sprint~1 addresses complementary needs. Patients and psychologists require a reliable and secure entry point into the platform, while administrators require clear privilege boundaries from the beginning of the system. For the development team, this sprint provides a reusable authentication kernel that prevents repeated security redesign in later iterations.

\subsection{Functional Analysis}

The functional requirements of Sprint~1 extend beyond basic registration and login. The platform must allow a user to create an account, verify it by email, authenticate with valid credentials, restore the authenticated profile after refresh, and navigate only to pages authorized for their role. These operations form the root dependency of all later workflows.

\begin{itemize}[leftmargin=*]
  \item[--] \emph{Registration}: collect essential identity data, validate the selected role, ensure email uniqueness, hash the password, and create a new user document.
  \item[--] \emph{Email verification}: validate the verification token or code and mark the account as verified before normal platform usage.
  \item[--] \emph{Login}: validate credentials, generate a JWT, and return the authenticated user context required by the client.
  \item[--] \emph{Authenticated profile retrieval}: expose an authenticated endpoint used by the client to restore state after page refresh.
  \item[--] \emph{Password recovery scaffolding}: provide the initial interface and request flow needed for later recovery hardening.
  \item[--] \emph{Role-based routing}: prevent users from accessing routes that do not correspond to their assigned role.
\end{itemize}

\subsection{Non-Functional Requirements}

The non-functional requirements of this sprint are mainly related to security, scalability, usability, and maintainability. They define the quality level expected from the authentication foundation.

\begin{itemize}[leftmargin=*]
  \item[--] \emph{Security}: passwords must be hashed before persistence; token verification must protect sensitive endpoints; role checks must be enforced server-side.
  \item[--] \emph{Scalability}: the authentication design must remain reusable by future modules without major refactoring.
  \item[--] \emph{Usability}: login and registration interfaces must remain clear and consistent, even at this early stage.
  \item[--] \emph{Maintainability}: concerns must be separated between models, controllers, middleware, and client state management.
\end{itemize}

\subsection{Business Rules}

The business rules of Sprint~1 define the constraints that govern identity and access control. These rules remain valid for later sprints, since all business modules depend on authenticated and authorized access.

\begin{itemize}[leftmargin=*]
  \item[--] Every account must be associated with exactly one valid role.
  \item[--] Email addresses are unique across all users.
  \item[--] Email verification state must be persisted and checked where required.
  \item[--] A user without a valid JWT cannot access protected resources.
  \item[--] A valid token does not automatically imply universal access; role-based constraints remain applicable.
\end{itemize}

\section{Architecture Overview}
\label{sec:sprint1-architecture}

The architectural view of Sprint~1 explains how authentication responsibilities are distributed across the frontend, backend, and database layers. The design remains intentionally simple in order to limit risk and prepare later extensions.

\subsection{Logical Architecture}

Sprint~1 relies on a three-layer web architecture composed of a React frontend, an Express application server, and a MongoDB database. This separation is particularly important because each layer corresponds to a different trust boundary. The client manages interface rendering and session presentation, the backend verifies identity and enforces authorization, and the database persists user information and role-related state.

From an architectural perspective, the sprint is intentionally limited. Its main contribution is not the number of implemented modules, but the order and correctness of the authentication responsibilities. Behavioral diagrams are therefore introduced later in the chapter after the structural and functional views.

\subsection{Module Interaction}

At this stage, module interaction remains focused on authentication. The authentication controller coordinates input validation, password verification, JWT signing, and response generation. Middleware decodes tokens and guards protected routes, while the frontend authentication context propagates authenticated state across protected pages. This separation helps prevent duplicated security logic and improves maintainability.

\section{Database Design}
\label{sec:sprint1-database}

The database design of Sprint~1 focuses on the \texttt{User} model. This model is intentionally concise, as its purpose is to support authentication, email verification, and role-based navigation without introducing premature domain-specific fields.

\subsection{User Model}

The \texttt{User} entity is the central persistence artifact of Sprint~1. It is designed to remain stable across future sprints and therefore includes the attributes required to identify a user and support role-sensitive access.

\begin{longtable}{@{} L{0.22\textwidth} L{0.16\textwidth} L{0.14\textwidth} L{0.38\textwidth} @{}}
\caption{Core \texttt{User} model used in Sprint~1.}
\label{tab:sprint1-user-model}\\
\toprule
\textbf{Attribute} &
\textbf{Type} &
\textbf{Constraint} &
\textbf{Purpose} \\
\midrule
\endfirsthead

\toprule
\textbf{Attribute} &
\textbf{Type} &
\textbf{Constraint} &
\textbf{Purpose} \\
\midrule
\endhead

\texttt{fullName} & String & Required & Human-readable identity displayed by the application. \\
\midrule
\texttt{email} & String & Required, unique & Primary authentication identifier. \\
\midrule
\texttt{password} & String & Required & Hashed password stored after \texttt{bcrypt} processing. \\
\midrule
\texttt{role} & Enum & Required & Role discriminator: patient, psychologist, or admin. \\
\midrule
\texttt{isVerified} & Boolean & Default false & Indicates whether the email verification flow has been completed. \\
\midrule
\texttt{createdAt} & Date & Automatic & Audit-oriented creation timestamp. \\
\midrule
\texttt{updatedAt} & Date & Automatic & Traceability of profile updates. \\

\bottomrule
\end{longtable}

\subsection{Security Considerations}

The database design is intentionally conservative. Sensitive credentials are not stored directly; only password hashes are persisted. Unique email constraints prevent account duplication, while role values are restricted to a predefined enumeration in order to avoid inconsistent privilege assignments.

To complement the tabular description of the data model, Figure~\ref{fig:sprint1-class} formalizes the relations between the user entity and the authentication components that process it.

\begin{figure}[H]
  \centering
  \includegraphics[width=0.92\textwidth]{assets/s1/d1.png}
  \caption{Class diagram associated with Sprint~1 authentication logic.}
  \label{fig:sprint1-class}
\end{figure}

The diagram shows that the \texttt{User} model remains the persistence centrepiece, while the controller, service, and middleware layers implement distinct responsibilities within the authentication pipeline.

\section{API Layer and Security Mechanisms}
\label{sec:sprint1-api}

The API layer provides the backend contract required for account creation, authentication, email verification, and identity restoration. Its surface remains concise in order to reduce complexity and improve security verification.

\subsection{Authentication Endpoints}

The authentication endpoints introduced in Sprint~1 are summarized in Table~\ref{tab:sprint1-endpoints}. Each endpoint corresponds to a specific stage of the account lifecycle.

\begin{longtable}{@{} L{0.28\textwidth} L{0.12\textwidth} L{0.52\textwidth} @{}}
\caption{Authentication endpoints introduced in Sprint~1.}
\label{tab:sprint1-endpoints}\\
\toprule
\textbf{Endpoint} &
\textbf{Method} &
\textbf{Purpose} \\
\midrule
\endfirsthead

\toprule
\textbf{Endpoint} &
\textbf{Method} &
\textbf{Purpose} \\
\midrule
\endhead

\texttt{/api/auth/register} & POST & Create a new account after validation, password hashing, and verification preparation. \\
\midrule
\texttt{/api/auth/login} & POST & Authenticate an existing user and issue a JWT. \\
\midrule
\path{/api/auth/verify-email} & POST & Validate the verification token or code and activate the account. \\
\midrule
\texttt{/api/auth/profile} & GET & Return the authenticated user context from the provided token. \\

\bottomrule
\end{longtable}

\subsection{Route Protection Strategy}

Route protection is enforced at two levels. On the backend, authentication middleware verifies the JWT and enriches the request with the decoded user identity before role middleware authorizes access. On the frontend, protected routes prevent accidental navigation to unauthorized pages and maintain a coherent user experience. Nevertheless, backend protection remains the authoritative security control because client-side restrictions can be bypassed.

\subsection{Security Controls Summary}

Table~\ref{tab:sprint1-security-controls} summarizes the security controls introduced during Sprint~1 and explains their expected effect on the platform.

\begin{longtable}{@{} L{0.24\textwidth} L{0.44\textwidth} L{0.24\textwidth} @{}}
\caption{Security controls introduced during Sprint~1.}
\label{tab:sprint1-security-controls}\\
\toprule
\textbf{Control} &
\textbf{Implementation Role} &
\textbf{Expected Effect} \\
\midrule
\endfirsthead

\toprule
\textbf{Control} &
\textbf{Implementation Role} &
\textbf{Expected Effect} \\
\midrule
\endhead

Password hashing & \texttt{bcrypt} pre-save or service-layer hashing & Prevent plaintext credential disclosure. \\
\midrule
Email verification & Verification endpoint and account-state flag & Prevent immediate use of unverified accounts. \\
\midrule
JWT verification & Authentication middleware & Reject forged or missing sessions. \\
\midrule
Role authorization & Dedicated authorization middleware & Prevent cross-role access to protected routes. \\
\midrule
Profile endpoint & Authenticated identity restoration & Preserve consistent client state after refresh. \\

\bottomrule
\end{longtable}

\section{Functional Specification}
\label{sec:sprint1-functional-spec}

The functional specification models Sprint~1 as an access-control sprint. The use case diagram shown in Figure~\ref{fig:sprint1-usecase} groups the main interactions around account creation, account activation, authentication, and protected access.

\begin{figure}[H]
  \centering
  \includegraphics[width=0.92\textwidth]{assets/s1/d2.png}
  \caption{Overall use case diagram for Sprint~1.}
  \label{fig:sprint1-usecase}
\end{figure}

The diagram emphasizes that every privileged action depends on prior authentication and that authorization is role-sensitive rather than globally uniform.

\subsection{Use Case Descriptions}

The following tables describe the main use cases implemented during Sprint~1. They clarify actors, goals, preconditions, success scenarios, and postconditions.

\begin{longtable}{@{} L{0.28\textwidth} L{0.66\textwidth} @{}}
\caption{Use case specification for user registration.}
\label{tab:uc-s1-register}\\
\toprule
\multicolumn{2}{@{}l@{}}{\textbf{UC-01 \quad Register a New User}} \\
\midrule
\endfirsthead

\toprule
\multicolumn{2}{@{}l@{}}{\textbf{UC-01 \quad Register a New User}} \\
\midrule
\endhead

Primary Actor & Visitor \\
\midrule
Goal & Create a new patient, psychologist, or admin account with valid credentials. \\
\midrule
Preconditions & The email address is not already associated with another account. \\
\midrule
Main Success Scenario & The user fills in the registration form; the backend validates inputs; the password is hashed; the user document is created; a successful response is returned to the client. \\
\midrule
Postconditions & A new user exists in the database with a hashed password and an assigned role. \\

\bottomrule
\end{longtable}

\begin{longtable}{@{} L{0.28\textwidth} L{0.66\textwidth} @{}}
\caption{Use case specification for email verification and account activation.}
\label{tab:uc-s1-verify}\\
\toprule
\multicolumn{2}{@{}l@{}}{\textbf{UC-02 \quad Verify Email and Activate Account}} \\
\midrule
\endfirsthead

\toprule
\multicolumn{2}{@{}l@{}}{\textbf{UC-02 \quad Verify Email and Activate Account}} \\
\midrule
\endhead

Primary Actor & Registered user \\
\midrule
Goal & Confirm account ownership and activate the account through the email verification flow. \\
\midrule
Preconditions & The account exists and a valid verification token or code has been issued. \\
\midrule
Main Success Scenario & The user opens the verification link or enters the verification code; the backend validates the token; the account verification flag is updated; the client is redirected to the authenticated flow. \\
\midrule
Alternate Flow & If the verification token is invalid or expired, the system rejects activation and prompts a retry or resend flow. \\

\bottomrule
\end{longtable}

\begin{longtable}{@{} L{0.28\textwidth} L{0.66\textwidth} @{}}
\caption{Use case specification for login.}
\label{tab:uc-s1-login}\\
\toprule
\multicolumn{2}{@{}l@{}}{\textbf{UC-03 \quad Authenticate User}} \\
\midrule
\endfirsthead

\toprule
\multicolumn{2}{@{}l@{}}{\textbf{UC-03 \quad Authenticate User}} \\
\midrule
\endhead

Primary Actor & Registered and verified user \\
\midrule
Goal & Access the platform through valid credentials and obtain an authenticated session token. \\
\midrule
Preconditions & The account exists, is verified when policy requires it, and the submitted password matches the stored hash. \\
\midrule
Main Success Scenario & The user submits email and password; the backend retrieves the account; \texttt{bcrypt.compare()} verifies the password; a JWT is signed and returned with the user payload. \\
\midrule
Alternate Flow & If credentials are invalid or the account is not yet activated, the system returns an authentication error and no token is issued. \\

\bottomrule
\end{longtable}

\begin{longtable}{@{} L{0.28\textwidth} L{0.66\textwidth} @{}}
\caption{Use case specification for route protection.}
\label{tab:uc-s1-protected}\\
\toprule
\multicolumn{2}{@{}l@{}}{\textbf{UC-04 \quad Access a Protected Route}} \\
\midrule
\endfirsthead

\toprule
\multicolumn{2}{@{}l@{}}{\textbf{UC-04 \quad Access a Protected Route}} \\
\midrule
\endhead

Primary Actor & Authenticated user \\
\midrule
Goal & Access a page or endpoint that is limited to a given role. \\
\midrule
Preconditions & A valid JWT is provided and the authenticated user's role satisfies the route policy. \\
\midrule
Main Success Scenario & Middleware verifies the token, attaches the user to the request, checks role authorization, and grants access to the protected resource. \\
\midrule
Alternate Flow & If the token is missing, expired, or incompatible with the route policy, the request is rejected. \\

\bottomrule
\end{longtable}

\subsection{Workflow Description}

\begin{figure}[H]
  \centering
  \includegraphics[width=0.92\textwidth]{assets/s1/register.png}
  \caption{Sequence diagram of registration in Sprint~1.}
  \label{fig:sprint1-sequence-registration}
\end{figure}

\begin{figure}[H]
  \centering
  \includegraphics[width=0.92\textwidth]{assets/s1/EmailVerification.png}
  \caption{Sequence diagram of the email verification.}
  \label{fig:sprint1-sequence-emailverif}
\end{figure}

\begin{figure}[H]
  \centering
  \includegraphics[width=0.92\textwidth]{assets/s1/LoginSequence.png}
  \caption{Sequence diagram of the login.}
  \label{fig:sprint1-sequence-login}
\end{figure}
\begin{figure}[H]
  \centering
  \includegraphics[width=0.92\textwidth]{assets/s1/profileRestore.png}
  \caption{Sequence diagram of the profile restoring.}
  \label{fig:sprint1-sequence-restore}
\end{figure}

\begin{figure}[H]
  \centering
  \includegraphics[width=0.92\textwidth]{assets/s1/protectedroute.png}
  \caption{Sequence diagram of Accessing a Protected Route.}
  \label{fig:sprint1-sequence-ProtectedRoute}
\end{figure}
\section{Interface Description}
\label{sec:sprint1-ui}

From a visual standpoint, Sprint~1 introduces only foundational interfaces. Their role is to support secure navigation and validate role-aware access rather than deliver complete business functionality.

\begin{itemize}[leftmargin=*]
  \item[--] \emph{Registration page}: captures identity, email, password, and selected role.
  \item[--] \emph{Email verification page}: confirms account activation through a token or verification code.
  \item[--] \emph{Login page}: handles credential submission and error feedback.
  \item[--] \emph{Password recovery page}: provides the initial recovery request interface.
  \item[--] \emph{Dashboard shells}: display basic role-aware entry points after authentication.
\end{itemize}

\subsection{Interface Screenshot}
\begin{figure}[H]
  \centering
  \includegraphics[width=0.82\textwidth]{assets/interfaces/interface1.png}
  \caption{Sprint~1 login interface.}
  \label{fig:sprint1-login-ui}
\end{figure}

\begin{figure}[H]
  \centering
  \includegraphics[width=0.82\textwidth]{assets/interfaces/interface2.png}
  \caption{Sprint~1 register interfaces.}
  \label{fig:sprint1-verification-ui}
\end{figure}

\begin{figure}[H]
  \centering
  \includegraphics[width=0.82\textwidth]{assets/interfaces/interface3.png}
  \caption{Sprint~1 recovery account}
  \label{fig:sprint1-dashboard-ui}
\end{figure}
\section{Validation Strategy}
\label{sec:sprint1-validation}

Validation focuses on the correctness of authentication logic, the robustness of route protection, and the coherence of client-server integration. The scenarios listed in Table~\ref{tab:sprint1-tests} are representative of the minimum controls required before moving to more complex modules.

\begin{longtable}{@{} L{0.14\textwidth} L{0.56\textwidth} L{0.22\textwidth} @{}}
\caption{Representative validation scenarios for Sprint~1.}
\label{tab:sprint1-tests}\\
\toprule
\textbf{Test ID} &
\textbf{Scenario} &
\textbf{Expected Result} \\
\midrule
\endfirsthead

\toprule
\textbf{Test ID} &
\textbf{Scenario} &
\textbf{Expected Result} \\
\midrule
\endhead

T1 & Register a new patient with valid data. & User is created and password is stored as a hash. \\
\midrule
T2 & Attempt registration with an existing email. & API rejects the request and preserves uniqueness. \\
\midrule
T3 & Verify account with valid token or code. & Verification state is updated successfully. \\
\midrule
T4 & Login with valid verified credentials. & JWT is returned and client state becomes authenticated. \\
\midrule
T5 & Login with invalid password or invalid verification state. & Authentication fails without token creation. \\
\midrule
T6 & Access protected route without token. & Request is rejected and the client is redirected. \\
\midrule
T7 & Access admin route with non-admin role. & Authorization middleware denies access. \\

\bottomrule
\end{longtable}

\subsection{Deployment Notes}

Sprint~1 was executed in a development environment using local frontend and backend instances, with automatic reload for rapid iteration. This setup was sufficient for validating integration between client routing, API contracts, and database persistence before introducing more complex infrastructure or production-oriented deployment constraints.

\begin{comment}
\section{Discussion and Limitations}
\label{sec:sprint1-discussion}

Sprint~1 deliberately prioritizes correctness over breadth. As a result, the interface remains visually basic and the sprint intentionally stops at the first authentication boundary. In practical terms, Sprint~1 does not yet address the broader account lifecycle expected in a production platform, such as password recovery, long-lived session management, audit-oriented account history, or stronger verification and abuse-prevention flows. In addition, dashboard shells introduced in this sprint are intentionally structural rather than functional: they confirm role segregation, but do not yet expose the specialized business tools delivered in later sprints. These limitations are consistent with the sprint goal, which is to establish a reliable JWT- and RBAC-based foundation before introducing more complex workflows.
\end{comment}

\section{Sprint Review and Retrospective}
\label{sec:sprint1-review}

Sprint~1 delivered the expected authentication baseline and reduced technical uncertainty for the rest of the project. The principal lesson from this sprint is that authorization rules must be treated first as a backend responsibility and only second as a frontend navigation convenience. Early investment in a stable user model, reusable middleware, and explicit role checks simplifies the implementation of subsequent sprints and reduces the risk of inconsistent access control.

\section{Conclusion}
\label{sec:sprint1-conclusion}

Sprint~1 provides the minimum secure foundation upon which the remaining modules of PsychPlatform are built. By implementing authentication, role-based access control, database initialization, and protected interface shells, the sprint establishes the technical conditions required for onboarding, booking, communication, and later AI-assisted services. Although functionally modest, it is architecturally decisive because it defines the platform's first trust boundary.
\clearpage
\chapter{Sprint 2: Psychologist Onboarding System}
\label{chap:sprint2-onboarding}

\section*{Introduction}
\phantomsection
\phantomsection
\addcontentsline{toc}{section}{Introduction}

Sprint~2 focuses on the design and implementation of a structured onboarding
pipeline for psychologists. In a mental-health platform, onboarding is more
than a registration form: it is a security boundary and a trust mechanism. The
platform must ensure that only verified and qualified professionals are able to
offer services, appear in search results, and access sensitive patient-facing
features.

This sprint introduces an end-to-end onboarding workflow covering: (i) profile
creation and completion, (ii) secure submission of professional credentials,
(iii) AI-assisted pre-screening to detect obvious inconsistencies or missing
elements, and (iv) an administrative review process that produces a final
approve/reject decision. A key requirement is auditability: every important
transition (submission, verification, admin decision) is recorded with
timestamps and actor identity to guarantee traceability and accountability.

\section{Sprint Objectives and Scope}
\label{sec:sprint2-objectives}

The main objective of Sprint~2 is to deliver a complete onboarding system that
transitions a psychologist from an initial account state to an \emph{Active}
state only after verification. The scope includes both user-facing components
and internal operational tooling:

\begin{figure}[H]
  \centering
  \includegraphics[width=0.84\textwidth]{assets/sprint2/objective2.jpg}
  \caption{Sprint~2 objectives and scope overview.}
  \label{fig:sprint2-objectives-scope}
\end{figure}

\begin{itemize}
    \item Collect comprehensive professional information (identity, practice
    details, specialisations, and experience).
    \item Enable secure upload and management of credential documents.
    \item Integrate an AI-assisted verification step that produces a structured
    result (risk score, extracted entities, and flags).
    \item Provide an admin review queue and a decision workflow with explicit
    reasoning and audit logging.
    \item Enforce access control so that only approved psychologists can use
    privileged features and become discoverable on the platform.
\end{itemize}

\textbf{Out of scope.} Sprint~2 does not attempt to fully automate professional
licence validation against external government registries (where such APIs are
often unavailable or inconsistent). Instead, the system is designed so that AI
pre-screening supports (but does not replace) human administrative judgement.

\textbf{Success Criteria.}
\begin{itemize}
    \item A psychologist can complete onboarding without admin intervention
    until the final approval step.
    \item Admin decisions are fully traceable (who, when, what, and why).
    \item The onboarding status model has clear state transitions and prevents
    invalid moves (e.g., approving an incomplete submission).
    \item Uploaded credentials are stored securely and only accessible to
    authorised roles.
\end{itemize}

\section{Sprint Planning}
\label{sec:sprint2-planning}

\subsection{Assumptions}
\begin{itemize}
    \item Sprint duration: 3 weeks.
    \item Agile team structure: Product Owner, backend developer(s), frontend
    developer(s), QA, and an admin representative for workflow validation.
    \item Security and privacy constraints are treated as first-class
    requirements (least privilege, encrypted storage, and audit logging).
\end{itemize}

\subsection{Deliverables}
\begin{itemize}
    \item Onboarding UI (profile form, document checklist, submission and status
    tracking).
    \item Backend APIs for onboarding records, document upload metadata, and
    status transitions.
    \item AI verification pipeline integration (job submission, result
    persistence, and failure handling).
    \item Admin dashboard (review queue, submission detail view, decision
    actions, and audit log view).
    \item Testing pipeline covering validation, permissions, and end-to-end
    onboarding scenarios.
\end{itemize}

\section{Stakeholders and Roles}
\label{sec:sprint2-stakeholders}

Sprint~2 formalises responsibilities across three main actors:

\begin{itemize}
    \item \textbf{Psychologist:} provides profile data and credentials, submits
    onboarding request, and monitors status updates and feedback.
    \item \textbf{Admin:} reviews submissions, interprets AI verification
    results, and approves or rejects onboarding with recorded reasoning.
    \item \textbf{AI Verification Service:} analyses submitted documents and
    returns structured evidence to support admin decisions (risk score, flags,
    and extracted fields).
\end{itemize}

\section{Sprint Backlog}
\label{sec:sprint2-backlog}

\renewcommand{\arraystretch}{1.4}

\subsection{User Stories (Onboarding Flow)}
\begin{table}[H]
\centering
\begin{tabularx}{\textwidth}{|c|X|X|}
\hline
\rowcolor{medgray}
\textbf{ID} & \textbf{Description} & \textbf{Acceptance Criteria} \\
\midrule
\hline

US-2.1 & Creation of professional profile &
\begin{itemize}[leftmargin=*]
\item Validate required fields (identity, contact, location, specialisation).
\item Allow saving as draft and resuming later.
\item Persist a clear ``completion'' indicator to guide the user.
\end{itemize} \\ \hline

US-2.2 & Upload credential documents &
\begin{itemize}[leftmargin=*]
\item Validate file type and size; reject unsupported uploads.
\item Categorise documents (licence, degree, ID, other).
\item Confirm secure storage and display upload status per document.
\end{itemize} \\ \hline

US-2.3 & Submit onboarding request &
\begin{itemize}[leftmargin=*]
\item Block submission if required fields/documents are missing.
\item Assign onboarding status and store submission timestamp.
\item Trigger AI verification and store a job reference.
\end{itemize} \\ \hline

US-2.4 & Track onboarding status &
\begin{itemize}[leftmargin=*]
\item Display status and next steps clearly (what is pending and why).
\item If rejected, show a human-readable rejection reason.
\item Provide guidance for resubmission when applicable.
\end{itemize} \\ \hline

\end{tabularx}
\caption{User Stories: Onboarding Flow}
\end{table}

\subsection{AI and Admin Workflow}
\begin{table}[H]
\centering
\begin{tabularx}{\textwidth}{|c|X|X|}
\hline
\rowcolor{medgray}
\textbf{ID} & \textbf{Feature} & \textbf{Description} \\
\midrule
\hline

US-2.5 & AI Verification & AI analyses documents and returns risk score, extracted fields, and flags. \\ \hline
US-2.6 & Admin Queue & Displays pending onboarding requests with filtering and prioritisation. \\ \hline
US-2.7 & Admin Decision & Admin approves/rejects with mandatory logging, reason, and timestamp. \\ \hline

\end{tabularx}
\caption{User Stories: AI and Admin Workflow}
\end{table}

\subsection{Technical Tasks}
\begin{table}[H]
\centering
\begin{tabularx}{\textwidth}{|X|X|}
\hline
\rowcolor{medgray}
\textbf{Task} & \textbf{Description} \\
\midrule
\hline

State Machine & Define onboarding states, transitions, and invariants. \\ \hline
Secure Storage & Implement encrypted credential storage and controlled access paths. \\ \hline
AI Integration & Integrate verification pipeline (async jobs, retries, persistence). \\ \hline
Admin Dashboard & Build moderation interface and decision endpoints with audit trail. \\ \hline
Testing & Add unit, integration, and end-to-end tests for onboarding scenarios. \\ \hline

\end{tabularx}
\caption{Technical Tasks}
\end{table}

% ============================================================
%  SPRINT 2 — FUNCTIONAL SPECIFICATION WITH USE CASE SUPPORT
% ============================================================

% ─────────────────────────────────────────────────────────────
\section{Functional Specification}
\label{sec:sprint2-functional}

% ─────────────────────────────────────────────────────────────
\subsection{Profile Data Collection}
\label{sec:uc-profile}

The onboarding profile is designed as a structured record with both required
and optional fields. Required fields focus on identity and professional
credibility; optional fields enhance matching and discovery:

\begin{itemize}
    \item \textbf{Identity and contact:} full name, verified email/phone,
    country/city, profile photo (optional, subject to moderation policy).
    \item \textbf{Professional details:} licence/registration number (when
    applicable), issuing authority, specialisations, spoken languages, years of
    experience, and short professional biography.
    \item \textbf{Operational preferences:} availability windows, online/offline
    capability, consultation types offered, and pricing policy (if the platform
    supports it).
\end{itemize}

Client-side and server-side validation rules are aligned to prevent inconsistent
states and to ensure the backend remains the source of truth. Draft saving is
supported to reduce onboarding abandonment and allow incremental completion.

% ── Textual Use Case Descriptions ────────────────────────────
\subsubsection*{Use Case Descriptions — Profile Data Collection}
\begin{table}[H]
\centering
\renewcommand{\arraystretch}{1.4}
\begin{tabularx}{\textwidth}{|l|X|}
\hline
\rowcolor{medgray}
\textbf{Field} & \textbf{Description} \\ \hline

Use Case ID & UC-01 \\ \hline
Use Case Name & Fill Profile Form \\ \hline
Actor(s) & Psychologist (primary) \\ \hline
Preconditions & Psychologist account exists; onboarding session is active \\ \hline

Main Flow &
1. Open onboarding form \newline
2. System displays required and optional fields \newline
3. Psychologist fills required fields \newline
4. System validates inputs (client + server) \newline
5. System saves profile as Draft \\ \hline

Alternative Flow &
Save as Draft: User can save progress at any time without full validation \\ \hline

Exceptions &
Validation Failure: System highlights errors and blocks progression \\ \hline

Postconditions &
Profile stored in Draft state with required fields completed \\ \hline

\end{tabularx}
\caption{Use Case UC-01 — Fill Profile Form}
\end{table}
\begin{table}[H]
\centering
\renewcommand{\arraystretch}{1.4}
\begin{tabularx}{\textwidth}{|l|X|}
\hline
\rowcolor{medgray}
\textbf{Field} & \textbf{Description} \\ \hline

Use Case ID & UC-02 \\ \hline
Use Case Name & Upload Profile Photo \\ \hline
Actor(s) & Psychologist (primary), Moderation System (secondary) \\ \hline
Preconditions & Profile form is open \\ \hline

Main Flow &
1. Select photo file \newline
2. Validate format, size, resolution \newline
3. Queue for moderation \newline
4. Approve and attach to profile \\ \hline

Exceptions &
Moderation Rejection: User is notified and can re-upload \\ \hline

Postconditions &
Approved photo linked to profile \\ \hline

\end{tabularx}
\caption{Use Case UC-02 — Upload Profile Photo}
\end{table}

% ─────────────────────────────────────────────────────────────
\subsection{Credential Submission and Document Handling}
\label{sec:uc-credentials}

Credential submission is implemented as a checklist-driven flow. Each document
upload creates a record containing:

\begin{itemize}
    \item Document type (licence, degree, ID, other).
    \item Original filename, MIME type, file size, and upload timestamp.
    \item Storage reference (e.g., encrypted object key) and integrity metadata
    (e.g., checksum) when supported.
    \item Visibility constraints enforcing that only the psychologist (for their
    own submission) and admins can access documents.
\end{itemize}

To maintain privacy and limit exposure of sensitive documents, the platform
should prefer short-lived, role-scoped access mechanisms (e.g., signed URLs or
streaming through the backend) rather than permanent public links.

% ── Textual Use Case Descriptions ────────────────────────────
\subsubsection*{Use Case Descriptions — Credential Submission and Document Handling}

\begin{table}[H]
\centering
\renewcommand{\arraystretch}{1.4}
\begin{tabularx}{\textwidth}{|l|X|}
\hline
\rowcolor{medgray}
\textbf{Field} & \textbf{Description} \\ \hline

Use Case ID & UC-03 \\ \hline
Use Case Name & Upload Credential Document \\ \hline
Actor(s) & Psychologist (primary), Storage System (secondary) \\ \hline
Preconditions & Profile in Draft; user authenticated \\ \hline

Main Flow &
1. Select document type \newline
2. Upload file \newline
3. Validate format and size \newline
4. Generate storage reference + checksum \newline
5. Store document with restricted access \newline
6. Mark checklist item complete \\ \hline

Alternative Flow &
Replace Document: New version created while preserving history \\ \hline

Exceptions &
Unsupported file: Upload rejected with constraints message \\ \hline

Postconditions &
Document securely stored and accessible only to authorized roles \\ \hline

\end{tabularx}
\caption{Use Case UC-03 — Upload Credential Document}
\end{table}

\begin{table}[H]
\centering
\renewcommand{\arraystretch}{1.4}
\begin{tabularx}{\textwidth}{|l|X|}
\hline
\rowcolor{medgray}
\textbf{Field} & \textbf{Description} \\ \hline

Use Case ID & UC-04 \\ \hline
Use Case Name & Access Document (Scoped) \\ \hline
Actor(s) & Psychologist (own documents), Admin (all documents) \\ \hline

Preconditions & Document record exists in storage \\ \hline

Main Flow &
1. Actor requests document access \newline
2. System verifies role-based permissions \newline
3. System generates a short-lived, scoped access URL \newline
4. Actor retrieves the document within validity window \\ \hline

Exceptions &
Unauthorised Access: System denies access and logs the event \\ \hline

Postconditions &
Document delivered securely; access event recorded in audit log \\ \hline

\end{tabularx}
\caption{Use Case UC-04 — Access Document (Scoped)}
\end{table}

% ─────────────────────────────────────────────────────────────
\subsection{Submission, Locking, and Resubmission Policy}
\label{sec:uc-submission}

Upon submission, the system validates completeness and transitions the record
into an immutable review state (policy-dependent). A common approach is to lock
critical fields and documents to avoid post-submission tampering; any changes
require explicit resubmission and generate a new audit event.

If rejected, the record includes an admin-provided reason and can either:
(i) allow edits and resubmission while preserving the historical audit trail, or
(ii) require a new onboarding attempt linked to the same user. Sprint~2 targets
the first option because it minimises friction while preserving traceability.

% ── Textual Use Case Descriptions ────────────────────────────
\subsubsection*{Use Case Descriptions — Submission, Locking, and Resubmission}

\begin{table}[H]
\centering
\renewcommand{\arraystretch}{1.4}
\begin{tabularx}{\textwidth}{|l|X|}
\hline
\rowcolor{medgray}
\textbf{Field} & \textbf{Description} \\ \hline

Use Case ID & UC-05 \\ \hline
Use Case Name & Submit Onboarding Application \\ \hline
Actor(s) & Psychologist (primary), System (secondary) \\ \hline
Preconditions & All required data and documents completed \\ \hline

Main Flow &
1. Trigger submission \newline
2. Validate completeness \newline
3. Lock critical fields \newline
4. Change state to Submitted \newline
5. Log event \newline
6. Notify admin \\ \hline

Exceptions &
Incomplete Application: Submission blocked \\ \hline

Postconditions &
Application in Submitted state, immutable \\ \hline

\end{tabularx}
\caption{Use Case UC-05 — Submit Onboarding Application}
\end{table}

\begin{table}[H]
\centering
\renewcommand{\arraystretch}{1.4}
\begin{tabularx}{\textwidth}{|l|X|}
\hline
\rowcolor{medgray}
\textbf{Field} & \textbf{Description} \\ \hline

Use Case ID & UC-06 \\ \hline
Use Case Name & Edit and Resubmit After Rejection \\ \hline
Actor(s) & Psychologist (primary) \\ \hline

Preconditions & Record is in \textit{Rejected} state; rejection reason is available \\ \hline

Main Flow &
1. Psychologist reviews rejection reason \newline
2. System unlocks editable fields and document slots \newline
3. Psychologist makes corrections \newline
4. Psychologist resubmits \newline
5. System transitions record to \textit{Submitted} and logs a new audit event \\ \hline

Postconditions &
Updated record re-enters review queue; full audit trail preserved \\ \hline

\end{tabularx}
\caption{Use Case UC-06 — Edit and Resubmit After Rejection}
\end{table}

% ============================================================
\section{AI-Assisted Verification}
\label{sec:sprint2-ai}

The AI verification step is designed as a pre-screening mechanism that provides
evidence to admins. Typical checks include:

\begin{itemize}
    \item \textbf{OCR and entity extraction:} extract name, licence number,
    issuing body, and validity dates.
    \item \textbf{Consistency checks:} compare extracted entities against
    profile fields (e.g., name match, country/authority plausibility).
    \item \textbf{Completeness checks:} detect missing required document types
    or unreadable uploads (low OCR confidence).
    \item \textbf{Risk scoring:} compute an overall risk score with supporting
    flags to guide review prioritisation.
\end{itemize}

Because AI systems can produce false positives and false negatives, the output
is treated as \emph{advisory} rather than determinative. The admin interface
must clearly separate factual extracted fields from risk heuristics and provide
enough context to justify the final decision.

% ── Textual Use Case Descriptions ────────────────────────────
\subsubsection*{Use Case Descriptions — AI-Assisted Verification}

\begin{table}[H]
\centering
\renewcommand{\arraystretch}{1.4}
\begin{tabularx}{\textwidth}{|l|X|}
\hline
\rowcolor{medgray}
\textbf{Field} & \textbf{Description} \\ \hline

Use Case ID & UC-07 \\ \hline
Use Case Name & Run AI Pre-Screening \\ \hline
Actor(s) & AI System (primary), Admin (consumer) \\ \hline
Preconditions & Application is Submitted \\ \hline

Main Flow &
1. Trigger AI job \newline
2. Perform OCR \newline
3. Extract entities \newline
4. Cross-check data \newline
5. Evaluate completeness \newline
6. Compute risk score \newline
7. Move to Pending Review \\ \hline

Alternative Flow &
Low OCR Confidence: Flag unreadable document \\ \hline

Exceptions &
AI Failure: Retry then escalate \\ \hline

Postconditions &
AI report attached to application \\ \hline

\end{tabularx}
\caption{Use Case UC-07 — Run AI Pre-Screening}
\end{table}

% ============================================================
\section{Administrative Review and Auditability}
\label{sec:sprint2-admin}

% ─────────────────────────────────────────────────────────────
\subsection{Admin Review Queue}
\label{sec:uc-queue}

The admin queue provides operational visibility and prioritisation. Essential
capabilities include filtering by status (e.g., submitted, pending review,
flagged), sorting by submission date, and quickly opening a detailed view.

% ── Textual Use Case Descriptions ────────────────────────────
\subsubsection*{Use Case Descriptions — Admin Review Queue}

\begin{table}[H]
\centering
\renewcommand{\arraystretch}{1.4}
\begin{tabularx}{\textwidth}{|l|X|}
\hline
\rowcolor{medgray}
\textbf{Field} & \textbf{Description} \\ \hline

Use Case ID & UC-08 \\ \hline
Use Case Name & Browse and Filter Review Queue \\ \hline
Actor(s) & Admin (primary) \\ \hline

Preconditions & Admin is authenticated with review privileges \\ \hline

Main Flow &
1. Admin opens review queue dashboard \newline
2. System displays applications in \textit{Pending Admin Review}, sorted by submission date (latest first) \newline
3. Admin applies filters (status, date range, risk level) \newline
4. System updates the list dynamically \newline
5. Admin selects a record to view details \\ \hline

Postconditions &
Admin identifies and opens a target application for detailed review \\ \hline

\end{tabularx}
\caption{Use Case UC-08 — Browse and Filter Review Queue}
\end{table}

% ─────────────────────────────────────────────────────────────
\subsection{Decision Workflow}
\label{sec:uc-decision}

Admin decisions must be explicit and recorded. Approvals transition the
psychologist account to an \emph{Active} state and enable platform privileges.
Rejections require a reason and transition the record to \emph{Rejected}. The
system records:

\begin{itemize}
    \item Admin identity and timestamp.
    \item Decision type (approve/reject).
    \item Rationale text and, when applicable, references to AI flags or
    inconsistencies.
\end{itemize}

% ── Textual Use Case Descriptions ────────────────────────────
\subsubsection*{Use Case Descriptions — Admin Decision Workflow}

\begin{table}[H]
\centering
\renewcommand{\arraystretch}{1.4}
\begin{tabularx}{\textwidth}{|l|X|}
\hline
\rowcolor{medgray}
\textbf{Field} & \textbf{Description} \\ \hline

Use Case ID & UC-09 \\ \hline
Use Case Name & Approve Application \\ \hline
Actor(s) & Admin (primary) \\ \hline
Preconditions & Application in Pending Review \\ \hline

Main Flow &
1. Select approve \newline
2. Add optional note \newline
3. Confirm decision \newline
4. Log decision \newline
5. Change state to Approved \newline
6. Activate account \newline
7. Notify user \\ \hline

Postconditions &
Account becomes Active \\ \hline

\end{tabularx}
\caption{Use Case UC-09 — Approve Application}
\end{table}

\begin{table}[H]
\centering
\renewcommand{\arraystretch}{1.4}
\begin{tabularx}{\textwidth}{|l|X|}
\hline
\rowcolor{medgray}
\textbf{Field} & \textbf{Description} \\ \hline

Use Case ID & UC-10 \\ \hline
Use Case Name & Reject Application \\ \hline
Actor(s) & Admin (primary), System (secondary) \\ \hline

Preconditions & Record is in \textit{Pending Admin Review} \\ \hline

Main Flow &
1. Admin selects \textit{Reject} \newline
2. System requires a rejection reason (mandatory) \newline
3. Admin may reference AI flags or inconsistencies \newline
4. Admin confirms decision \newline
5. System logs admin ID, timestamp, decision type, and reason \newline
6. System transitions record to \textit{Rejected} \newline
7. Psychologist is notified with rejection reason \\ \hline

Postconditions &
Record is \textit{Rejected}; psychologist can edit and resubmit (UC-06) \\ \hline

\end{tabularx}
\caption{Use Case UC-10 — Reject Application}
\end{table}

% ─────────────────────────────────────────────────────────────
\subsection{Audit Log}
\label{sec:uc-audit}

An audit log is maintained for all state transitions and security-sensitive
events (document upload, submission, AI result received, admin decision,
resubmission). This log supports operational debugging, dispute resolution, and
compliance requirements.

% ── Textual Use Case Descriptions ────────────────────────────
\subsubsection*{Use Case Descriptions — Audit Log}

\begin{table}[H]
\centering
\renewcommand{\arraystretch}{1.4}
\begin{tabularx}{\textwidth}{|l|X|}
\hline
\rowcolor{medgray}
\textbf{Field} & \textbf{Description} \\ \hline

Use Case ID & UC-11 \\ \hline
Use Case Name & Record Audit Event \\ \hline
Actor(s) & System (primary, automated) \\ \hline

Preconditions &
A state transition or security-sensitive action is triggered (e.g., document upload, submission, AI result, admin decision, resubmission) \\ \hline

Main Flow &
1. System detects a loggable event \newline
2. System captures event details (type, actor, timestamp, record reference, metadata) \newline
3. System appends the event to the immutable audit log \\ \hline

Postconditions &
Event is permanently recorded and available for debugging, dispute resolution, and compliance \\ \hline

\end{tabularx}
\caption{Use Case UC-11 — Record Audit Event}
\end{table}

\begin{table}[H]
\centering
\renewcommand{\arraystretch}{1.4}
\begin{tabularx}{\textwidth}{|l|X|}
\hline
\rowcolor{medgray}
\textbf{Field} & \textbf{Description} \\ \hline

Use Case ID & UC-12 \\ \hline
Use Case Name & Query Audit Log \\ \hline
Actor(s) & Admin (primary) \\ \hline

Preconditions & Admin is authenticated with audit access privileges \\ \hline

Main Flow &
1. Admin opens audit log view or applies filters (date range, event type, actor) \newline
2. System retrieves matching log entries \newline
3. System displays results in chronological order \newline
4. Admin inspects event details \\ \hline

Postconditions &
Admin retrieves required audit evidence \\ \hline

\end{tabularx}
\caption{Use Case UC-12 — Query Audit Log}
\end{table}

% ── Consolidated Use Case Diagram ───────────────────────────
\subsubsection*{ Use Case Diagram — Psychologist Onboarding}

\begin{figure}[H]
  \centering
  \includegraphics[width=0.95\textwidth,height=0.72\textheight,keepaspectratio]{assets/sprint2/usecase2.png}
  \caption{use case diagram for Sprint~2 psychologist onboarding.}
  \label{fig:sprint2-consolidated-use-case}
\end{figure}

% ============================================================
\section{Onboarding State Machine}
\label{sec:sprint2-state-machine}

Sprint~2 enforces onboarding progression through a controlled state machine.
The exact naming may vary in implementation, but the conceptual states are
stable:

\begin{table}[H]
\centering
\begin{tabularx}{\textwidth}{|l|X|X|}
\hline
\rowcolor{medgray}
\textbf{State} & \textbf{Meaning} & \textbf{Typical Transitions} \\
\midrule
\hline

Draft & Profile and/or documents are being prepared. &
Draft $\rightarrow$ Submitted (after completeness validation). \\ \hline

Submitted & Submission is complete and awaiting verification. &
Submitted $\rightarrow$ Verifying; Submitted $\rightarrow$ Rejected (admin). \\ \hline

Verifying & AI verification job is running or pending results. &
Verifying $\rightarrow$ Pending Admin Review; Verifying $\rightarrow$ Verification Failed (retry). \\ \hline

Pending Admin Review & Admin review is required before activation. &
Pending Admin Review $\rightarrow$ Approved; Pending Admin Review $\rightarrow$ Rejected. \\ \hline

Approved & Onboarding accepted; psychologist can be activated. &
Approved $\rightarrow$ Active (system action). \\ \hline

Rejected & Onboarding rejected; edits and resubmission allowed. &
Rejected $\rightarrow$ Draft (edit); Rejected $\rightarrow$ Submitted (resubmit). \\ \hline

\end{tabularx}
\caption{Onboarding state machine (conceptual model)}
\end{table}

The state machine prevents invalid operations such as approving without a
submission, submitting without required documents, or exposing privileged
features before activation.
\section{System Design}
\label{sec:sprint2-design}
\subsection{Class Diagram}
\begin{figure}[H]
  \centering
  \includegraphics[width=0.85\textwidth]{assets/sprint2/ClassDiagram.png}
  \label{fig:DiagramClass}
  \caption{Class Diagram: Psychologist Onboarding}
\end{figure}

\subsection{Sequence Diagram}
\begin{figure}[H]
\centering
  \includegraphics[width=0.85\textwidth]{assets/sprint2/PB-03.png}
  \label{fig:PB-03}
  \caption{Sequence Diagram: PB-03}
\end{figure}
\begin{figure}[H]
\centering
  \includegraphics[width=0.85\textwidth]{assets/sprint2/PB-04.png}
  \label{fig:PB-04}
  \caption{Sequence Diagram: PB-04}
\end{figure}
\begin{figure}[H]
\centering
  \includegraphics[width=0.85\textwidth]{assets/sprint2/PB-05.png}
  \label{fig:PB-05}
  \caption{Sequence Diagram: PB-05}
\end{figure}
\clearpage
\section{Testing and Validation Strategy}
\label{sec:sprint2-testing}

\begin{longtable}{@{} L{0.22\textwidth} L{0.68\textwidth} @{}}
\toprule
\textbf{Testing Area} &
\textbf{Validation Scope} \\
\midrule
\endfirsthead

\toprule
\textbf{Testing Area} &
\textbf{Validation Scope} \\
\midrule
\endhead

Unit testing & Validate profile field rules and error messages; test state transition guards for valid and invalid moves; verify RBAC permissions for psychologist and admin operations. \\
\midrule
Integration testing & Validate API communication between client and backend for onboarding actions; verify document storage interactions, including upload metadata and access controls; test AI pipeline integration through job creation, result persistence, and retries. \\
\midrule
End-to-end testing & Validate the complete onboarding workflow from draft to admin approval; test admin review flows with both approve and reject outcomes; verify access control so only approved psychologists become discoverable and gain privileged access. \\
\midrule
Test environment & Use a staging environment configured with production-like security controls and monitoring; rely on deterministic fixtures or mocked responses for external AI services to ensure reliable test execution. \\
\bottomrule
\caption{Testing and validation strategy for Sprint~2.}
\label{tab:sprint2-testing-strategy}\\
\end{longtable}
\begin{longtable}{@{} L{0.11\textwidth} L{0.48\textwidth} L{0.23\textwidth} C{0.12\textwidth} @{}}
\toprule
\textbf{Test ID} &
\textbf{Scenario} &
\textbf{Expected Result} &
\textbf{Result} \\
\midrule
\endfirsthead

\toprule
\textbf{Test ID} &
\textbf{Scenario} &
\textbf{Expected Result} &
\textbf{Result} \\
\midrule
\endhead

T2-1 & Psychologist submits onboarding profile with valid required fields & Profile is saved and moves to the expected review-ready state. & \textcolor{successgreen}{Passed} \\
\midrule
T2-2 & Psychologist submits incomplete or invalid profile data & Submission is rejected and field-level validation errors are returned. & \textcolor{successgreen}{Passed} \\
\midrule
T2-3 & Psychologist uploads credential documents & Document metadata is stored and access remains restricted to authorized users. & \textcolor{successgreen}{Passed} \\
\midrule
T2-4 & AI-assisted credential analysis job is created & Analysis task is queued, processed, and persisted with the related verification request. & \textcolor{successgreen}{Passed} \\
\midrule
T2-5 & Admin approves a valid psychologist profile & Psychologist account becomes approved and can access privileged psychologist features. & \textcolor{successgreen}{Passed} \\
\midrule
T2-6 & Admin rejects a psychologist profile with justification & Profile status becomes rejected and the rejection reason is stored for user feedback. & \textcolor{successgreen}{Passed} \\
\midrule
T2-7 & Unapproved psychologist attempts privileged access & Request is denied and the psychologist remains hidden from discoverable listings. & \textcolor{successgreen}{Passed} \\
\bottomrule
\caption{Representative validation scenarios for Sprint~2.}
\label{tab:sprint2-tests}\\
\end{longtable}

\section{Sprint Review and Retrospective}
\label{sec:sprint2-review}

\renewcommand{\arraystretch}{1.4}
\begin{table}[H]
\centering
\begin{tabularx}{\textwidth}{|l|X|}
\hline
\rowcolor{medgray}
\textbf{Category} & \textbf{Details} \\ \hline

Sprint Review &
\begin{itemize}
    \item Demonstration of onboarding UI (drafting, uploading, submission)
    \item Demonstration of AI verification results in admin interface
    \item Admin decision workflow with audit trail visibility
\end{itemize} \\ \hline

Challenges &
\begin{itemize}
    \item AI integration complexity (latency, failures, confidence handling)
    \item Security constraints for document storage and access
    \item Balancing user experience with strict validation rules
\end{itemize} \\ \hline

Improvements &
\begin{itemize}
    \item Improve user feedback (clear checklist, missing-item guidance)
    \item Optimise AI processing (async jobs, retries, better status UX)
    \item Extend admin tools (bulk actions, prioritisation, analytics)
\end{itemize} \\ \hline

\end{tabularx}
\caption{Sprint 2 Review and Retrospective Summary}
\end{table}
% ============================================================
\section{UI Demonstration}
\label{sec:sprint2-ui-demo}
% ─────────────────────────────────────────────────────────────

% ─────────────────────────────────────────────────────────────
% ─────────────────────────────────────────────────────────────
\begin{figure}[H]
    \centering
    \includegraphics[width=0.85\textwidth]{assets/interfaces/interface5.png}
    \caption{Credential upload interface with checklist tracking}
\end{figure}
\begin{figure}[H]
    \centering
    \includegraphics[width=0.85\textwidth]{assets/interfaces/interface6.png}
    \caption{Admin's verification review interface}
\end{figure}.
\begin{figure}[H]
    \centering
    \includegraphics[width=0.85\textwidth]{assets/interfaces/interface7.png}
    \caption{Admin's verification interface}
\end{figure}.
\begin{figure}[H]
    \centering
    \includegraphics[width=0.85\textwidth]{assets/interfaces/interface8.png}
    \caption{Admin's Audit log interface}
\end{figure}.

\subsection*{Discussion}

The presented interfaces confirm the alignment between system design and user
interaction. Emphasis was placed on usability, validation clarity, and security,
ensuring a reliable and intuitive onboarding experience.
\section{Conclusion}
\label{sec:sprint2-conclusion}

Sprint~2 delivers a robust onboarding pipeline that combines secure credential
handling, AI-assisted pre-screening, and an accountable administrative approval
workflow. By enforcing a controlled state machine and a complete audit trail,
the platform establishes the trust and governance needed to safely connect
patients with verified psychologists.
\clearpage
\chapter{Sprint 3: Search, Scheduling, and Booking Lifecycle}
\label{chap:sprint3-booking}

\section*{Introduction}
\phantomsection
\addcontentsline{toc}{section}{Introduction}

Sprint~3 introduces the first complete patient-facing operational workflow of the platform. After Sprint~1 established authentication and role-based access control, and Sprint~2 structured psychologist onboarding and verification, this sprint focuses on connecting patients with approved psychologists through search, availability consultation, and appointment booking. The implemented scope covers psychologist discovery, filtering, geospatial search, calendar availability, booking requests, status transitions, cancellation, and completion handling. Since these functionalities depend on time-sensitive and potentially concurrent data, particular attention is given to lifecycle correctness, access control, and the consistency of calendar-related information.

\section{Sprint Objectives and Scope}
\label{sec:sprint3-objectives}

The purpose of this sprint is to transform the platform from a verified user-management system into an operational consultation scheduling environment. The objective is not limited to displaying psychologists, but also includes the controlled management of availability and booking transitions.

\begin{figure}[H]
  \centering
  \includegraphics[width=0.84\textwidth]{assets/s3/objective3.jpg}
  \caption{Sprint~3 objectives and scope overview.}
  \label{fig:sprint3-objectives-scope}
\end{figure}

\subsection{Primary Objectives}

The primary objectives of Sprint~3 are defined according to the needs of both patients and psychologists, while preserving the technical constraints required by the booking lifecycle.

\begin{itemize}[leftmargin=*]
  \item[--] Allow patients to search for psychologists using structured filters.
  \item[--] Support geospatial search and distance-aware ordering when patient location is available.
  \item[--] Enable psychologists to define, update, and maintain availability calendars, blocks, and exceptions.
  \item[--] Implement the booking lifecycle, including request, confirmation, cancellation, decline, and completion states.
  \item[--] Support dynamic booking over availability blocks, including the coherent split of remaining time after confirmation when required by the scheduling policy.
  \item[--] Preserve data consistency across schedules, bookings, and role-based permissions.
\end{itemize}

\subsection{Scope Boundary}

Sprint~3 is deliberately limited to discovery and booking workflows. Session activation, real-time messaging, report generation, and AI-assisted interaction are excluded from this sprint and reserved for later development phases. This separation keeps the sprint technically coherent and makes it possible to validate the appointment lifecycle before adding communication and consultation features.

\section{Sprint Planning}
\label{sec:sprint3-planning}

Sprint planning for this phase focused on the stability of the search--availability--booking chain. Since the sprint introduces the first high-interaction workflow for patients, the planning process emphasized data coherence, state modeling, and explicit validation rules.

\subsection{Sprint Context}

This sprint introduces a stateful domain workflow that differs from previous authentication and onboarding operations. Booking features require consistent behavior under repeated user interaction, especially when several users may attempt to consult or reserve overlapping availability. Therefore, the planning phase prioritized the modeling of domain entities, the clarification of lifecycle transitions, and the definition of authorization boundaries.

\subsection{Stakeholders and Roles}

The sprint involves three main roles. Each role contributes to the appointment workflow from a different perspective.

\begin{itemize}[leftmargin=*]
  \item[--] Patient: searches for psychologists, consults profiles and availability, and submits booking requests.
  \item[--] Psychologist: publishes availability, manages calendar slots, and responds to booking requests.
  \item[--] System: enforces booking invariants, prevents inconsistent time-slot allocation, and records status changes.
\end{itemize}

\subsection{Execution Priorities}

The execution priorities were organized to reduce implementation risk. Search accuracy alone is insufficient if the availability and booking lifecycle is unstable. For this reason, the sprint first validates searchable psychologist data, then establishes availability management, formalizes booking state transitions, and finally integrates the patient-facing search and request flows.

\section{Sprint Backlog}
\label{sec:sprint3-backlog}

The Sprint~3 backlog, presented in Table~\ref{tab:sprint3-backlog}, summarizes the selected product backlog items related to search, location-aware discovery, calendar management, and booking lifecycle control.

\begin{longtable}{@{} L{0.11\textwidth} L{0.20\textwidth} L{0.43\textwidth} C{0.11\textwidth} C{0.07\textwidth} @{}}
\caption{Sprint~3 backlog dedicated to patient engagement and appointment management.}
\label{tab:sprint3-backlog}\\
\toprule
\textbf{ID} &
\textbf{Title} &
\textbf{User Story Summary} &
\textbf{Priority} &
\textbf{SP} \\
\midrule
\endfirsthead

\toprule
\textbf{ID} &
\textbf{Title} &
\textbf{User Story Summary} &
\textbf{Priority} &
\textbf{SP} \\
\midrule
\endhead

PB-06 & Search \& Discovery & Patient searches psychologists by specialty, language, and profile criteria. & High & 5 \\
\midrule
PB-07 & Geospatial Search & Patient ranks or filters psychologists according to geographic proximity. & Medium & 5 \\
\midrule
PB-08 & Calendar Management & Psychologist creates and updates availability slots. & High & 5 \\
\midrule
PB-09 & Booking Lifecycle & Patient requests bookings and psychologists confirm, decline, cancel, or complete them. & High & 8 \\

\bottomrule
\end{longtable}

\section{Requirement Analysis}
\label{sec:sprint3-analysis}

The requirement analysis clarifies the functional chain implemented during Sprint~3 and the constraints that govern it. The objective is to ensure that the interface remains simple for users while the backend preserves strict scheduling and authorization rules.

\subsection{Functional Analysis}

Sprint~3 combines search, scheduling, and booking into one coherent service chain. Each function contributes to the same objective: enabling a patient to identify an appropriate psychologist and initiate a valid appointment request.

\begin{itemize}[leftmargin=*]
  \item[--] Search and filtering: return approved psychologists according to specialization, location, language, and selected profile criteria.
  \item[--] Geospatial logic: compute distance-aware relevance when the patient provides a location.
  \item[--] Availability management: allow psychologists to create, update, and remove time slots or larger availability blocks while preventing overlapping conflicts.
  \item[--] Booking workflow: support request submission, confirmation, decline, cancellation, and completion through explicit status values.
  \item[--] Split-slot handling: when a patient books a sub-window inside a larger availability interval, the remaining free time must be re-materialized as coherent future availability according to the selected scheduling policy.
\end{itemize}

\subsection{Non-Functional Requirements}

The non-functional requirements of this sprint are directly linked to the reliability of the appointment workflow. They address consistency, performance, security, and traceability.

\begin{itemize}[leftmargin=*]
  \item[--] Consistency: the same slot must not be booked concurrently by several patients.
  \item[--] Performance: search and pagination must remain responsive even when several profile filters are applied.
  \item[--] Security: only authenticated and authorized actors may access calendars and bookings related to them.
  \item[--] Traceability: booking state changes must remain visible and auditable.
\end{itemize}

\subsection{Business Constraints}

In addition to technical constraints, Sprint~3 is governed by business rules that preserve the correctness of the appointment lifecycle.

\begin{itemize}[leftmargin=*]
  \item[--] Only approved psychologists can appear in search results and accept bookings.
  \item[--] A booking request must always reference an existing and valid slot or bookable sub-window.
  \item[--] A cancelled or declined booking cannot transition directly to completed.
  \item[--] A confirmed booking must lock the requested interval and prevent double booking under concurrent access.
  \item[--] When confirmation splits an availability block, the resulting free segments must remain coherent and non-overlapping.
  \item[--] Availability publication and booking confirmation must remain scoped to the psychologist who owns the schedule.
\end{itemize}

\subsection{Transition to Design}

The requirement analysis shows that Sprint~3 is governed less by interface complexity than by lifecycle discipline. For this reason, the design treats booking as a stateful domain process rather than as a simple set of CRUD operations. This choice improves the clarity of transitions and reduces the risk of inconsistent appointment states.

\section{Architecture Overview}
\label{sec:sprint3-architecture}

The architecture of Sprint~3 extends the platform by introducing stronger interaction between profile data, availability data, and booking records. The architectural view therefore focuses on how these modules cooperate to provide reliable scheduling behavior.

\subsection{Module Interaction}

Sprint~3 strengthens the coupling between the profile module, the availability module, and the booking module. Search consumes psychologist profile data, availability exposes future bookable slots, and booking uses both profile eligibility and slot validity to authorize request creation and lifecycle transitions.

Because this sprint is dominated by stateful interactions rather than by complex structural layering, behavioral diagrams are particularly useful. The detailed execution flows are therefore presented later in Section~\ref{sec:sprint3-workflow}, after the use-case and class-level views have clarified the functional and structural foundations.

\subsection{Database Design}

The sprint requires persistent entities beyond the user account itself. Table~\ref{tab:sprint3-entities} summarizes the principal entities involved in search, availability publication, and booking management.

\begin{longtable}{@{} L{0.22\textwidth} L{0.24\textwidth} L{0.46\textwidth} @{}}
\caption{Principal persistent entities introduced or activated by Sprint~3.}
\label{tab:sprint3-entities}\\
\toprule
\textbf{Entity} &
\textbf{Key Attributes} &
\textbf{Role in the Sprint} \\
\midrule
\endfirsthead

\toprule
\textbf{Entity} &
\textbf{Key Attributes} &
\textbf{Role in the Sprint} \\
\midrule
\endhead

\makecell[l]{\texttt{Psychologist}\\\texttt{Profile}} &
specialty, languages, location, approval status &
Makes approved professionals searchable and filterable. \\
\midrule

\texttt{AvailabilitySlot} &
psychologistId, date, startTime, endTime, status &
Stores declared availability, blocked periods, or derived free intervals. \\
\midrule

\texttt{Booking} &
patientId, psychologistId, slotId, status, timestamps &
Represents the booking lifecycle from request to completion and anchors slot consumption. \\

\bottomrule
\end{longtable}

The domain entities of Sprint~3 are sufficiently interdependent to justify an explicit structural model. The class diagram in Figure~\ref{fig:sprint3-class} makes the ownership rules and relationships easier to interpret before the behavioral workflows are discussed.

\begin{figure}[H]
  \centering
  \includegraphics[width=0.92\textwidth]{assets/s3/d2.png}
  \caption{Class diagram representing the core scheduling and booking entities of Sprint~3, including the relationships between patients, psychologists, availability slots, and booking records.}
  \label{fig:sprint3-class}
\end{figure}

This structural view highlights that booking is a relation-rich domain object. It links patients, psychologists, and availability slots while preserving an explicit lifecycle state.

\section{Functional Specification}
\label{sec:sprint3-functional}

The functional specification translates the requirements of Sprint~3 into UML-oriented views and detailed use-case descriptions. It clarifies the interactions between patients, psychologists, and the system.

The use case diagram in Figure~\ref{fig:sprint3-usecase} summarizes the main interactions supported by this sprint. It positions the booking lifecycle as the central domain process, surrounded by profile search and calendar management.

\begin{figure}[H]
  \centering
  \includegraphics[width=0.92\textwidth]{assets/s3/d1.png}
  \caption{Overall use case diagram for Sprint~3, covering psychologist search, availability management, and the booking lifecycle.}
  \label{fig:sprint3-usecase}
\end{figure}

The use cases show two complementary perspectives. The patient navigates from search to booking, whereas the psychologist manages the service supply by maintaining availability and responding to appointment requests.

\subsection{Detailed Use Cases}

The following tables specify the principal use cases implemented during this sprint. They focus on search, availability management, and booking lifecycle processing.

\begin{longtable}{@{} L{0.28\textwidth} L{0.66\textwidth} @{}}
\caption{Use case specification for psychologist search.}
\label{tab:uc-s3-search}\\
\toprule
\multicolumn{2}{@{}l@{}}{\textbf{UC-06 \quad Search and Filter Psychologists}} \\
\midrule
\endfirsthead

\toprule
\multicolumn{2}{@{}l@{}}{\textbf{UC-06 \quad Search and Filter Psychologists}} \\
\midrule
\endhead

Primary Actor & Patient \\
\midrule
Goal & Identify approved psychologists matching selected criteria. \\
\midrule
Preconditions & Patient is authenticated and searchable profiles exist. \\
\midrule
Main Success Scenario & Patient submits search filters; the backend applies profile constraints and pagination; matching psychologists are returned. \\
\midrule
Postconditions & A filtered result set is displayed to the patient. \\

\bottomrule
\end{longtable}

\begin{longtable}{@{} L{0.28\textwidth} L{0.66\textwidth} @{}}
\caption{Use case specification for calendar management.}
\label{tab:uc-s3-availability}\\
\toprule
\multicolumn{2}{@{}l@{}}{\textbf{UC-08 \quad Manage Availability}} \\
\midrule
\endfirsthead

\toprule
\multicolumn{2}{@{}l@{}}{\textbf{UC-08 \quad Manage Availability}} \\
\midrule
\endhead

Primary Actor & Psychologist \\
\midrule
Goal & Publish and update available consultation slots. \\
\midrule
Preconditions & Psychologist account is authenticated and approved. \\
\midrule
Main Success Scenario & Psychologist creates or updates slots or larger availability blocks; the backend validates time coherence and overlap rules, then stores the resulting availability state. \\
\midrule
Alternate Flow & Conflicting, invalid, or policy-incompatible time ranges are rejected. \\

\bottomrule
\end{longtable}

\begin{longtable}{@{} L{0.28\textwidth} L{0.66\textwidth} @{}}
\caption{Use case specification for the booking lifecycle.}
\label{tab:uc-s3-booking}\\
\toprule
\multicolumn{2}{@{}l@{}}{\textbf{UC-09 \quad Request and Process Booking}} \\
\midrule
\endfirsthead

\toprule
\multicolumn{2}{@{}l@{}}{\textbf{UC-09 \quad Request and Process Booking}} \\
\midrule
\endhead

Primary Actors & Patient, Psychologist \\
\midrule
Goal & Create a booking request and move it through explicit lifecycle states. \\
\midrule
Preconditions & A valid availability slot exists and both actors are authorized. \\
\midrule
Main Success Scenario & Patient selects a slot or sub-window and submits a booking request; the system creates a booking in pending state; the psychologist confirms or declines the request; on confirmation, the reserved interval is locked and any remaining free interval is split according to policy; later transitions allow cancellation or completion according to defined rules. \\
\midrule
Postconditions & The booking status reflects the latest valid lifecycle transition, slot integrity is preserved, and the workflow remains auditable. \\

\bottomrule
\end{longtable}

\section{Workflow Description}
\label{sec:sprint3-workflow}

The workflow description complements the structural and functional views by detailing how user actions are executed over time. The central behavioral artifact of Sprint~3 is the booking transition chain.

The first sequence diagram, shown in Figure~\ref{fig:sprint3-sequence-search}, clarifies how discovery and availability cooperate before a booking request is submitted.

\begin{figure}[H]
  \centering
  \includegraphics[width=0.92\textwidth]{assets/s3/d3.png}
  \caption{Sequence diagram describing psychologist search, profile selection, and slot retrieval during Sprint~3.}
  \label{fig:sprint3-sequence-search}
\end{figure}

This sequence shows that the search stage is not isolated from the scheduling domain. It prepares the booking process by ensuring that the patient selects an approved psychologist and consults valid availability data.

\begin{figure}[H]
  \centering
  \includegraphics[width=0.92\textwidth]{assets/s3/d4.png}
  \caption{Sequence diagram representing the booking request and confirmation workflow, including the validation of availability and lifecycle transition rules.}
  \label{fig:sprint3-sequence}
\end{figure}

As illustrated in Figure~\ref{fig:sprint3-sequence}, the booking lifecycle depends on coordinated checks across profile eligibility, slot validity, split-slot consistency, and role permissions. This avoids treating booking as a simple form submission and reinforces the consistency of the scheduling process.

To further formalize lifecycle correctness, the state transitions are modeled explicitly in Figure~\ref{fig:sprint3-state}.

\begin{figure}[H]
  \centering
  \includegraphics[width=0.7\textwidth]{assets/s3/d5.png}
  \caption{State model governing the valid lifecycle transitions of a booking request in Sprint~3.}
  \label{fig:sprint3-state}
\end{figure}

The state model prevents ambiguous or contradictory outcomes, such as completing a booking that was never confirmed or confirming a booking whose slot has already been cancelled.

\section{API Layer}
\label{sec:sprint3-api}

The API layer exposes search, calendar, and booking operations through protected routes. These routes are governed by authentication, role-based permissions, and ownership checks.

Table~\ref{tab:sprint3-api} presents the representative endpoint families implemented or activated during Sprint~3. The objective is to keep each endpoint family aligned with one clear responsibility.

\begin{longtable}{@{} L{0.36\textwidth} C{0.14\textwidth} L{0.42\textwidth} @{}}
\caption{Representative API responsibilities for Sprint~3.}
\label{tab:sprint3-api}\\
\toprule
\textbf{Endpoint Family} &
\textbf{Method} &
\textbf{Responsibility} \\
\midrule
\endfirsthead

\toprule
\textbf{Endpoint Family} &
\textbf{Method} &
\textbf{Responsibility} \\
\midrule
\endhead

\texttt{/api/psychologists} &
GET &
Search and filter approved psychologists with pagination. \\
\midrule

\makecell[l]{\texttt{/api/psychologists/}\\\texttt{nearby}} &
GET &
Return geospatially ranked psychologists relative to a patient location. \\
\midrule

\makecell[l]{\texttt{/api/calendar}\\\texttt{/availability}\\\texttt{/:psychId}} &
GET &
Expose bookable availability for a selected psychologist. \\
\midrule

\makecell[l]{\texttt{/api}\\\texttt{/availability}} &
\makecell[c]{\scriptsize POST/PUT/\\\scriptsize DELETE} &
Create, update, or remove availability slots or larger availability blocks. \\
\midrule

\makecell[l]{\texttt{/api}\\\texttt{/bookings}} &
\makecell[c]{\scriptsize POST/GET/\\\scriptsize PATCH} &
Create booking requests and process lifecycle transitions, including confirmation-driven slot locking. \\

\bottomrule
\end{longtable}

\section{Interface Screenshot}
\label{sec:sprint3-ui}

\begin{figure}[H]
  \centering
  \includegraphics[width=0.84\textwidth]{assets/interfaces/interface9.png}
  \caption{Sprint~3 psychologist search and filtering interface.}
  \label{fig:sprint3-search-ui}
\end{figure}
\begin{figure}[H]
  \centering
  \includegraphics[width=0.84\textwidth]{assets/interfaces/interface12.png}
  \caption{Sprint~3 psychologist availability management interface.}
  \label{fig:sprint3-calendar-ui}
\end{figure}
\section{Security Considerations}
\label{sec:sprint3-security}

Sprint~3 extends security concerns beyond identity management. Search results must only expose approved professionals, availability data must not leak beyond authorized contexts, and booking records must remain visible only to the actors directly involved. These rules rely on the role-based access control introduced earlier and are strengthened through ownership checks on slot and booking resources. This security model is essential because appointment records may reveal sensitive information about patient behavior, consultation preferences, and psychologist availability.

\section{Validation Strategy}
\label{sec:sprint3-validation}

The validation strategy focuses on the correctness of the search--availability--booking chain. Table~\ref{tab:sprint3-tests} summarizes representative test scenarios that verify filtering, geospatial ranking, slot conflict detection, booking creation, concurrency control, and access protection.

\begin{longtable}{@{} L{0.12\textwidth} L{0.46\textwidth} L{0.24\textwidth} C{0.12\textwidth} @{}}
\caption{Representative test scenarios for Sprint~3.}
\label{tab:sprint3-tests}\\
\toprule
\textbf{Test ID} &
\textbf{Scenario} &
\textbf{Expected Result} &
\textbf{Result} \\
\midrule
\endfirsthead

\toprule
\textbf{Test ID} &
\textbf{Scenario} &
\textbf{Expected Result} &
\textbf{Result} \\
\midrule
\endhead

T3-1 & Search with specialization and language filters & Matching approved psychologists are returned with pagination. & \textcolor{successgreen}{Passed} \\
\midrule
T3-2 & Geospatial search with patient coordinates & Results are ranked or filtered according to distance logic. & \textcolor{successgreen}{Passed} \\
\midrule
T3-3 & Create overlapping availability slots or blocks & The conflicting availability definition is rejected. & \textcolor{successgreen}{Passed} \\
\midrule
T3-4 & Create booking on free slot & Booking is created in pending state. & \textcolor{successgreen}{Passed} \\
\midrule
T3-5 & Confirm booking on a sub-window inside a larger interval & Status becomes confirmed and remaining free intervals are split coherently. & \textcolor{successgreen}{Passed} \\
\midrule
T3-6 & Two patients request the same target slot concurrently & Only one booking succeeds; the other is rejected or conflicts safely. & \textcolor{successgreen}{Passed} \\
\midrule
T3-7 & Attempt unauthorized booking access & Request is denied by access-control rules. & \textcolor{successgreen}{Passed} \\

\bottomrule
\end{longtable}

\subsection{Deployment Notes}

Deployment for Sprint~3 focused on verifying the interoperability of search, calendar, and booking modules in a staging-like environment. Particular attention was given to seeded data quality, date handling consistency, and the ability to reproduce complete request-confirmation-cancellation workflows. At this stage, validation remains primarily functional and integration-oriented; broader stress testing, transaction-level concurrency testing, and real user evaluation should be considered in later iterations.

\begin{comment}
\section{Discussion and Limitations}
\label{sec:sprint3-discussion}

Although Sprint~3 delivers a coherent appointment workflow, the updated
documentation highlights several operational limits that remain important. The
first concerns \textbf{time handling}: availability generation, booking
requests, and confirmation flows depend on rigorous time-zone normalization,
and these rules become harder to maintain as the platform serves geographically
distributed users. The second concerns \textbf{concurrency}: when multiple
patients target the same slot, conflict prevention must rely on transactional
guards, locks, or uniqueness constraints to guarantee that only one booking is
accepted. The third concerns \textbf{privacy and geospatial precision}: the
location subsystem should avoid exposing exact coordinates when approximate
ranking is sufficient for discovery. Finally, real-time status updates and map
interactions remain useful but non-essential enrichments; the sprint's actual
core contribution lies in the correctness of the search--availability--booking
state chain.
\end{comment}

\section{Sprint Review and Retrospective}
\label{sec:sprint3-review}

Sprint~3 achieved a coherent patient engagement workflow by linking search, availability consultation, and booking request processing. The sprint also confirmed the importance of explicit state modeling for time-sensitive features. The main retrospective insight is that appointment-related functionalities should be governed by clear invariants from the beginning, rather than being implemented as isolated interface actions and corrected afterward. Future refinements should focus on stronger concurrency guarantees, more systematic time-zone handling, and the replacement of interface placeholders with validated screenshots from the final implementation.

\section{Conclusion}
\label{sec:sprint3-conclusion}

Sprint~3 transforms the platform into an operational scheduling system by linking psychologist discovery, availability management, and appointment requests within a single consistent lifecycle. Its main contribution is the formalization of booking behavior as a controlled state process, supported by structured APIs, role-based access checks, and representative validation scenarios. This sprint therefore prepares the platform for subsequent session execution, communication, and consultation-support features.

\chapter{Sprint 4: Session Management, Messaging, and Clinical Summarization}
\label{chap:sprint4-MMC}

% ------------------------------------------------------------------
\section{Introduction}
% ------------------------------------------------------------------
Sprint~4 consolidates the operational consultation loop by integrating session activation, real-time communication, voice messaging, report generation, chatbot interaction, and AI-assisted clinical summarization. Because these features span several subsystems, the sprint is treated as an integration-focused iteration.

% ------------------------------------------------------------------
\section{Sprint Objectives and Scope}
% ------------------------------------------------------------------
\label{sec:sprint4-objectives}
The objective of Sprint~4 is to support the controlled execution of the consultation process after a booking has been confirmed. The scope includes payment-based session activation, secure access-code handling, real-time communication, voice exchange, session closure, report production, and AI-assisted preparation before the psychologist meets the patient.

\begin{figure}[H]
  \centering
  \includegraphics[width=0.84\textwidth]{assets/sprint4/objective4.jpg}
  \caption{Sprint~4 objectives and scope overview.}
  \label{fig:sprint4-objectives-scope}
\end{figure}

\begin{itemize}[leftmargin=*]
  \item[--] Enable payment confirmation and email-based access code issuance, including code persistence, expiration control, and validation.
  \item[--] Activate consultation sessions using the issued access code while enforcing authorization and session state transitions.
  \item[--] Provide real-time text messaging within a shared consultation interface using WebSocket-based communication.
  \item[--] Support voice message capture, upload, storage, and playback as a complement to text interaction.
  \item[--] Allow psychologists to terminate a session and generate a PDF report suitable for clinical record-keeping.
  \item[--] Integrate patient chatbot interaction and generate a structured clinical summary from chatbot history before the session.
\end{itemize}

\textbf{Out of scope.} Sprint~4 does not introduce new discovery mechanisms, new booking rules, public ratings, document-query governance, multilingual interface management, or administrative analytics. It assumes that the appointment has already been created in Sprint~3 and focuses on what happens between payment confirmation, session participation, and clinical closure.

\textbf{Success Criteria.}
\begin{itemize}[leftmargin=*]
  \item[--] A confirmed booking can be activated only through a valid payment-linked access code.
  \item[--] Session participants can exchange real-time text messages and voice messages within the authorized consultation context.
  \item[--] The psychologist can close the session and produce a traceable PDF report.
  \item[--] The chatbot history can be transformed into a structured clinical summary without exposing unauthorized patient data.
\end{itemize}

\subsection*{Out of scope (explicitly deferred)}

The following elements are intentionally deferred in order to keep Sprint~4 focused on the core integration loop.

\begin{itemize}[leftmargin=*]
  \item[--] Complex payment provider integrations beyond the minimal confirmation flow described in this sprint, such as full card processing, refunds, or chargebacks.
  \item[--] Human moderation workflows for clinical summaries, since summaries are system-generated and can be revised in later sprints.
  \item[--] Long-term media CDN integration, as audio storage is handled through the current server-managed upload strategy.
\end{itemize}

% ------------------------------------------------------------------
\section{Sprint Planning}
% ------------------------------------------------------------------
Sprint~4 is planned as a multi-track iteration with coordinated integration checkpoints. This structure is necessary because session activation, communication, reporting, and AI processing must all operate consistently within the same lifecycle.

\begin{itemize}[leftmargin=*]
  \item[--] \emph{Track A --- Session/Payment:} implement code generation, email delivery, and activation state-machine updates.
  \item[--] \emph{Track B --- Communication:} implement and validate WebSocket event contracts, message persistence, and interface binding.
  \item[--] \emph{Track C --- AI Pipeline:} implement chatbot interaction endpoints, safe history persistence, and clinical summarization.
  \item[--] \emph{Track D --- Reporting and Compliance:} implement session termination semantics, PDF generation, and artifact access control.
\end{itemize}

Given the cross-cutting nature of these changes, the sprint plan emphasizes shared API contracts between frontend and backend, explicit state transitions, deterministic persistence behavior in MongoDB, and strict authorization checks at all sensitive endpoints.

% ------------------------------------------------------------------
\section{Assumptions}
% ------------------------------------------------------------------
Sprint~4 relies on a set of technical assumptions inherited from earlier sprints, including the availability of baseline access control, session-related platform services, and external dependencies required for messaging, file handling, and AI processing.

\begin{table}[H]
\centering
\begin{tabular}{@{} L{0.28\textwidth} L{0.64\textwidth} @{}}
\toprule
\textbf{Assumption} &
\textbf{Description} \\
\midrule
MERN stack baseline &
MongoDB Atlas is configured, Express APIs are accessible at \texttt{http://localhost:5000}, and the React client runs at \texttt{http://localhost:3000}. \\
Role-based access control &
JWT authentication and role restrictions for \texttt{patient}, \texttt{psychologist}, and \texttt{admin} are already implemented and consistently enforced through middleware. \\
Email delivery &
The server is configured with an email provider, such as Nodemailer with SMTP credentials stored in \texttt{.env}, to send transactional access-code messages. \\
WebSocket infrastructure &
Socket.io is configured on the backend, and the frontend can connect to it. Session or conversation rooms can be used to scope real-time events. \\
PDF generation support &
Server-side PDF generation dependencies, such as PDFKit, are available, and PDF artifacts can be stored and retrieved securely. \\
AI provider connectivity &
The platform has credentials and network access to an AI provider, such as Groq or Google Generative AI, with support for transient failure handling through timeouts, retries, or fallback messages. \\
\bottomrule
\end{tabular}
\caption{Technical assumptions required for Sprint~4 implementation.}
\label{tab:sprint4-assumptions}
\end{table}

% ------------------------------------------------------------------
\section{Deliverables}
% ------------------------------------------------------------------
The deliverables of Sprint~4 cover the operational features required to activate a session, conduct interaction between participants, close the session, and prepare clinical context through AI-supported processing.

\begin{itemize}[leftmargin=*]
  \item[--] Patient payment confirmation workflow and access code delivery by email (PB-10).
  \item[--] Session activation endpoints and client interface flows (PB-10).
  \item[--] Psychologist session termination and PDF report generation (PB-11).
  \item[--] Real-time text messaging and message persistence (PB-12).
  \item[--] Voice message recording, uploading, and playback within the session interface (PB-13).
  \item[--] Patient AI chatbot interaction interface and backend endpoints (PB-14).
  \item[--] System-generated structured clinical summary from chatbot history (PB-15).
  \item[--] Documentation updates and validation strategy for integrated features.
\end{itemize}

% ------------------------------------------------------------------
\section{Stakeholders and Roles}
% ------------------------------------------------------------------
The stakeholders in Sprint~4 engage with the session lifecycle from distinct operational perspectives, including patient access, psychologist interaction, and system-managed communication or AI support.

\begin{table}[H]
\centering
\begin{tabular}{@{} L{0.24\linewidth} L{0.33\linewidth} L{0.33\linewidth} @{}}
\toprule
\textbf{Stakeholder / Role} &
\textbf{Responsibility} &
\textbf{Interests} \\
\midrule
Patient &
Confirms payment, activates the session, exchanges messages, and uses the AI chatbot. &
Rapid activation, privacy, and usability. \\
Psychologist &
Conducts the session, sends and receives messages, ends the session, and generates the report. &
Reliable communication, clinical context, and record integrity. \\
System (Backend Services) &
Enforces business rules, state transitions, persistence, and notifications. &
Security, correctness, and auditability. \\
AI Service &
Generates chatbot responses and clinical summaries. &
Safe prompting, bounded inputs, and consistent outputs. \\
Admin (Operational) &
Monitors operational issues, handles disputes, and supports compliance. &
Traceability, access control, and observability. \\
\bottomrule
\end{tabular}
\caption{Stakeholders and roles involved in Sprint~4.}
\label{tab:sprint4-stakeholders}
\end{table}

% ------------------------------------------------------------------
\section{Sprint Backlog}
% ------------------------------------------------------------------
The Sprint~4 backlog consolidates the implementation scope associated with session activation, communication, reporting, and AI-assisted preparation. Each item corresponds to a specific part of the consultation lifecycle.

\begin{table}[H]
\centering
\small
\begin{tabular}{@{} L{0.07\textwidth} L{0.14\textwidth} L{0.30\textwidth} C{0.11\textwidth} C{0.06\textwidth} L{0.16\textwidth} @{}}
\toprule
\textbf{ID} &
\textbf{Title} &
\textbf{User Story Summary} &
\textbf{Priority} &
\textbf{SP} &
\textbf{Primary Components} \\
\midrule
PB-10 &
Session Payment Confirmation \& Access Code &
Patient confirms payment and receives an email access code to activate the session. &
High &
8 &
React UI, Express routes, Email service, Session model. \\
PB-11 &
Session Termination \& PDF Report &
Psychologist ends a session and generates a PDF report for clinical records. &
Medium &
5 &
Express routes, PDF generation service, Storage, Authorization. \\
PB-12 &
Real-Time Text Messaging &
Session participants exchange real-time text messages in the shared interface. &
High &
8 &
Socket.io, Message model, React chat UI. \\
PB-13 &
Voice Messaging &
Participants send voice messages using a microphone interface. &
Medium &
5 &
Media capture UI, Upload endpoint, Audio storage, Playback. \\
PB-14 &
AI Chatbot Interaction &
Patient interacts with the AI chatbot to express mental state before the session. &
High &
8 &
React chatbot UI, Chatbot routes, AI provider integration. \\
PB-15 &
Clinical Summary Generation &
System generates a structured clinical summary from chatbot history. &
High &
8 &
Summarization service, Prompt templates, Storage, Retrieval. \\
\bottomrule
\end{tabular}
\caption{Sprint~4 backlog covering session activation, reporting, messaging, and AI-assisted preparation.}
\label{tab:sprint4-backlog}
\end{table}
% ==================================================================
\section{Functional Specification}
% ==================================================================
The functional specification for Sprint~4 formalizes the main session-related use cases, including activation, reporting, text and voice communication, chatbot interaction, and clinical summary generation.

% ------------------------------------------------------------------
\subsection{Session Activation and Payment}
% ------------------------------------------------------------------
The session activation workflow connects payment confirmation to access-code generation and controlled session activation. This sequence ensures that access to the consultation channel is granted only after the required payment and verification steps have completed.
\subsubsection*{Use Case Descriptions --- Session Activation}

\begin{table}[H]
\centering
\begin{tabularx}{\textwidth}{@{} L{0.28\textwidth} X @{}}
\toprule
\multicolumn{2}{@{}l@{}}{\textbf{UC-10 \quad Confirm Payment \& Receive Access Code}} \\
\midrule
Primary Actor & Patient \\
\midrule
Supporting Actors & System (Payment/Session Service); Email Service \\
\midrule
Goal & Confirm payment, generate an access code, and deliver it to the patient by email. \\
\midrule
Preconditions & Patient is authenticated; session request exists; session is in \texttt{pending\_payment} state. \\
\midrule
Trigger & Patient clicks \textit{``Confirm Payment''} in the interface. \\
\midrule
Main Success Scenario &
\begin{minipage}[t]{\linewidth}\begin{itemize}[leftmargin=*, nosep]
  \item[--] Patient submits confirmation.
  \item[--] System validates session eligibility.
  \item[--] System generates access code.
  \item[--] System persists code with TTL or expiry.
  \item[--] System sends code by email.
  \item[--] Interface displays the \textit{``Code sent''} state.
\end{itemize}\end{minipage} \\
\midrule
Extensions / Alternate Flows &
A1 --- Email failure: System retains the code and retries; the interface displays a recoverable error.
A2 --- Session not eligible: Interface displays the reason and prevents confirmation. \\
\midrule
Postconditions & Access code exists and is associated with the session; patient is notified; session remains inactive until code is validated. \\
\midrule
Non-Functional Notes & Code must be cryptographically random; logs must not store the full code; confirmation must be rate-limited to prevent abuse. \\
\bottomrule
\end{tabularx}
\caption{UC-10 --- Confirm payment and receive access code.}
\label{tab:uc10}
\end{table}

\begin{table}[H]
\centering
\begin{tabularx}{\textwidth}{@{} L{0.28\textwidth} X @{}}
\toprule
\multicolumn{2}{@{}l@{}}{\textbf{UC-10b \quad Activate Session Using Access Code}} \\
\midrule
Primary Actor & Patient \\
\midrule
Supporting Actors & System (Session Service) \\
\midrule
Goal & Validate the access code and transition the session to active state. \\
\midrule
Preconditions & Patient is authenticated; patient possesses the code; session is in \texttt{pending\_activation} or \texttt{pending\_payment\_confirmed} state. \\
\midrule
Trigger & Patient submits the code in the activation interface. \\
\midrule
Main Success Scenario &
\begin{minipage}[t]{\linewidth}\begin{itemize}[leftmargin=*, nosep]
  \item[--] Patient submits code.
  \item[--] System validates code, expiry, and session association.
  \item[--] System transitions session to \texttt{active}.
  \item[--] System grants access to the session interface.
\end{itemize}\end{minipage} \\
\midrule
Extensions / Alternate Flows &
A1 --- Invalid or expired code: Interface shows an error and allows retry within rate limits.
A2 --- Session already active or completed: Interface shows the current status and redirects appropriately. \\
\midrule
Postconditions & Session becomes active; session participants can exchange messages. \\
\midrule
Non-Functional Notes & Brute-force attempts must be reduced through rate limiting and attempt counters; the code may be stored hashed if required by compliance policy. \\
\bottomrule
\end{tabularx}
\caption{UC-10b --- Activate session using access code.}
\label{tab:uc10b}
\end{table}

% ------------------------------------------------------------------
\subsection{Session Management and Reporting}
% ------------------------------------------------------------------
This part of the sprint covers controlled session termination and the generation of consultation reports. The workflow ensures that reporting occurs as a valid post-session operation rather than as an independent or premature action.

\subsubsection*{Use Case Descriptions --- Session Termination \& Report Generation}

\begin{table}[H]
\centering
\begin{tabularx}{\textwidth}{@{} L{0.28\textwidth} X @{}}
\toprule
\multicolumn{2}{@{}l@{}}{\textbf{UC-11 \quad End Session \& Generate PDF Report}} \\
\midrule
Primary Actor & Psychologist or Patient \\
\midrule
Supporting Actors & System (Session Service); PDF Generation Service \\
\midrule
Goal & Close a session and generate a clinical PDF report. \\
\midrule
Preconditions & The psychologist or patient is authenticated; the session is \texttt{active}; the authenticated user is assigned to the session. \\
\midrule
Trigger & Psychologist or patient clicks \textit{``End Session''} in the session interface. \\
\midrule
Main Success Scenario &
\begin{minipage}[t]{\linewidth}\begin{itemize}[leftmargin=*, nosep]
  \item[--] Psychologist or patient requests to end the session.
  \item[--] System validates that the requester belongs to the active session.
  \item[--] System transitions the session to \texttt{completed}.
  \item[--] System generates the PDF report with metadata and structured content.
  \item[--] System stores the report artifact securely.
  \item[--] Interface provides a link or download option for the report.
\end{itemize}\end{minipage} \\
\midrule
Extensions / Alternate Flows &
A1 --- PDF generation failure: Session still completes; report generation can be retried independently.
A2 --- Unauthorized attempt: System returns 403 and writes to the audit log. \\
\midrule
Postconditions & Session is completed; report artifact is available to authorized users, including the assigned psychologist, the session patient, and admin. \\
\midrule
Non-Functional Notes & PDF files must not leak sensitive data through public endpoints; storage must enforce access control and a predictable retention policy. \\
\bottomrule
\end{tabularx}
\caption{UC-11 --- End session and generate PDF report.}
\label{tab:uc11}
\end{table}

% ------------------------------------------------------------------
\subsection{Real-Time Messaging System}
% ------------------------------------------------------------------
The real-time messaging workflow enables both participants to exchange messages within an active session while preserving session scoping, persistence, and access control.

\subsubsection*{Use Case Descriptions --- Messaging}

\begin{table}[H]
\centering
\begin{tabularx}{\textwidth}{@{} L{0.28\textwidth} X @{}}
\toprule
\multicolumn{2}{@{}l@{}}{\textbf{UC-12 \quad Real-Time Text Chat}} \\
\midrule
Primary Actors & Patient; Psychologist \\
\midrule
Supporting Actors & System (Socket Service); Message Persistence Service \\
\midrule
Goal & Exchange real-time text messages within an active session. \\
\midrule
Preconditions & Participants are authenticated; session is active; both users are authorized for the session room. \\
\midrule
Trigger & User types a message and presses \textit{``Send''}. \\
\midrule
Main Success Scenario &
\begin{minipage}[t]{\linewidth}\begin{itemize}[leftmargin=*, nosep]
  \item[--] Client emits \texttt{send\_message}.
  \item[--] Server validates room membership.
  \item[--] Server persists message to MongoDB.
  \item[--] Server broadcasts \texttt{receive\_message} to room participants.
  \item[--] Clients render the message and update unread indicators.
\end{itemize}\end{minipage} \\
\midrule
Extensions / Alternate Flows &
A1 --- Temporary disconnect: Client queues the message and retries; server handles duplicates idempotently.
A2 --- Session not active: Server rejects the event and notifies the client. \\
\midrule
Postconditions & Message is persisted and visible to both participants in the session room. \\
\midrule
Non-Functional Notes & Low latency is required; message ordering must be preserved per room; payload size limits and sanitization are required to reduce XSS risk. \\
\bottomrule
\end{tabularx}
\caption{UC-12 --- Real-time text chat.}
\label{tab:uc12}
\end{table}

% ------------------------------------------------------------------
\subsection{Voice Messaging}
% ------------------------------------------------------------------
Voice messaging extends the session communication model by supporting short audio exchanges between participants. This capability is integrated into the same protected interaction context as text messaging.

\subsubsection*{Use Case Descriptions --- Voice Messaging}

\begin{table}[H]
\centering
\begin{tabularx}{\textwidth}{@{} L{0.28\textwidth} X @{}}
\toprule
\multicolumn{2}{@{}l@{}}{\textbf{UC-13 \quad Voice Messaging}} \\
\midrule
Primary Actors & Patient; Psychologist \\
\midrule
Supporting Actors & System (Upload Service); Message Service \\
\midrule
Goal & Record and exchange short voice messages during an active session. \\
\midrule
Preconditions & Browser supports microphone access; participant is authenticated and authorized; session is active. \\
\midrule
Trigger & User presses the microphone button and confirms upload. \\
\midrule
Main Success Scenario &
\begin{minipage}[t]{\linewidth}\begin{itemize}[leftmargin=*, nosep]
  \item[--] Client records audio.
  \item[--] Client uploads audio blob to the backend.
  \item[--] Backend validates MIME type and file size, then stores the audio file.
  \item[--] Backend creates a message record of type \texttt{voice}.
  \item[--] Backend broadcasts message event to room participants.
  \item[--] Recipients can play audio inline in the interface.
\end{itemize}\end{minipage} \\
\midrule
Extensions / Alternate Flows &
A1 --- Microphone denied: Interface provides guidance and a text fallback option.
A2 --- Upload too large: Interface requests a shorter recording and displays size constraints. \\
\midrule
Postconditions & Voice message is stored and playable by authorized session participants. \\
\midrule
Non-Functional Notes & Strict size limits must be enforced; private storage paths must not be exposed; streaming-safe audio formats should be preferred. \\
\bottomrule
\end{tabularx}
\caption{UC-13 --- Voice messaging.}
\label{tab:uc13}
\end{table}

% ------------------------------------------------------------------
\subsection{AI Chatbot Interaction}
% ------------------------------------------------------------------
The chatbot interaction workflow allows the patient to express concerns before the session while generating structured conversational data that can later support summary generation.

\subsubsection*{Use Case Descriptions --- AI Chatbot}

\begin{table}[H]
\centering
\begin{tabularx}{\textwidth}{@{} L{0.28\textwidth} X @{}}
\toprule
\multicolumn{2}{@{}l@{}}{\textbf{UC-14 \quad Patient Chatbot Interaction}} \\
\midrule
Primary Actor & Patient \\
\midrule
Supporting Actors & System (Chatbot Service); AI Service \\
\midrule
Goal & Allow patient to express mental state before a session through a guided chatbot. \\
\midrule
Preconditions & Patient is authenticated; chatbot module is available; user consent to AI features is acknowledged when required by policy. \\
\midrule
Trigger & Patient opens the chatbot interface and submits a message. \\
\midrule
Main Success Scenario &
\begin{minipage}[t]{\linewidth}\begin{itemize}[leftmargin=*, nosep]
  \item[--] Patient sends message.
  \item[--] System appends message to chatbot history.
  \item[--] System calls AI provider with a bounded context window.
  \item[--] System stores AI response.
  \item[--] Interface renders response and updates history display.
\end{itemize}\end{minipage} \\
\midrule
Extensions / Alternate Flows &
A1 --- AI provider error or timeout: Interface shows \textit{``temporarily unavailable''} and retains the patient message for retry.
A2 --- Content policy violation: Response is refused and the user is guided with a neutral message. \\
\midrule
Postconditions & Chatbot history is persisted as structured turns and becomes available for clinical summary generation. \\
\midrule
Non-Functional Notes & Prompt size must be bounded; sensitive content must not appear in logs; per-user session isolation must be enforced. \\
\bottomrule
\end{tabularx}
\caption{UC-14 --- Patient chatbot interaction.}
\label{tab:uc14}
\end{table}

% ------------------------------------------------------------------
\subsection{Clinical Summary Generation}
% ------------------------------------------------------------------
Clinical summary generation transforms pre-session chatbot history into structured information intended to support psychologist preparation. The output is positioned as consultation support material rather than a replacement for clinical judgment.

\subsubsection*{Use Case Descriptions --- Clinical Summary}

\begin{table}[H]
\centering
\begin{tabularx}{\textwidth}{@{} L{0.28\textwidth} X @{}}
\toprule
\multicolumn{2}{@{}l@{}}{\textbf{UC-15 \quad Clinical Summary Generation}} \\
\midrule
Primary Actor & System (automated pipeline) \\
\midrule
Supporting Actors & AI Service; Psychologist as consumer \\
\midrule
Goal & Generate a structured clinical summary from patient chatbot history. \\
\midrule
Preconditions & Patient chatbot history exists; system is authorized to process it; session context exists or is linkable. \\
\midrule
Trigger & Automatic job at session creation or activation, or explicit \textit{``Generate Summary''} action depending on implementation. \\
\midrule
Main Success Scenario &
\begin{minipage}[t]{\linewidth}\begin{itemize}[leftmargin=*, nosep]
  \item[--] System fetches recent chatbot turns.
  \item[--] System normalizes and filters content.
  \item[--] System calls AI provider with a structured summary prompt.
  \item[--] System validates output structure.
  \item[--] System stores summary in MongoDB.
  \item[--] Psychologist can view summary in patient or session interface.
\end{itemize}\end{minipage} \\
\midrule
Extensions / Alternate Flows &
A1 --- Token or payload overflow: System truncates history and retries with reduced context.
A2 --- Malformed AI output: System stores a safe fallback message and flags the record for regeneration. \\
\midrule
Postconditions & Structured clinical summary is available to authorized psychologists and linked to the session. \\
\midrule
Non-Functional Notes & Output must be clearly labeled as AI-generated; defensive output parsing is required; versioned prompt metadata should be stored for traceability. \\
\bottomrule
\end{tabularx}
\caption{UC-15 --- Clinical summary generation.}
\label{tab:uc15}
\end{table}

\subsubsection*{Overall Use Case Diagram --- Sprint~4}
\begin{figure}[H]
  \centering
  \includegraphics[width=0.92\textwidth]{assets/sprint4/sprin4usecase.png}
  \caption{Overall use case diagram for Sprint~4 session activation, communication, report generation, and AI-assisted preparation.}
  \label{fig:sprint4-overall-usecase}
\end{figure}

% ==================================================================
\section{System Design}
% ==================================================================
Sprint~4 maintains a modular MERN architecture with explicit separation between the client interface, backend routing, service logic, and MongoDB persistence. This layering supports maintainability while ensuring that sensitive operations remain under server-side control.

\begin{itemize}[leftmargin=*]
  \item[--] \emph{React Client:} UI state, microphone capture, WebSocket client, form validation, and view rendering.
  \item[--] \emph{Express Backend:} API endpoints, authorization middleware, file upload handling, and route grouping.
  \item[--] \emph{Services Layer:} payment code issuance, email delivery, WebSocket event validation, AI prompts and summarization, and PDF generation.
  \item[--] \emph{MongoDB Models:} sessions, messages, notifications, chatbot history, clinical summaries, and reports.
\end{itemize}

\subsection{Technical Architecture and Implementation}

\subsubsection{Real-time Messaging Architecture}
\begin{itemize}[leftmargin=*]
  \item[--] The Socket.IO gateway runs alongside the Express API or as a separate horizontally scalable gateway. Authentication occurs during the Socket.IO handshake: the client provides an access token (JWT), which the gateway validates and maps to a user id. After validation, the socket is joined to rooms such as \texttt{session:\{sessionId\}} and \texttt{user:\{userId\}}.
  \item[--] For multi-instance scaling, the Socket.IO Redis adapter is used to propagate events across nodes. Presence information is kept ephemeral in a Redis store while message persistence is written to MongoDB for auditability and history retrieval.
\end{itemize}

\subsubsection{Session Code Generation and Validation}
\begin{itemize}[leftmargin=*]
  \item[--] Session codes are generated server-side as high-entropy tokens. After payment confirmation, the server issues a single-use session code bound to \texttt{sessionId}, \texttt{userId}, and an issuance timestamp. Codes are hashed before being stored to allow safe verification without raw token persistence.
  \item[--] Validation accepts code submissions via a secure \texttt{POST /api/sessions/\{id\}/join} endpoint that checks: code hash match, expiry window, and that the requester is a participant on the booking. Successful validation transitions the session to \texttt{ACTIVE} for the participant and issues socket join permission.
\end{itemize}

\subsubsection{Payment Integration and Webhooks}
\begin{itemize}[leftmargin=*]
  \item[--] Payments use a provider (e.g., Stripe) with server-side creation of payment intents and verification via webhook handlers. Webhook handlers are idempotent and verify signatures before changing booking/payment states. On success, the payment handler triggers issuance of the session code and notifies the participants.
\end{itemize}

\subsubsection{Voice and Attachment Handling}
\begin{itemize}[leftmargin=*]
  \item[--] Clients upload voice artifacts using presigned URLs. The backend records metadata (duration, mime-type, checksum) and the messaging system references the artifact URL in persisted message documents. Transcoding and duration extraction are performed by background workers where needed.
\end{itemize}

\subsubsection{PDF Report Generation}
\begin{itemize}[leftmargin=*]
  \item[--] Report generation is performed by a worker that composes session metadata, selectable conversation excerpts, clinician notes, and templates into a PDF using a rendering library (e.g., Puppeteer or wkhtmltopdf). PDFs are stored in secure object storage and access is controlled via signed download links.
\end{itemize}

\subsubsection{AI Integration Points}
\begin{itemize}[leftmargin=*]
  \item[--] The AI chatbot is integrated through a constrained service that mediates requests to the LLM provider. The service enforces prompt templates, safety filters, and maintains a minimal request/response log for audit. During sessions, the chatbot can provide supportive prompts and a therapist briefing generator which summarizes salient emotional signals --- always flagged as AI-assistive and subject to clinician review.
\end{itemize}

\subsubsection{AI Chatbot: In-Session Integration --- Implementation Details}
The chatbot architecture is designed as a modular subsystem built around a low-latency turn API. It is supported by background workers that handle heavier tasks such as embedding generation, PDF/OCR ingestion, and summarization. The system also includes a vector index service and a dedicated policy and risk management service. The turn API authenticates users, loads the relevant \texttt{IntakeSession} context, performs retrieval and safety checks, and manages response generation through a protected LLM interface.

\paragraph{Runtime Pipeline (Per Turn)}
\begin{enumerate}[leftmargin=*, label=\arabic*.]
  \item Authenticate the user and load the session context.
  \item Normalize the language input, including partial Darija and Arabizi mapping, and preprocess the message while acknowledging that dialect coverage is not yet exhaustive.
  \item Retrieve relevant dialect lexicon entries and permission-filtered document chunks.
  \item Run deterministic phrase detectors together with an LLM-based risk classifier to detect possible risks and trigger escalation if necessary.
  \item Build a stage-aware prompt that includes persona constraints, retrieved context, and safety instructions.
  \item Generate the response through the guarded LLM service, then filter the output and attach provenance metadata.
  \item Store the generated \texttt{ChatbotMessage}, update the \texttt{IntakeSession}, queue summarization and compression tasks, and emit events for notifications or clinician review.
\end{enumerate}

\paragraph{Safety and Escalation}
The system follows a hybrid safety approach that combines high-precision regex and phrase detectors with an LLM-based contextual risk classifier. When high-risk situations are detected, the platform creates \texttt{risk\_alert} notifications, stores \texttt{AuditEvent} records, and provides clinicians with summaries supported by evidence snippets.

\paragraph{Emotional Analysis}
For each interaction, the system computes emotional indicators using lexical analysis, LLM-derived vectors, and prosodic features when audio is available. These signals are combined into an emotional timeline used in clinician briefings and post-session analytics.

\paragraph{Human-in-the-Loop}
Clinicians can review summaries and emotional analysis results to verify their accuracy. Their corrections are integrated into an active-learning pipeline that supports label storage, periodic re-indexing, and model calibration. Outputs with low confidence scores are automatically flagged for mandatory clinician review.

\paragraph{Observability and Auditing}
Every retrieval and generation request records lightweight provenance information, including model identifiers and versions, retrieval references, confidence scores, and sanitized input snapshots. This information supports reproducibility, debugging, and incident analysis while protecting sensitive user data.

\subsubsection{Operational Recommendations for Chatbot in Sessions}
\begin{itemize}[leftmargin=*]
  \item[--] Keep the turn API stateless and push heavy processing to workers; monitor queue depths and LLM call latencies.
  \item[--] Maintain a model registry that records which model version produced each embedding or generation result.
  \item[--] Provide clinician controls to disable AI assistance per-session and to request human-only mode when preferred.
\end{itemize}

\subsubsection{Security and Compliance}
\begin{itemize}[leftmargin=*]
  \item[--] All message and attachment accesses are permission-checked. Signed URLs include short TTLs and are scoped to the requester. Worker-to-storage access uses per-job short-lived credentials.
  \item[--] Webhook endpoints (payments, transcription callbacks) verify provider signatures and use idempotency keys to prevent duplicated processing.
\end{itemize}

\subsubsection{Observability and Operational Recommendations}
\begin{itemize}[leftmargin=*]
  \item[--] Instrument socket events, join/leave rates, message delivery latencies, and worker queue depths. Maintain dashboards and alerts for failed PDF generations, webhook failures, and high message error rates.
\end{itemize}

\subsubsection{Suggested Next Steps}
\begin{itemize}[leftmargin=*]
  \item[--] Add automated integration tests for the socket handshake and room authorization.
  \item[--] Implement an operational runbook for handling payment webhook retries and disputed payments.
\end{itemize}


\subsection{Class Diagram}

The Sprint~4 class diagram summarizes the main domain entities required for session
activation, real-time communication, voice messages, chatbot history, clinical summaries,
and PDF report generation. It highlights how the consultation session acts as the central
aggregate linking patients, psychologists, messages, access codes, chatbot turns, summaries,
and generated report artifacts.

\begin{figure}[H]
  \centering
  \includegraphics[width=0.95\textwidth]{assets/sprint4/classDiagram.png}
  \caption{Class diagram for Sprint~4 session activation, messaging, chatbot interaction, clinical summary generation, and report artifacts.}
  \label{fig:sprint4-class-diagram}
\end{figure}

The recommended Sprint~4 domain entities remain aligned with the operational needs of
session activation, messaging, chatbot history, and report generation.

\begin{itemize}[leftmargin=*]
  \item[--] \texttt{Session}: \texttt{patientId}, \texttt{psychologistId}, \texttt{status}, \texttt{accessCode}, \texttt{accessCodeExpiresAt}, timestamps.
  \item[--] \texttt{Message}: \texttt{sessionId}, \texttt{senderId}, \texttt{type} (\texttt{text|voice}), \texttt{content}, \texttt{assetPath}, timestamps.
  \item[--] \texttt{ChatbotTurn}: \texttt{userId}, \texttt{role} (\texttt{user|assistant}), \texttt{content}, timestamps.
  \item[--] \texttt{ClinicalSummary}: \texttt{userId}, \texttt{sessionId}, \texttt{summaryJson}, \texttt{model}, \texttt{promptVersion}, timestamps.
  \item[--] \texttt{ReportArtifact}: \texttt{sessionId}, \texttt{createdBy}, \texttt{pdfPath}, \texttt{hash}, timestamps.
\end{itemize}


\clearpage
\subsection*{Sequence Diagram}
\thispagestyle{plain}
\vspace*{\fill}
\begin{figure}[H]
  \centering
  \makebox[\textwidth][c]{%
    \includegraphics[width=\textwidth,height=0.62\textheight,keepaspectratio]{assets/sprint4/seq.png}%
  }
  \vspace{0.35cm}
  \caption{Sequence Diagram --- ``Activate Session and Start Messaging''}
\end{figure}
\vspace*{\fill}
\clearpage
% ==================================================================
\section{Testing and Validation Strategy}
% ==================================================================
The validation strategy for Sprint~4 is designed to cover both synchronous and asynchronous behavior, including API execution, event-driven messaging, media upload, report generation, and AI-assisted processing.

\begin{longtable}{@{} L{0.22\textwidth} L{0.68\textwidth} @{}}
\toprule
\textbf{Testing Area} &
\textbf{Validation Scope} \\
\midrule
\endfirsthead

\toprule
\textbf{Testing Area} &
\textbf{Validation Scope} \\
\midrule
\endhead

Unit testing & Validate access code generation, including entropy, format, and expiry handling; verify allowed session state transitions and rejection of invalid transitions; check AI prompt bounding and structural validation for summary outputs; validate deterministic PDF generation and report hash computation. \\
\midrule
Integration testing & Verify cooperation across API, MongoDB, messaging, email, storage, and AI modules; confirm persistence of sessions, codes, messages, chatbot turns, summaries, and reports; validate sandbox email queuing without exposing sensitive codes; test voice upload MIME/type validation and message event creation. \\
\midrule
End-to-end testing & Confirm the full consultation journey: payment confirmation, mocked email access code reception, session activation, real-time message delivery, psychologist session closure, PDF report generation and authorization, chatbot interaction, summary generation, and psychologist summary access. \\
\midrule
Test environment & Use a dedicated MongoDB test database, sandbox SMTP or provider credentials, deterministic AI mock responses for continuous integration, local Socket.io validation for room scoping, and local filesystem storage for media and PDF artifacts with cleanup. \\
\bottomrule
\caption{Testing and validation strategy for Sprint~4.}
\label{tab:sprint4-testing-strategy}\\
\end{longtable}

\begin{longtable}{@{} L{0.11\textwidth} L{0.48\textwidth} L{0.23\textwidth} C{0.12\textwidth} @{}}
\toprule
\textbf{Test ID} &
\textbf{Scenario} &
\textbf{Expected Result} &
\textbf{Result} \\
\midrule
\endfirsthead

\toprule
\textbf{Test ID} &
\textbf{Scenario} &
\textbf{Expected Result} &
\textbf{Result} \\
\midrule
\endhead

T4-1 & Patient confirms payment and receives an access code by email & A valid session code is generated, queued for email delivery, and linked to the correct booking. & \textcolor{successgreen}{Passed} \\
\midrule
T4-2 & Patient activates a session using a valid code & Session status changes to active and the user gains access to the consultation room. & \textcolor{successgreen}{Passed} \\
\midrule
T4-3 & Patient attempts activation with an expired or invalid code & Activation is rejected and no session access is granted. & \textcolor{successgreen}{Passed} \\
\midrule
T4-4 & Patient and psychologist exchange real-time messages & Messages are persisted, delivered through Socket.IO, and scoped to the correct session room. & \textcolor{successgreen}{Passed} \\
\midrule
T4-5 & Patient uploads a voice message during the session & File type and size are validated, media is stored, and a voice message event is created. & \textcolor{successgreen}{Passed} \\
\midrule
T4-6 & Patient completes chatbot interaction before consultation in Arabic, French, and English & AI summary is generated in each supported language flow and becomes visible to the authorized psychologist. & \textcolor{successgreen}{Passed} \\
\midrule
T4-7 & Patient completes chatbot interaction using Darija dialect & Dialectal input is partially interpreted through the Darija context index, but coverage still requires further improvement and validation. & \textcolor{warnorange}{In progress} \\
\midrule
T4-8 & Psychologist ends the consultation and generates a report & Session is completed, PDF report is generated, and access remains permission-controlled. & \textcolor{successgreen}{Passed} \\
\midrule
T4-9 & Unauthorized user attempts to access messages or reports & Request is denied and sensitive consultation data is not exposed. & \textcolor{successgreen}{Passed} \\
\bottomrule
\caption{Representative validation scenarios for Sprint~4.}
\label{tab:sprint4-tests}\\
\end{longtable}

% ==================================================================
\section{Sprint Review and Retrospective}
% ==================================================================
Sprint~4 is reviewed against functional completeness, workflow correctness, operational responsiveness, and the protection of sensitive actions. These criteria reflect the integrated and stateful nature of the sprint.

\begin{itemize}[leftmargin=*]
  \item[--] \emph{Functional completeness:} all PB-10 to PB-15 flows are demonstrable in an integrated environment.
  \item[--] \emph{Reliability and correctness:} no session activation without a valid code, no messaging without authorization, and report generation occurs only when the session is completed.
  \item[--] \emph{Performance:} real-time messaging latency remains acceptable for clinical interaction, and AI interactions are bounded so they do not block the interface indefinitely.
  \item[--] \emph{Security and privacy:} media and PDF endpoints enforce access control, logs do not expose sensitive content, and AI processing respects isolation boundaries.
\end{itemize}

Retrospective focus areas include the stability and documentation of WebSocket contracts, the effectiveness of AI prompt bounding against provider errors, and the maintainability of report generation and storage for future compliance needs.

% ==================================================================
\section{UI Demonstration}
\begin{figure}[H]
  \centering
  \includegraphics[width=0.84\textwidth]{assets/interfaces/interface10.png}
  \caption{Sprint~3 booking payement.}
  \label{fig:sprint3-booking-ui}
\end{figure}
\begin{figure}[H]
  \centering
  \includegraphics[width=0.84\textwidth]{assets/interfaces/interface13.png}
  \caption{Sprint~3 pdf AI Summary.}
  \label{fig:sprint3-pdfAiSummary}
\end{figure}
\begin{figure}[H]
  \centering
  \includegraphics[width=0.76\textwidth]{assets/interfaces/interface14.png}
  \caption{Sprint~3 Real Time Chatting.}
  \label{fig:sprint3-Chat-ui}
\end{figure}
\begin{figure}[H]
  \centering
  \includegraphics[width=0.76\textwidth]{assets/interfaces/interface15.png}
  \caption{Sprint~3 ChatBot.}
  \label{fig:sprint3-chatbot-ui}
\end{figure}
% ==================================================================
\section{Conclusion}
% ==================================================================
Sprint~4 establishes the platform's operational readiness for online consultation by integrating session activation, protected communication, report generation, and AI-assisted pre-session analysis. The sprint preserves the existing layering strategy while coordinating several subsystems within one controlled workflow.

\chapter{Sprint 5: Platform Enrichment, Governance, and Continuous Engagement}
\label{chap:sprint5-enrichment}

\section*{Introduction}
\phantomsection
\addcontentsline{toc}{section}{Introduction}

Sprint~5 extends \emph{PsychPlatform} beyond the core consultation lifecycle toward a more
mature and operationally complete platform. Whereas previous sprints focused on identity
management, psychologist onboarding, booking, sessions, and AI-assisted intake, this sprint
introduces enrichment features related to professional workflow, administrative governance,
multilingual accessibility, and long-term user engagement. The sprint integrates a psychologist
dashboard, document-based AI querying, ratings and feedback, an extended admin panel,
internationalization with right-to-left support, a floating assistant, and a notification
subsystem, while preserving the same principles of permission awareness, auditability, and
privacy protection.

% ─────────────────────────────────────────────────────────────────────────────
\section{Sprint Objectives and Scope}
\label{sec:sprint5-objectives}

Sprint~5 aims to improve the platform after the implementation of the main consultation
workflow. Its objective is not to redefine the core clinical process, but to strengthen
usability, supervision, accessibility, and engagement through additional controlled modules.

\begin{figure}[H]
  \centering
  \includegraphics[width=0.84\textwidth]{assets/s5/objective5.png}
  \caption{Placeholder --- Sprint~5 objectives and scope overview.}
  \label{fig:sprint5-objectives-scope}
\end{figure}

\subsection{Primary Objectives}

The primary objectives of this sprint are organized around platform enrichment, user
support, governance, and accessibility.

\begin{itemize}[leftmargin=*]
  \item[--] Provide psychologists with a dashboard covering analytics, sessions, reports, and
        workload indicators.
  \item[--] Enable document AI querying with strict permission boundaries and grounded,
        source-cited answers.
  \item[--] Implement a ratings and feedback system to improve trust and discovery relevance.
  \item[--] Extend the admin panel beyond onboarding to include user management, approvals,
        monitoring, and moderation support.
  \item[--] Deliver internationalization for EN/FR/AR, including RTL layout support for Arabic.
  \item[--] Add a floating assistant to provide contextual help and AI-driven guidance across
        the UI.
  \item[--] Provide a robust notifications system combining real-time and in-app delivery with
        queued retry logic.
\end{itemize}

\subsection{Success Criteria}

The sprint is considered successful when the following conditions are verifiably met.

\begin{itemize}[leftmargin=*]
  \item[--] Psychologists can monitor sessions and access reports from a single dashboard.
  \item[--] Users can query authorized documents safely with explainable, source-linked results
        and audit logs.
  \item[--] Ratings are tied to completed sessions, cannot be spammed, and are moderated
        by policy.
  \item[--] Admins can manage users and operational flows with full traceability.
  \item[--] The UI is fully usable in EN/FR/AR and supports RTL for Arabic.
  \item[--] Notifications are delivered reliably with user preferences and read/unread state.
\end{itemize}

\subsection{Scope Boundary}

Sprint~5 enriches the platform but does not redefine the core clinical flow. Its purpose is
to strengthen usability, supervision, and operational maturity on top of the capabilities
already delivered in earlier sprints.

% ─────────────────────────────────────────────────────────────────────────────
\section{Sprint Planning}
\label{sec:sprint5-planning}

Sprint planning for this iteration is guided by the cross-cutting nature of the new modules.
Unlike previous sprints that focused on a dominant business workflow, this sprint touches
several complementary concerns that must remain consistent with the existing platform
architecture.

\subsection{Sprint Context}

Sprint~5 is characterized by horizontal integration. Unlike earlier sprints, which each
introduced one dominant business flow, this sprint touches multiple cross-cutting concerns:
observability, moderation, AI governance, language support, and notification delivery.
Planning therefore prioritized modularity and permission boundaries.

\subsection{Dependencies}

Sprint~5 builds explicitly on earlier work. Sprint~4 delivered sessions, messaging, reports,
and AI chatbot integration; Sprint~3 delivered the booking lifecycle. Only active and
approved psychologists can interact with the dashboard and document query features.

\subsection{Stakeholders and Roles}

The stakeholders involved in Sprint~5 interact with different enrichment services. Their
roles are distributed across professional use, patient engagement, administrative supervision,
AI processing, and notification delivery.

\begin{itemize}[leftmargin=*]
  \item[--] \emph{Psychologist:} consults dashboard insights and interacts with document query tools.
  \item[--] \emph{Patient:} receives notifications, submits ratings, and uses a multilingual interface.
  \item[--] \emph{Admin:} supervises users, moderation actions, and operational events.
  \item[--] \emph{AI Service:} supports grounded question answering on uploaded documents.
  \item[--] \emph{Notification Service:} persists and distributes domain events to the appropriate
        recipients.
\end{itemize}

\subsection{Planned Deliverables}

The following deliverables were committed during sprint planning.

\begin{itemize}[leftmargin=*]
  \item[--] Psychologist dashboard views and analytics endpoints.
  \item[--] Document upload, AI indexing, and permission-scoped query endpoints.
  \item[--] Ratings and feedback UI, service endpoints, and aggregation logic.
  \item[--] Admin panel modules: user management, approvals, monitoring, and moderation tools.
  \item[--] i18n infrastructure and RTL-ready UI components (EN/FR/AR).
  \item[--] Floating assistant component integrated into key screens.
  \item[--] Notifications pipeline: event generation, delivery workers, in-app center, and
        real-time Socket.IO updates.
\end{itemize}

\subsection{Execution Strategy}

Given the breadth of Sprint~5, implementation was organized as a set of parallel yet bounded
streams: dashboard aggregation, AI document querying, ratings and moderation, localization,
assistant integration, and notification delivery. This separation reduced cross-module coupling
while preserving a common governance layer based on permission checks and audit logging.

The execution followed a value-first incremental approach: notifications and the dashboard
skeleton were delivered early for visible progress; data governance and permission scopes were
defined before document AI querying; i18n foundations and RTL styling were put in place before
translating all screens; and the floating assistant started as contextual FAQ/navigation
support before any AI depth was added.

% ─────────────────────────────────────────────────────────────────────────────
\section{Sprint Backlog}
\label{sec:sprint5-backlog}

\subsection{User Stories}

Table~\ref{tab:sprint5-userstories} presents the full set of user stories addressed in
Sprint~5, covering all seven enrichment modules.

\begin{longtable}{@{} L{0.10\textwidth} L{0.50\textwidth} L{0.32\textwidth} @{}}
\toprule
\textbf{ID} &
\textbf{User Story} &
\textbf{Key Acceptance Criteria} \\
\midrule
\endfirsthead

\toprule
\textbf{ID} &
\textbf{User Story} &
\textbf{Key Acceptance Criteria} \\
\midrule
\endhead

US-5.1 & As a psychologist, I want a dashboard that summarises my activity so that I can manage my work efficiently. &
Shows upcoming sessions, recent sessions, and pending tasks; displays analytics (sessions this week/month, cancellation rate); links to session details and report library. \\
\midrule
US-5.2 & As a psychologist, I want to view my session history and reports so that I can reference prior work securely. &
Session list supports date-range and status filters with pagination; reports accessible only to authorised participants. \\
\midrule
US-5.3 & As a user, I want to upload documents to a personal workspace so that I can search or ask questions about them. &
PDF format enforced with size limits; documents stored securely with owner and scope tags; upload triggers a chunking/embedding job. \\
\midrule
US-5.4 & As a user, I want to ask questions about my documents so that I can retrieve relevant information quickly. &
System returns answer plus referenced excerpts; access is permission-scoped; queries and results are logged for auditing. \\
\midrule
US-5.5 & As a patient, I want to leave a rating and feedback after my session so that I can share my experience. &
Rating only submittable for a \texttt{COMPLETED} session; one rating per session; average rating and count displayed in discovery under moderation rules. \\
\midrule
US-5.6 & As the system/admin, I want safeguards so that ratings cannot be abused. &
Basic spam prevention and prohibited content checks enforced; admin can hide/remove a review with an audited reason. \\
\midrule
US-5.7 & As an admin, I want to manage users so that I can ensure platform governance. &
Search users, view profiles, and change account status with justification; view psychologist onboarding status; all actions recorded in audit logs. \\
\midrule
US-5.8 & As an admin, I want a monitoring panel so that I can detect issues quickly. &
View failed payments, webhooks, AI jobs, PDFs, and notification failures; filterable audit log viewer for sensitive actions. \\
\midrule
US-5.9 & As a user, I want to change the app language so that I can use the platform comfortably. &
EN/FR/AR translations exist for key flows; language preference persisted per user and applied on next login. \\
\midrule
US-5.10 & As an Arabic-speaking user, I want the UI to render correctly in RTL so that it is readable and consistent. &
Layout flips appropriately; dates, times, and numbers formatted per locale. \\
\midrule
US-5.11 & As a user, I want a floating assistant so that I can get help while navigating the platform. &
Accessible from most screens; preserves conversation state; answers FAQ-style queries and routes users to relevant pages; sensitive data exposure prevented by access rules. \\
\midrule
US-5.12 & As a user, I want an in-app notifications centre so that I can track events without relying on email/SMS. &
Notifications have read/unread state and are grouped by type; users can open related objects from notifications. \\
\midrule
US-5.13 & As a user, I want important events to appear in real-time so that I can react quickly. &
Booking/session status changes trigger notifications; delivery uses Socket.IO when online, with polling fallback when offline. \\
\bottomrule
\caption{Sprint~5 user stories with acceptance criteria.}
\label{tab:sprint5-userstories}\\
\end{longtable}

\subsection{Engineering Backlog Summary}

Table~\ref{tab:sprint5-backlog} summarises the high-level engineering items and their
relative effort estimates.

\begin{longtable}{@{} L{0.10\textwidth} L{0.17\textwidth} L{0.45\textwidth} C{0.10\textwidth} C{0.07\textwidth} @{}}
\toprule
\textbf{ID} &
\textbf{Title} &
\textbf{User Story Summary} &
\textbf{Priority} &
\textbf{SP} \\
\midrule
\endfirsthead

\toprule
\textbf{ID} &
\textbf{Title} &
\textbf{User Story Summary} &
\textbf{Priority} &
\textbf{SP} \\
\midrule
\endhead

PB-16 & Psychologist Dashboard & Provide role-specific professional overview widgets and direct access to operational tasks. & High & 5 \\
\midrule
PB-17 & Document AI Querying & Allow secure upload and grounded question answering over documents. & High & 8 \\
\midrule
PB-18 & Ratings and Feedback & Enable post-session rating and moderated textual feedback. & Medium & 5 \\
\midrule
PB-19 & Admin Panel & Extend supervision, audit visibility, and user management capabilities. & High & 8 \\
\midrule
PB-20 & i18n and RTL & Support EN/FR/AR localization and Arabic right-to-left rendering. & High & 5 \\
\midrule
PB-21 & Floating Assistant & Offer persistent contextual guidance across pages. & Medium & 3 \\
\midrule
PB-22 & Notifications & Persist and deliver domain notifications in-app and in real time. & High & 5 \\
\bottomrule
\caption{Sprint~5 engineering backlog focused on platform enrichment and governance.}
\label{tab:sprint5-backlog}\\
\end{longtable}

\subsection{Definition of Done}

A sprint item is considered done when all of the following conditions are satisfied.

\begin{itemize}[leftmargin=*]
  \item[--] New modules are role-protected (patient/psychologist/admin).
  \item[--] AI document queries are permission-scoped and traceable (sources/excerpts returned and logged).
  \item[--] Ratings are session-bound with moderation tooling and audit logs in place.
  \item[--] Internationalization works end-to-end for EN/FR/AR; RTL layout passes UI checks.
  \item[--] Notifications are durable (stored in the database) and visible (real-time and in-app).
  \item[--] All deliverables are deployed to the staging environment for stakeholder review.
\end{itemize}

% ─────────────────────────────────────────────────────────────────────────────
\section{Requirement Analysis}
\label{sec:sprint5-analysis}

Sprint~5 introduces enrichment features that increase platform adoption and operational
maturity. The main challenges are: (1)~permission-scoped AI retrieval, (2)~correctness and
fairness of ratings, (3)~reliable notifications, and (4)~consistent multilingual UI behaviour
including RTL.

\subsection{Functional Analysis}

The sprint combines several modules with distinct interaction patterns but a shared
requirement for controlled access.

\begin{longtable}{@{} L{0.27\textwidth} L{0.65\textwidth} @{}}
\toprule
\textbf{Module} &
\textbf{Functional Analysis} \\
\midrule
\endfirsthead

\toprule
\textbf{Module} &
\textbf{Functional Analysis} \\
\midrule
\endhead

\emph{Dashboard} & Aggregate sessions, reports, and operational information into a
psychologist-focused overview with deep links to details and availability adjustments. \\
\midrule
\emph{Document AI querying} & Upload documents to secure storage; extract text (OCR
where needed); index content via chunking and embeddings for semantic retrieval;
generate answers grounded in retrieved chunks; return source excerpts for
verification; enforce access scopes (private, session-linked, admin-only). \\
\midrule
\emph{Ratings} & Attach one rating to a completed session; aggregate ratings per
psychologist; expose moderation operations with reasons and audit logs. \\
\midrule
\emph{Admin panel} & User listing, search, status changes, and role checks;
operational monitoring across payments, AI jobs, PDFs, sockets, and notifications;
audit log viewer with actor/action/entity/time-range filters. \\
\midrule
\emph{Internationalization and RTL} & Translation keys for UI text; locale-based
formatting for date/time/number; direction switching (LTR for EN/FR, RTL for AR)
with mirrored layout; preference persisted and applied consistently. \\
\midrule
\emph{Floating assistant} & UI overlay available across pages; supports lightweight
help and actions; uses safe prompts and routes to relevant features without
exposing restricted data. \\
\midrule
\emph{Notifications} & Generation on booking/session/report/admin events; in-app
persistence with read/unread state and deep-link payload; real-time delivery via
Socket.IO to user-specific rooms with API fallback. \\
\bottomrule
\caption{Functional analysis of Sprint~5 modules.}
\label{tab:sprint5-functional-analysis}\\
\end{longtable}

\subsection{Non-Functional Requirements}

The non-functional requirements are mainly related to privacy, reliability, usability, and
auditability.

\begin{longtable}{@{} L{0.27\textwidth} L{0.65\textwidth} @{}}
\toprule
\textbf{Requirement Area} &
\textbf{Requirement} \\
\midrule
\endfirsthead

\toprule
\textbf{Requirement Area} &
\textbf{Requirement} \\
\midrule
\endhead

\emph{Privacy and security} & Strict scoping for documents and AI retrieval; stored AI
prompts minimised; files encrypted at rest. \\
\midrule
\emph{Reliability} & Notification delivery retries with durable event storage and
idempotent consumers. \\
\midrule
\emph{Performance} & Dashboard aggregation queries optimised; document indexing
asynchronous; pagination across all lists. \\
\midrule
\emph{Usability} & Consistent UX across languages; RTL styling validated; floating
assistant non-intrusive. \\
\midrule
\emph{Auditability} & Admin actions, document access, AI queries, and rating
moderation all logged. \\
\bottomrule
\caption{Non-functional requirements for Sprint~5.}
\label{tab:sprint5-non-functional-requirements}\\
\end{longtable}

\subsection{Policy and Governance Notes}

\begin{itemize}[leftmargin=*]
  \item[--] AI document Q\&A must be constrained to returned sources to avoid hallucinations;
        grounded responses must always include references.
  \item[--] Ratings must not expose sensitive health details; content policy and moderation
        must be enforced.
  \item[--] Notification payloads should avoid leaking sensitive information in previews
        (configurable redaction).
\end{itemize}

\subsection{Technology and Integration Highlights}
 
\subsubsection*{Dashboard State Management}
 
On the server side, aggregation endpoints compute and return consistent Data Transfer Objects
(DTOs) that consolidate sessions, report statuses, upcoming tasks, and key metrics in a single
request, avoiding the N+1 query patterns that would arise from per-widget calls. On the client
side, each dashboard widget manages its own loading, error, and empty-state lifecycle
independently, so a single failing widget does not block the rest of the view. Role-based
visibility rules are enforced both at the API level (psychologist scope is validated on every
request) and at the component level (widgets are conditionally rendered based on the
authenticated role). A short-lived cache layer prevents redundant re-fetching when the user
navigates between dashboard sub-views within the same session.
 
\subsubsection*{Document AI Querying Pipeline}
 
The document AI pipeline is divided into two distinct phases: ingestion and query-time
retrieval.
 
\textbf{Ingestion phase.} When a user uploads a PDF document, the backend stores the raw file
in secure object storage and creates a \texttt{DocumentRecord} with \texttt{indexingStatus =
UPLOADED}. A background indexing worker then takes over: it extracts text using
\texttt{pdf-parse} (with optional OCR fallback for scanned documents), splits the text into
overlapping chunks using a \texttt{RecursiveCharacterTextSplitter} (default chunk size of 800
characters, overlap of 100), generates a dense embedding vector per chunk using the Google
Gemini \texttt{text-embedding-004} model, and upserts each chunk with its embedding and
permission-scoped metadata into MongoDB Atlas Vector Search (collection
\texttt{patient\_doc\_chunks}, index \texttt{vector\_index}). Once all chunks are indexed the
document status transitions to \texttt{INDEXED} and the user is notified that querying is
available.
 
\textbf{Query-time retrieval.} When a user submits a question, the backend first verifies that
the user holds a valid access scope for the target document (private ownership, session
participant, or admin policy). The query text is then embedded with the same
\texttt{text-embedding-004} model to produce a query vector. A similarity search against
\texttt{vector\_index} retrieves the top-$k$ most relevant chunks while applying the
permission filter as a metadata pre-filter, ensuring that chunks outside the user's scope are
never surfaced. The retrieved chunks, together with the original question, are passed to the
LLM answer service which generates a grounded response referencing only the provided evidence.
The answer is returned alongside the source chunk identifiers and excerpt text so that the
user can verify every claim. The entire interaction is persisted as a \path{DocumentQueryLog}
entry for auditability.
 
\textbf{Resilient retrieval fallback.} The retrieval strategy degrades gracefully when the
vector infrastructure is unavailable: the system first tries Atlas Vector Search; if that
fails, it falls back to a MongoDB text-index search; and finally to a regex-based scan. This
three-tier approach ensures that users can still retrieve relevant content even in degraded
infrastructure conditions.
 
\subsubsection*{Ratings System}
 
Eligibility is enforced by a compound unique index on \texttt{(sessionId, patientId)} at the
database level, making duplicate submissions idempotent rather than dependent on application
logic alone. When a rating is created or updated, an aggregation job recomputes the
psychologist's average score and review count. These aggregated values are stored as a
separate \texttt{RatingAggregate} document so that discovery queries can read a single
pre-computed record rather than computing averages on the fly across potentially thousands of
reviews. Moderation endpoints available to admins allow hiding or removing a review; every
moderation action is recorded in the \texttt{AuditLog} with the admin identity, the action
type, and the justification, ensuring accountability and appeal traceability. Rate-limiting
middleware on the submission endpoint prevents burst abuse.
 
\subsubsection*{Admin Panel}
 
All admin API routes are protected by a dedicated role guard that rejects requests from
non-admin tokens before any business logic executes. Every state-modifying action (account
suspension, onboarding status override, review moderation) writes an \texttt{AuditLog} entry
atomically within the same database transaction as the primary update, guaranteeing that no
governance action goes unrecorded. The monitoring view is powered by a lightweight aggregation
endpoint that queries counters for failed payment webhooks, failed AI indexing jobs, PDF
generation errors, socket connection drops, and notification delivery failures over a
configurable time window, presenting operational health without requiring an external
observability tool. The audit log viewer supports server-side filtering by actor, action type,
entity type, entity identifier, and time range, and returns paginated results to avoid
transferring large log volumes to the client.
 
\subsubsection*{Internationalization and RTL}
 
The frontend uses a key-based i18n library (\texttt{react-i18next}) where all user-visible
strings are referenced by translation keys rather than hard-coded values. Translation bundles
for EN, FR, and AR are loaded dynamically on locale switch to avoid bloating the initial
bundle. The selected locale is stored in the user's \texttt{UserPreference} document and
applied as a default on login; it is also propagated via an \texttt{Accept-Language} header
on API requests so that server-generated messages (validation errors, email notifications) can
be localised consistently.
 
For Arabic RTL rendering, the application sets \texttt{dir="rtl"} on the root HTML element
and adopts CSS logical properties throughout (\texttt{margin-inline-start/end},
\texttt{padding-inline-start/end}, \texttt{inset-inline-start/end}) so that layout directions
flip automatically without duplicating style rules. Icons with directional semantics (arrows,
progress indicators) are mirrored using a CSS class injected at the root level. Numeric and
date formatting delegates to the browser's \texttt{Intl.NumberFormat} and
\texttt{Intl.DateTimeFormat} APIs, passing the active locale code to ensure correct
numeral systems and calendar formats.
 
\subsubsection*{Floating Assistant}
 
The floating assistant is implemented as a persistent overlay React component that mounts once
at the application root and remains alive across page transitions, preserving conversation
history without re-initialisation. When the user submits a query, the frontend sends the
message to a dedicated \texttt{POST /api/assistant} endpoint along with the current route
context (page identifier, visible resource type). The backend resolves the user's role and
permitted scopes before constructing the prompt, ensuring that the assistant cannot surface
data from resources the user is not authorised to view. The prompt is constructed with an
explicit system instruction that constrains the assistant to FAQ-style answers and navigation
suggestions; it does not have direct access to patient records or clinical notes. Interactions
are optionally logged (with user consent) as non-identifiable traces for quality improvement.
 
\subsubsection*{Notifications Pipeline}
 
The notifications subsystem is built around a publish-subscribe pattern with durable storage.
When a domain event occurs (booking confirmed, session started, report available, admin
decision issued), the originating service publishes a typed event to an internal event bus.
A dedicated notification worker consumes these events, creates a \texttt{Notification} record
in MongoDB, and then attempts real-time delivery via Socket.IO by emitting a
\texttt{notification:new} event to the \texttt{user:\{userId\}} room. If the user is not
connected at that moment, the notification remains stored and is retrieved the next time
the client calls \texttt{GET /api/notifications}.
 
For optional out-of-band channels (email, SMS), a secondary queue worker picks up the same
event and dispatches the message through the configured provider. Both workers implement
exponential backoff retry logic with a configurable maximum attempt count; failures after
exhausting retries are recorded as \texttt{NotificationDeliveryAttempt} entries with their
error codes for operational monitoring. Notification payloads are designed with a
minimal-exposure principle: they contain only the event type, a human-readable summary, and
a deep-link reference to the related domain object, without embedding raw clinical or
financial data.
\subsection{Module Dependency Notes}

Sprint~5 depends explicitly on earlier work. Dashboard widgets consume session and report data
from Sprint~4; ratings rely on completed-session semantics; document AI querying extends the
storage and permission model introduced in the platform core; and notifications derive from
events emitted by booking, session, reporting, and administrative modules.

% ─────────────────────────────────────────────────────────────────────────────
\section{Architecture Overview}
\label{sec:sprint5-architecture}

The architecture of Sprint~5 connects enrichment modules to existing platform services. This
view is important because dashboard aggregation, AI document querying, notifications, ratings,
and administration operate across several data sources and access-control boundaries.

\subsection{Module Interaction}

Sprint~5 connects several pre-existing platform modules through cross-cutting services. The
dashboard aggregates data from sessions and reports; document AI querying relies on secure
storage and retrieval logic; and notifications are triggered by events emitted from booking,
session, report, or admin modules.

Since Sprint~5 is driven by cross-cutting runtime interactions, the most meaningful technical
diagrams are sequence diagrams rather than static component inventories. These behavioral
views are therefore grouped later in the workflow section, after the chapter establishes the
use-case and class perspectives.

\subsection{Database and Indexing Design}

Several new persistence concerns emerge in Sprint~5.
Table~\ref{tab:sprint5-artifacts} summarises the main artifacts used to support ratings,
notifications, document querying, and auditability.

\begin{longtable}{@{} L{0.20\textwidth} L{0.20\textwidth} L{0.52\textwidth} @{}}
\toprule
\textbf{Artifact} &
\textbf{Representative Fields} &
\textbf{Purpose} \\
\midrule
\endfirsthead

\toprule
\textbf{Artifact} &
\textbf{Representative Fields} &
\textbf{Purpose} \\
\midrule
\endhead

\texttt{Rating} &
sessionId, patientId, psychologistId, score, comment, status &
Stores post-session feedback with moderation support. \\
\midrule
\texttt{Notification} &
userId, type, payload, read, createdAt &
Persists user-targeted domain notifications. \\
\midrule
\texttt{DocumentRecord} &
ownerId, scope, fileRef, extractedText, indexingStatus &
Tracks uploaded documents and their AI-queryable state. \\
\midrule
\makecell[l]{\texttt{Document}\\\texttt{QueryLog}} &
userId, documentId, queryText, scopeFilter, createdAt &
Preserves traceability of AI document questions and scope decisions. \\
\midrule
\texttt{AuditLog} &
actorId, action, entityType, entityId, timestamp &
Preserves traceability of sensitive administrative or moderation actions. \\
\bottomrule
\caption{Core persistence and indexing artifacts associated with Sprint~5.}
\label{tab:sprint5-artifacts}\\
\end{longtable}

Since Sprint~5 introduces several governance-oriented artifacts, a class-level representation
is useful to show how persistent records support moderation, notification delivery, and
traceable AI interactions.

\begin{figure}[H]
  \centering
  \includegraphics[width=0.92\textwidth]{assets/s5/d2.png}
  \caption{Class diagram for the principal enrichment artifacts introduced in Sprint~5.}
  \label{fig:sprint5-class}
\end{figure}

The diagram emphasises that enrichment features are supported by explicit audit and
traceability structures rather than transient UI state alone.

% ─────────────────────────────────────────────────────────────────────────────
\section{Functional Specification}
\label{sec:sprint5-functional}

The functional specification presents the main user interactions introduced in Sprint~5. The
use case diagram in Figure~\ref{fig:sprint5-usecase} groups these interactions around
professional productivity, governance, and user engagement.

\begin{figure}[H]
  \centering
  \includegraphics[width=0.92\textwidth]{assets/s5/usecases1.jpeg}
  \caption{Overall use case diagram for Sprint~5.}
  \label{fig:sprint5-usecase}
\end{figure}

The grouping clarifies that Sprint~5 enriches the platform at multiple levels without
replacing the core flows delivered earlier.

\subsection{Selected Use Cases}

The following use case specifications describe the principal interactions introduced in
Sprint~5.

\begin{longtable}{@{} L{0.28\textwidth} L{0.66\textwidth} @{}}
\toprule
\multicolumn{2}{@{}l@{}}{\textbf{UC-16 \quad Consult Psychologist Dashboard}} \\
\midrule
\endfirsthead
\toprule
\multicolumn{2}{@{}l@{}}{\textbf{UC-16 \quad Consult Psychologist Dashboard}} \\
\midrule
\endhead

Primary Actor & Psychologist \\
\midrule
Goal & Access a role-specific operational summary of sessions, reports, and upcoming actions. \\
\midrule
Preconditions & Psychologist is authenticated and authorised to access professional features. \\
\midrule
Main Success Scenario & Psychologist opens the dashboard; the system loads widgets (upcoming sessions, recent sessions, reports pending, key metrics); psychologist selects an item to navigate to details. \\
\midrule
Alternate Flow & No data: dashboard shows empty-state guidance. \\
\midrule
Postconditions & Dashboard activity may be logged (non-sensitive analytics). \\
\bottomrule
\caption{Use case specification for dashboard consultation.}
\label{tab:uc-s5-dashboard}\\
\end{longtable}

\begin{longtable}{@{} L{0.28\textwidth} L{0.66\textwidth} @{}}
\toprule
\multicolumn{2}{@{}l@{}}{\textbf{UC-17 \quad Query Uploaded Documents with AI}} \\
\midrule
\endfirsthead
\toprule
\multicolumn{2}{@{}l@{}}{\textbf{UC-17 \quad Query Uploaded Documents with AI}} \\
\midrule
\endhead

Primary Actor & Authorised user \\
\midrule
Goal & Ask grounded questions over uploaded documents and receive source-supported answers. \\
\midrule
Preconditions & The user has access to the document scope and indexing has completed successfully. \\
\midrule
Main Success Scenario & User initiates upload; backend validates metadata and returns an upload URL; document is stored and an indexing job is created; user awaits indexing; user submits a question; system retrieves relevant chunks within permitted scope; system generates a grounded answer; system returns answer plus source excerpt references. \\
\midrule
Alternate Flows & Unsupported document: system rejects and explains constraints. Indexing failure: job marked failed and retry is available. No relevant content found: system returns ``no evidence found'' and suggests query refinements. \\
\midrule
Postconditions & Query is logged (privacy-aware) and auditable; answer remains linked to evidence excerpts. \\
\bottomrule
\caption{Use case specification for document AI querying.}
\label{tab:uc-s5-rag}\\
\end{longtable}

\begin{longtable}{@{} L{0.28\textwidth} L{0.66\textwidth} @{}}
\toprule
\multicolumn{2}{@{}l@{}}{\textbf{UC-18 \quad Rate Psychologist After Session}} \\
\midrule
\endfirsthead
\toprule
\multicolumn{2}{@{}l@{}}{\textbf{UC-18 \quad Rate Psychologist After Session}} \\
\midrule
\endhead

Primary Actor & Patient \\
\midrule
Goal & Leave a rating and optional textual feedback after a completed session. \\
\midrule
Preconditions & Patient completed a session with the psychologist. \\
\midrule
Main Success Scenario & Patient opens completed session details; submits rating (1--5) and optional feedback text; system validates eligibility (one per session, anti-abuse rules); system stores review and updates aggregates. \\
\midrule
Alternate Flow & Ineligible rating attempt: system blocks and explains reason. \\
\midrule
Postconditions & Review is visible per policy and moderation rules. \\
\bottomrule
\caption{Use case specification for rating submission.}
\label{tab:uc-s5-rating}\\
\end{longtable}

\begin{longtable}{@{} L{0.28\textwidth} L{0.66\textwidth} @{}}
\toprule
\multicolumn{2}{@{}l@{}}{\textbf{UC-19 \quad Admin User Management and Monitoring}} \\
\midrule
\endfirsthead
\toprule
\multicolumn{2}{@{}l@{}}{\textbf{UC-19 \quad Admin User Management and Monitoring}} \\
\midrule
\endhead

Primary Actor & Admin \\
\midrule
Goal & Manage user accounts, moderate content, and review operational health indicators. \\
\midrule
Preconditions & Admin is authenticated with the required privileges. \\
\midrule
Main Success Scenario & Admin searches for a user and opens profile details; changes account status with justification; reviews operational dashboards for failures and alerts; uses audit logs to investigate sensitive actions. \\
\midrule
Postconditions & All admin actions are fully audited. \\
\bottomrule
\caption{Use case specification for admin user management and monitoring.}
\label{tab:uc-s5-admin}\\
\end{longtable}

\begin{longtable}{@{} L{0.28\textwidth} L{0.66\textwidth} @{}}
\toprule
\multicolumn{2}{@{}l@{}}{\textbf{UC-21 \quad Use the Floating Assistant}} \\
\midrule
\endfirsthead
\toprule
\multicolumn{2}{@{}l@{}}{\textbf{UC-21 \quad Use the Floating Assistant}} \\
\midrule
\endhead

Primary Actors & Patient, Psychologist \\
\midrule
Goal & Obtain contextual help and lightweight navigation assistance without exposing
      restricted clinical data. \\
\midrule
Preconditions & User is authenticated and the assistant feature is enabled for the current
               environment. \\
\midrule
Main Success Scenario & User opens the assistant widget; the system interprets the contextual
                        request; safe suggestions, links, or guided actions are returned; the
                        user navigates to the relevant page or function. \\
\bottomrule
\caption{Use case specification for floating assistant interaction.}
\label{tab:uc-s5-assistant}\\
\end{longtable}

\begin{longtable}{@{} L{0.28\textwidth} L{0.66\textwidth} @{}}
\toprule
\multicolumn{2}{@{}l@{}}{\textbf{UC-22 \quad Generate and Deliver Notifications}} \\
\midrule
\endfirsthead
\toprule
\multicolumn{2}{@{}l@{}}{\textbf{UC-22 \quad Generate and Deliver Notifications}} \\
\midrule
\endhead

Primary Actors & System, Notification Service \\
\midrule
Goal & Persist user-relevant events and deliver them in real time when possible. \\
\midrule
Preconditions & A triggering domain event occurs and a target user can be resolved. \\
\midrule
Main Success Scenario & Domain event occurs (booking confirmed, report ready, etc.); notifications
                        service creates a notification record; if user is online the system emits a
                        real-time notification event via Socket.IO; user reads notification and
                        state becomes read accordingly. \\
\midrule
Alternate Flow & Real-time delivery fails: user retrieves notifications via polling API. \\
\midrule
Postconditions & Notification persists and is visible in the notification centre. \\
\bottomrule
\caption{Use case specification for notification delivery.}
\label{tab:uc-s5-notifications}\\
\end{longtable}

% ─────────────────────────────────────────────────────────────────────────────
\section{Workflow Description}
\label{sec:sprint5-workflow}

The most technically significant workflows in Sprint~5 concern dashboard aggregation,
contextual guidance, AI document querying, and notification delivery. The following UML
diagrams are presented directly in this sprint in order to keep the structural and behavioral
views aligned with the discussion.

\begin{figure}[H]
  \centering
  \includegraphics[width=0.92\textwidth]{assets/s5/d3.png}
  \caption{Sequence diagram of dashboard aggregation and contextual assistant support in Sprint~5.}
  \label{fig:sprint5-sequence-dashboard}
\end{figure}

This first sequence highlights how the dashboard and assistant cooperate at runtime to
provide operational visibility and safe user guidance.

The second sequence focuses on the AI document-query pipeline because it combines storage,
indexing, permission filters, retrieval, and generation.

\begin{figure}[H]
  \centering
  \includegraphics[width=0.92\textwidth]{assets/s5/d4.png}
  \caption{Sequence diagram of the document AI querying pipeline.}
  \label{fig:sprint5-sequence-rag}
\end{figure}

This sequence clarifies that generation is only the final step of a broader pipeline. Access
control, indexing readiness, and evidence retrieval are all required before the AI service
produces a response.

Because notifications constitute another cross-cutting concern, their event flow is
represented separately.

\begin{figure}[H]
  \centering
  \includegraphics[width=0.92\textwidth]{assets/s5/d5.png}
  \caption{Sequence view of notification persistence and delivery.}
  \label{fig:sprint5-sequence-notif}
\end{figure}

The notification sequence highlights the dual nature of the subsystem: notifications are both
durable records stored for later consultation and transient events delivered to connected
clients.

% ─────────────────────────────────────────────────────────────────────────────
\section{API Layer}
\label{sec:sprint5-api}

The API layer exposes the enrichment services delivered during Sprint~5.
Table~\ref{tab:sprint5-api} summarises the representative endpoint families and their
responsibilities.

\begin{longtable}{@{} L{0.30\textwidth} C{0.13\textwidth} L{0.49\textwidth} @{}}
\toprule
\textbf{Endpoint Family} &
\textbf{Method} &
\textbf{Responsibility} \\
\midrule
\endfirsthead

\toprule
\textbf{Endpoint Family} &
\textbf{Method} &
\textbf{Responsibility} \\
\midrule
\endhead

\texttt{/api/dashboard} & GET & Return psychologist dashboard aggregates and shortcut data. \\
\midrule
\texttt{/api/documents} & \makecell[c]{POST/\\GET} & Upload and list document resources with scope metadata. \\
\midrule
\texttt{/api/document-query} & POST & Execute a grounded question-answering request over indexed documents. \\
\midrule
\makecell[l]{\texttt{/api}\\\texttt{/ratings}} & \makecell[c]{\scriptsize POST/GET/\\\scriptsize PATCH} & Create ratings, expose aggregates, and moderate feedback. \\
\midrule
\texttt{/api/admin} & \makecell[c]{GET/\\PATCH} & Support user supervision and operational review. \\
\midrule
\texttt{/api/assistant} & POST & Return safe contextual help and navigation suggestions. \\
\midrule
\texttt{/api/notifications} & \makecell[c]{GET/\\PATCH} & Retrieve and mark notification records as read. \\
\bottomrule
\caption{Representative API responsibilities for Sprint~5.}
\label{tab:sprint5-api}\\
\end{longtable}

% ─────────────────────────────────────────────────────────────────────────────
\section{Security Considerations}
\label{sec:sprint5-security}

Security in Sprint~5 is inseparable from governance. Document AI querying must be constrained
by ownership or session scope; rating submission must remain linked to legitimate completed
sessions; admin actions must be logged; and notification payloads must avoid unnecessary
exposure of sensitive information.

\subsection{Governance Controls}

Table~\ref{tab:sprint5-governance-controls} summarises the governance-oriented controls
required to preserve secure and traceable platform enrichment.

\begin{longtable}{@{} L{0.24\textwidth} L{0.44\textwidth} L{0.24\textwidth} @{}}
\toprule
\textbf{Concern} &
\textbf{Control Mechanism} &
\textbf{Expected Outcome} \\
\midrule
\endfirsthead

\toprule
\textbf{Concern} &
\textbf{Control Mechanism} &
\textbf{Expected Outcome} \\
\midrule
\endhead

Document access & Scope-aware authorization and query logging & Prevent unauthorized semantic retrieval. \\
\midrule
Ratings fairness & Session-bound submission rules and moderation workflows & Preserve trustworthiness of public feedback. \\
\midrule
Admin actions & Audit logging and restricted admin endpoints & Ensure traceability of governance operations. \\
\midrule
Notifications privacy & Minimal payload design and target-user resolution & Reduce exposure of sensitive event details. \\
\bottomrule
\caption{Governance-oriented security controls in Sprint~5.}
\label{tab:sprint5-governance-controls}\\
\end{longtable}

% ─────────────────────────────────────────────────────────────────────────────
\section{Interface Screenshots}
\label{sec:sprint5-ui}
\begin{figure}[H]
  \centering
  \includegraphics[width=0.84\textwidth]{assets/interfaces/interface16.png}
  \caption{Sprint~5 psychologist dashboard interface.}
  \label{fig:sprint5-dashboard-ui}
\end{figure}
\begin{figure}[H]
  \centering
  \includegraphics[width=0.84\textwidth]{assets/interfaces/interface17.png}
  \caption{Sprint~5 document upload and AI query interface.}
  \label{fig:sprint5-document-ui}
\end{figure}
\begin{figure}[H]
  \centering
  \includegraphics[width=0.84\textwidth]{assets/interfaces/interface18.png}
  \caption{Sprint~5 notifications centre interface.}
  \label{fig:sprint5-notifications-ui}
\end{figure}
\begin{figure}[H]
  \centering
 \includegraphics[width=0.84\textwidth]{assets/interfaces/interface19.png}
  \caption{Sprint~5 floating assistant interface.}
  \label{fig:sprint5-assistant-ui}
\end{figure}
\begin{figure}[H]
  \centering
  \begin{subfigure}[t]{0.48\textwidth}
    \centering
    \includegraphics[width=\linewidth]{assets/interfaces/interface20.png}
  \end{subfigure}\hfill
  \begin{subfigure}[t]{0.48\textwidth}
    \centering
    \includegraphics[width=\linewidth]{assets/interfaces/interface21.png}
  \end{subfigure}
  \caption{Sprint~5 multilingual and RTL interface view.}
  \label{fig:sprint5-i18n-ui}
\end{figure}
\clearpage
\section{Tests}
\label{sec:sprint5-tests}

Sprint~5 requires tests across AI retrieval boundaries, localization correctness, and
notification reliability.

\subsection{Unit Testing Strategy and Scope}

Unit tests cover the following areas.

\begin{itemize}[leftmargin=*]
  \item[--] Dashboard aggregation logic: correct filters and time windows.
  \item[--] Ratings eligibility rules: completed-only, one per session, update window enforcement.
  \item[--] Document indexing pipeline utilities: chunking behaviour and metadata handling.
  \item[--] Permission scoping for retrieval: document visibility rules per scope type.
  \item[--] Notification formatting and state transitions: read/unread logic.
  \item[--] i18n: translation key presence and locale formatting utilities.
  \item[--] RTL: direction-aware layout rules where unit-testable (component snapshots).
\end{itemize}

External services (vector store, AI provider, notification gateway) are mocked in all unit
tests.

\subsection{Integration Testing Across Components}

Integration tests focus on the following cross-component flows.

\begin{itemize}[leftmargin=*]
  \item[--] Document upload $\rightarrow$ indexing job $\rightarrow$ query endpoint returns
        sources and respects permission scope.
  \item[--] Ratings creation affects psychologist aggregates and visibility in discovery
        endpoints.
  \item[--] Admin moderation actions hide reviews and produce auditable log entries.
  \item[--] Notifications: persistent storage, real-time delivery to connected users, and
        fallback retrieval via API when offline.
  \item[--] i18n/RTL: language switch persistence, correct RTL rendering on key screens, and
        correct date/time formatting.
\end{itemize}
\clearpage
\subsection{Manual End-to-End Scenario Validation}

Table~\ref{tab:sprint5-tests} presents the representative validation scenarios used to
confirm end-to-end correctness.

\begin{longtable}{@{} L{0.11\textwidth} L{0.48\textwidth} L{0.23\textwidth} C{0.12\textwidth} @{}}
\toprule
\textbf{Test ID} &
\textbf{Scenario} &
\textbf{Expected Result} &
\textbf{Result} \\
\midrule
\endfirsthead

\toprule
\textbf{Test ID} &
\textbf{Scenario} &
\textbf{Expected Result} &
\textbf{Result} \\
\midrule
\endhead

T5-1 & Load psychologist dashboard & Aggregated widgets and lists returned only for authorised psychologist users. & \textcolor{successgreen}{Passed} \\
\midrule
T5-2 & Query indexed document within scope & Answer returned with grounded excerpts and access checks enforced. & \textcolor{successgreen}{Passed} \\
\midrule
T5-3 & Query document outside scope & Request rejected and no content exposed. & \textcolor{successgreen}{Passed} \\
\midrule
T5-4 & Submit rating after completed session & Rating accepted once and linked to the correct session. & \textcolor{successgreen}{Passed} \\
\midrule
T5-5 & Admin suspends a user with justification & User status updated and audit log entry created. & \textcolor{successgreen}{Passed} \\
\midrule
T5-6 & Open floating assistant on a restricted page & Suggestions remain safe and do not expose unauthorised data. & \textcolor{successgreen}{Passed} \\
\midrule
T5-7 & Switch locale from English to Arabic & Interface text changes and layout direction becomes RTL. & \textcolor{successgreen}{Passed} \\
\midrule
T5-8 & Trigger a notification-worthy domain event & Notification persisted and delivered to the intended user channel. & \textcolor{successgreen}{Passed} \\
\midrule
T5-9 & Upload a document, wait for indexing, query with cross-user access attempt & Sources shown to owner; cross-user access blocked. & \textcolor{successgreen}{Passed} \\
\bottomrule
\caption{Representative validation scenarios for Sprint~5.}
\label{tab:sprint5-tests}\\
\end{longtable}

\section{Conclusion}
\label{sec:sprint5-conclusion}

Sprint~5 enriches \emph{PsychPlatform} with governance, usability, and engagement features
required by a more mature digital health platform. The psychologist dashboard improves
professional workflow visibility by aggregating sessions, reports, and workload indicators
into a single operational view. Document AI querying increases productivity while respecting
privacy through permission-scoped retrieval and source-grounded generation. The ratings
system strengthens trust by tying feedback to completed sessions and providing
moderation-ready controls. The extended admin panel supports governance through user
management, operational monitoring, and a filterable audit log. Internationalization with RTL
broadens accessibility to EN/FR/AR users and ensures layout consistency for Arabic speakers.
The floating assistant improves navigation and contextual support across the platform. The
notifications subsystem, combining durable storage with real-time Socket.IO delivery and
queue-based retry logic, provides timely and reliable engagement across all user roles. By
combining these capabilities under a shared permission-aware and auditable governance layer,
Sprint~5 extends the platform's operational quality while preserving the security principles
established in earlier sprints.

% ================= GENERAL CONCLUSION =================

\chapter*{General Conclusion}
\phantomsection
\addcontentsline{toc}{chapter}{General Conclusion}

\vspace{0.5cm}

This dissertation presented the design, implementation, and validation of
\textbf{PsychPlatform}, an AI-assisted psychological consultation platform
built on the MERN stack and developed across five SCRUM sprints over a
fifteen-week academic year.

\section*{Summary of Contributions}
\phantomsection
\addcontentsline{toc}{section}{Summary of Contributions}

The platform addresses a documented and consequential gap in the digital
mental health landscape: the absence of a unified system that connects
AI-assisted patient preparation, verified practitioner discovery, structured
clinical intelligence, and secure session management within a single coherent
workflow. Each of the five sprints delivered a distinct and verifiable
increment toward this objective.

\textbf{Sprint~1} established the security foundation of the platform through
JWT-based authentication, bcrypt password hashing, role-based access control,
and email verification, creating the trust boundary upon which all subsequent
modules depend.

\textbf{Sprint~2} introduced the psychologist onboarding pipeline,
combining OCR-based document extraction, LLM-assisted credential analysis,
a controlled state machine governing the verification lifecycle, and a
human-in-the-loop administrative review workflow. This pipeline operates
without dependency on national health system registries, making it
contextually applicable in Tunisia and North Africa where such registries
are not available.

\textbf{Sprint~3} delivered the patient engagement layer — multi-criteria
psychologist search, geospatial proximity discovery using MongoDB's
\texttt{2dsphere} index, interactive calendar management, and a stateful
booking lifecycle with split-slot handling and concurrency-safe confirmation.

\textbf{Sprint~4} represented the most technically dense integration effort,
assembling session activation via one-time cryptographic codes, real-time
Socket.IO messaging with voice support, PDF session report generation,
and the dual RAG chatbot pipeline. The latter constitutes the most
technically distinctive contribution of the project: two parallel retrieval
pipelines — one proactively enriching every chatbot turn with
preliminary Darija psychological context and clinical PDF chunks, the
other reactively answering explicit document queries with permission-scoped,
source-cited responses — both sharing the Gemini \texttt{text-embedding-004}
infrastructure and a three-tier degradation fallback strategy.

\textbf{Sprint~5} completed the platform with governance and engagement
features: the psychologist dashboard, a rating and moderation system,
an extended administrative panel with a filterable audit log, full
internationalisation supporting English, French, and Arabic with dynamic
right-to-left layout rendering, a persistent floating assistant, and a
durable real-time notification subsystem.

\section*{Acknowledged Limitations}
\phantomsection
\addcontentsline{toc}{section}{Acknowledged Limitations}

The platform has several limitations that the authors consider important
to state explicitly.

The \textbf{Darija RAG coverage} remains preliminary. The current retrieval
index improves the chatbot's dialect awareness relative to a general-purpose
LLM, but its lexical coverage is not exhaustive, its recall on unseen
dialectal expressions has not been formally evaluated, and its reliability
as a clinical-grade support tool requires further validation. Test scenario
T4-7 was marked as in-progress at submission time, reflecting this honest
assessment.

The \textbf{AI evaluation} is qualitative. The clinical summarisation engine,
the credential verification pipeline, and the document Q\&A system have been
validated through functional test scenarios but not through quantitative
metrics. No precision, recall, or human-rated quality scores have been
computed for the AI outputs. This limits the strength of the empirical claims
made on behalf of the AI components.

The \textbf{clinical knowledge base} consists of twenty psychoeducation
documents, which is a reasonable proof-of-concept scope but insufficient
for deployment as a clinical support tool. Expansion and curation by mental
health professionals would be required before any clinical use.

The platform has been validated in a development environment. It has not been
subjected to production-scale load testing, real-user usability evaluation,
or formal security audit, all of which would be required before deployment in
a real healthcare context.

\section*{Future Work}
\phantomsection
\addcontentsline{toc}{section}{Future Work}

Several directions are identified for future development.

\textbf{Darija corpus expansion and evaluation.} The most immediate priority
is the systematic enrichment of the Darija psychological expression index
through collaboration with linguists and mental health practitioners, followed
by a formal retrieval evaluation on a held-out test set.

\textbf{Quantitative AI evaluation.} A structured human evaluation protocol
for clinical summary quality — involving rated assessments by psychology
students or practitioners on a sample of generated summaries — would
substantially strengthen the empirical basis of the contribution.

\textbf{GDPR and data protection compliance.} As the platform handles
sensitive mental health data, a formal data protection impact assessment,
explicit informed consent mechanisms for AI-processed interactions, and
a right-to-erasure implementation would be required before any deployment
involving real patients.

\textbf{Video consultation.} The current session interface supports text
and voice messaging. Integration of WebRTC-based video consultation would
complete the teleconsultation feature set expected of a production platform.

\textbf{Production deployment and load testing.} Containerisation using
Docker, cloud deployment, and stress testing under concurrent session
load would allow the non-blocking architecture's performance claims to
be validated empirically.

\vspace{0.8cm}

\noindent
In conclusion, PsychPlatform demonstrates that the integration of dual RAG
pipelines, LLM-assisted clinical intelligence, and a complete consultation
lifecycle within a single full-stack platform is technically feasible at
the level of a supervised bachelor project. The platform does not replace
the psychologist; its purpose — to make every consultation more clinically
productive by ensuring that both patient and practitioner arrive better
prepared — has been architecturally realised and validated across all five
sprint increments. The work provides a sound foundation for continued
development toward a clinically deployable and linguistically inclusive
digital mental health service.

% ================= BIBLIOGRAPHY =================
\cleardoublepage
\phantomsection
\addcontentsline{toc}{chapter}{Bibliography}
\bibliographystyle{plain}
\bibliography{ref}
\end{document}